\documentclass[11pt,a4paper]{article}
\usepackage[T1]{fontenc}
\usepackage[utf8]{inputenc}
\usepackage{lmodern,amsmath,amssymb,amsthm,mathtools}
\usepackage[margin=26mm]{geometry}
\usepackage{microtype,booktabs,longtable,array,enumitem,listings,xcolor}
\usepackage[colorlinks=true,linkcolor=blue!40!black,citecolor=blue!40!black,urlcolor=blue!40!black]{hyperref}
\setlength{\emergencystretch}{4em}
\setlength{\parskip}{2pt}
\setlist{itemsep=3pt,topsep=5pt}
\setcounter{tocdepth}{1}
\newtheorem{theorem}{Theorem}[section]
\newtheorem{lemma}[theorem]{Lemma}
\newtheorem{proposition}[theorem]{Proposition}
\newtheorem{corollary}[theorem]{Corollary}
\newtheorem*{thmA}{Theorem A}
\newtheorem*{thmB}{Theorem B}
\newtheorem*{propC}{Proposition C}
\theoremstyle{definition}\newtheorem{definition}[theorem]{Definition}
\theoremstyle{remark}\newtheorem{remark}[theorem]{Remark}
\newtheorem{example}[theorem]{Example}
\newcommand{\Qb}{\overline{\mathbb Q}}
\newcommand{\E}{\mathbf E}\newcommand{\Ebar}{\overline{\mathbf E}}
\newcommand{\cR}{\mathcal R}\newcommand{\cM}{\mathcal M}
\newcommand{\Ein}{\operatorname{Ein}}
\newcommand{\F}{\mathfrak F}\newcommand{\FF}{\mathbb F}
\newcommand{\C}{\mathbb C}\newcommand{\Z}{\mathbb Z}
\newcommand{\Q}{\mathbb Q}\newcommand{\R}{\mathbb R}
\newcommand{\Tr}{\operatorname{Tr}}\newcommand{\Nm}{\operatorname{N}}
\DeclareMathOperator{\trdeg}{trdeg}\DeclareMathOperator{\Res}{Res}
\DeclareMathOperator{\Frac}{Frac}\DeclareMathOperator{\ord}{ord}
\lstset{basicstyle=\ttfamily\scriptsize,breaklines=true,columns=fullflexible,keepspaces=true,showstringspaces=false}
\hypersetup{pdftitle={Winding Embeddings and Prime-Power Pisot Floors},pdfauthor={KEIO},pdfsubject={Unified research manuscript with explicit hypotheses and further extensions}}
\title{Winding Embeddings and Prime-Power Pisot Floors\\[4pt]
\large Integral Intersections, Composite Windows, and Local Obstructions}
\author{KEIO}
\date{24 October 2026\\\small Unified and extended research manuscript}
\begin{document}
\maketitle
\begin{abstract}
We consolidate three research manuscripts on winding embeddings of
Euler--Bessel value fields and on composite values in prime-power Pisot
sequences. The analytic part develops polynomial dependence on winding,
exact integral-intersection arguments for two full Bessel blocks, three
harmonic coordinates, and affine projections, with imported independence
and nonvanishing inputs stated explicitly. The arithmetic part gives
simultaneous composite windows, shifted coefficient sieves, Pell-family
certificates, and field obstructions for rational powers of Mills' constant.
The extensions include a composite-progression theorem for polynomial
observables in every degree and every prime-power clock, a finite exact
winding criterion in bounded degree, and a closed formula for the
adjacent-derivative determinant of a normalized Bessel equation.
We also separate the two floor branches in finite-prime survival.
All reproduced proofs and finite certificates are contained in one source
file. Claims depending on unavailable earlier manuscripts remain explicitly
qualified. Neither Euler's constant nor Mills' constant is proved
transcendental, and the full cubic Pisot problem is not resolved.
\end{abstract}
\tableofcontents
\clearpage
\section{Scope, dependencies, and results of this edition}
\label{unified:scope}
This paper combines the supplied manuscripts \cite{WL,SourceExt,SourcePisot}
into one English article, with one bibliography and globally unique labels.
The two subjects share a concern with the exact reach of algebraic and local
certificates, but their proofs are independent; no transfer between them is
asserted merely because they occur in one paper. Source mathematical results
and proofs are retained, apart from duplicate statements replaced by internal
references and corrected descriptions of scope and provenance.

\paragraph{Logical status.}
The continued phase construction is derived from Fischler--Rivoal's preprint
and standard connection formulas in Section~\ref{foundation:section}.
Its extension to a continuous winding parameter does not require joint
algebraic independence. The exact analytic intersection classifications,
however, use the joint-independence assertions explicitly listed as (I5)
and their complex-scale or multi-Ein analogues imported from \cite{EB}.
That manuscript is not among the attachments. These classifications are
therefore proved here \emph{under those stated inputs}, rather than being
presented as independently established unconditional arithmetic theorems.
The fractional-order algebraicity-locus statement likewise invokes the
unprovided all-derivative nonvanishing theorem of \cite{BR}.
The determinant identities added below are proved directly and do not
depend on either manuscript.

The Pisot progression arguments, the coefficient sieve, the reproduced
Pell certificates, and the new prime-clock theorem have proofs in this
paper using standard published tools. The discriminant-$1524$ corollary
retains its explicitly identified earlier trace theorem as an assumption.
Finite prime patterns that invoke Schinzel's Hypothesis H remain conditional.
Historical status claims about earlier numbered research notes are records
of the supplied manuscript, not independent verification of absent files.

\paragraph{Further results proved in this edition.}
The composite-window mechanism extends from cubic traces to polynomial
observables of root-power coefficients in arbitrary degree, for every prime
clock. In particular, a Pisot number of any degree with a noninteger
three-adic trace limit has arbitrarily wide fixed windows of simultaneous
composite floors on an eventual arithmetic progression of indices.
The analytic additions give a bounded number of exact winding tests under
the stated independence inputs and an unconditional closed derivative
determinant. Additional finite-sieve constructions clarify why the two
floor branches have different local field restrictions.
These are extensions of the supplied arguments; no historical priority
or new resolution of a previously named global conjecture is asserted.

\paragraph{Notation.}
The analytic coefficient field is $k=\overline{\mathbb Q}$ and
$\overline{\mathbf E}$ denotes integral closure of the ring of $E$-values.
The arithmetic finite field is written $\mathbb F_p$; the calligraphic
Euler--Bessel field $\mathfrak F$ is a different object.
All positive rational powers of Mills' constant refer to its real branch.
The words ``composite progression'' always concern indices, with finite
initial exceptions removed after growth is established.
\clearpage
\part{Winding embeddings and special-function value fields}
% Standalone foundation insert; uses standard theorem environments.
% Bibliography keys FR and DLMF refer to Fischler--Rivoal 2026 and NIST DLMF.
% Labels use foundation:. Root may namespace the labels when integrating.
\section{A direct foundation for the winding construction}
\label{foundation:section}
This section separates the phase construction from the joint algebraic
independence hypothesis.  The former follows from the published-source
construction of Fischler--Rivoal \cite[Theorem 2.3 and Proposition 3.1]{FR}
and the standard connection identities below.  It does not require the
unprovided Euler--Bessel manuscript.  Joint independence, where used in
later intersection arguments, remains a separate input.

Write $k=\overline{\mathbb Q}$ and let
$\mathcal R=\mathbb C[k_{\rm add}]$ be the ring of finite exponential sums.
Fix a direction $-\pi/2<\theta<\pi/2$, using the right-hand convention at
an anti-Stokes direction.  The cited phase construction assigns canonical
components $f_\rho$ to an E-function $f$, with
$f(z)=\sum_\rho f_\rho(z)e^{\rho z}$.  Each component admits analytic
continuation along paths in $\mathbb C^\times$: the products
$f_\rho(z)e^{\rho z}$ solve an E-operator whose only possible finite
singularity is zero.  Denote continuation around a positively oriented
loop about zero, repeated $j\in\mathbb Z$ times, by a superscript $[j]$.

\begin{proposition}[Continued phase morphisms]\label{foundation:continued}
For every integer $j$ the rule
\[
 \Phi^{[j]}(f(1))(X)=\sum_\rho f_\rho^{[j]}(1)e^{\rho X}
\]
is a well-defined injective $k$-algebra homomorphism
$\mathbf E\longrightarrow\mathcal R$ satisfying
$\Phi^{[j]}(\xi)(1)=\xi$ and
$\Phi^{[j]}(e^a)=e^{aX}$ for algebraic $a$.
\end{proposition}
\begin{proof}
If $f(1)=g(1)$, Beukers' division theorem, as used in the proof of
\cite[Proposition 3.1]{FR}, gives an E-function $h$ with
$f-g=(z-1)h$.  Uniqueness of the canonical expansions gives
$f_\rho-g_\rho=(z-1)h_\rho$ componentwise.  Continuing this identity
along the same loop and evaluating at $1$ gives
$f_\rho^{[j]}(1)=g_\rho^{[j]}(1)$.  Thus the rule is well defined.
The component identities for sums and products continue along the loop;
for products the phase-$\rho$ component is the finite convolution
$\sum_{\alpha+\beta=\rho}f_\alpha g_\beta$.
Consequently the rule respects addition and multiplication.  Constants
are fixed.  Continuing the identity
$f(z)=\sum_\rho f_\rho(z)e^{\rho z}$ and using that $f$ is entire gives
$\sum_\rho f_\rho^{[j]}(1)e^\rho=f(1)$.  Evaluation at $X=1$ is therefore
a left inverse, proving injectivity.  The exponential identity follows
from its single constant component.
\end{proof}

\begin{lemma}[Normality and subgroup intersections]
\label{foundation:normal}
The ring $\mathcal R$ is integrally closed in its fraction field.
For every torsion-free abelian group $G$ and subgroup $\Lambda$,
\[
 \mathbb C[G]\cap\operatorname{Frac}(\mathbb C[\Lambda])
       =\mathbb C[\Lambda].
\]
Consequently, if $\xi\in\operatorname{Frac}(\mathbf E)$ is integral over
$\mathbf E$, then $\Phi^{[j]}(\xi)\in\mathcal R$ for every $j$.
\end{lemma}
\begin{proof}
For the normality assertion, put the finite supports of a fraction and
of the coefficients of its monic integral equation in one finitely
generated subgroup of $k_{\rm add}$.  That subgroup is free abelian.
Its group algebra is a Laurent polynomial ring over $\mathbb C$ and is
integrally closed.  The fraction therefore belongs to this ring, hence
to $\mathcal R$.
For the intersection assertion, write $u=f/g\in\mathbb C[G]$ with
$f,g\in\mathbb C[\Lambda]$ and $g\ne0$.  Decompose the finite sum $u$
according to its $G/\Lambda$ cosets.  Multiplication by $g$ preserves
every coset.  The identity $gu=f$ and the fact that the group algebra
of a torsion-free abelian group is a domain force every nonidentity
coset component of $u$ to vanish.  Thus $u\in\mathbb C[\Lambda]$.
Finally, applying a continued phase homomorphism to a monic equation
for $\xi$ makes its image integral over $\mathcal R$; normality applies.
\end{proof}

\subsection{Positive-argument component formulas}
For $t>0$ choose $\theta=-\pi/4$ and put $p=\pi i$.
The connection identity
$K_0(ze^{\pi i})=K_0(z)-\pi iI_0(z)$ shows that the two exponential
components of $I_0(tz)$ are
\[
 F_+(z)=I_0(tz)-K_0(tz)/p,
 \qquad F_-(z)=K_0(tz)/p,
\]
with phases $t$ and $-t$, respectively; thus the actual component
functions are $e^{-tz}F_+(z)$ and $e^{tz}F_-(z)$.
They have the canonical Gevrey asymptotic expansions in a common large
sector about $-\pi/4$, by the standard $K$ asymptotics and connection
formula \cite[\S\S10.34, 10.40]{DLMF}.  The power series identity
\[
 Y(tz)=(\gamma+\log(tz/2))I_0(tz)+K_0(tz)
\]
then gives the phase components
$g(z)F_+(z)$ and $(g(z)+p)F_-(z)$, where
$g(z)=\gamma+\log(tz/2)$.

Under $j$ positive loops,
\[
 K_0\longmapsto K_0-2\pi ij I_0,\qquad
 K_1\longmapsto K_1+2\pi ij I_1,\qquad
 g\longmapsto g+2\pi ij.
\]
The first two identities are \cite[10.34.5]{DLMF}; differentiation of
the two $I_0$ components gives the corresponding $I_1$ components.
Writing $I=I_0(t)$, $J=I_1(t)$, $K=K_0(t)$, and $L=K_1(t)$ therefore yields
\begin{align*}
 \Phi^{[j]}(I_0(t))&=A_je^{tX}+B_je^{-tX},\\
 \Phi^{[j]}(I_1(t))&=A'_je^{tX}+B'_je^{-tX},\\
 \Phi^{[j]}(Y(t))&=(g+j\tau)A_je^{tX}
                    +(g+j\tau+p)B_je^{-tX},
\end{align*}
where $g=\gamma+\log(t/2)$, $\tau=2\pi i$, and
\begin{align*}
 A_j&=e^{-t}((1+2j)I-K/p),&B_j&=e^t(K/p-2jI),\\
 A'_j&=e^{-t}((1+2j)J+L/p),&B'_j&=e^t(-2jJ-L/p).
\end{align*}
These are the positive-scale formulas used in the winding line.
For $a>0$, the identity
$\Ein(az)=\gamma+\log(az)+E_1(az)$, together with the phases $0,-a$,
gives directly
\[
 \Phi^{[j]}(\Ein(a))=\gamma+\log a+j\tau+
 (\Ein(a)-\gamma-\log a-j\tau)e^{a(1-X)}.
\]
The identity and its branch convention are recorded in
\cite[6.2.4]{DLMF}.  These derivations concern positive arguments;
complex-scale normalizations must be justified with a common summation
direction and the corresponding branches separately.

\subsection{Continuous winding without algebraic independence}
\begin{proposition}[Interpolation of genuine embeddings]
\label{foundation:interpolation}
Let $u_1,\ldots,u_N$ be generators of a subfield $F$ of
$\operatorname{Frac}(\mathbf E)$, and suppose their continued phase
images are rational functions $\Psi_i(z,s)$ in a finite phase torus
$z$ and a parameter $s$, at infinitely many distinct parameter values
$s_j$.  Suppose also that, for every $s_0\in\mathbb C$, all
$\Psi_i(z,s_0)$ are defined as rational functions, are regular at
$z=\boldsymbol 1$, and satisfy $\Psi_i(\boldsymbol1,s_0)=u_i$.
Then substitution defines a field embedding $\Phi_{s_0}:F\to
\mathbb C(z)$ for every $s_0$, regardless of algebraic relations among
the $u_i$.
\end{proposition}
\begin{proof}
If $P(u_1,\ldots,u_N)=0$ for $P\in k[T_1,\ldots,T_N]$, then the genuine
embeddings give $P(\Psi(z,s_j))=0$ at infinitely many $s_j$.
Clearing rational denominators shows that $P(\Psi(z,s))=0$ identically:
a polynomial in $s$ over $\mathbb C(z)$ cannot have infinitely many
distinct roots.  For a fixed $s_0$, the assumed regular substitution
therefore preserves every algebraic relation.  Conversely, when
$Q(u_1,\ldots,u_N)\ne0$, the rational function $Q(\Psi(z,s_0))$ is nonzero
because its value at $z=\boldsymbol1$ is that same nonzero number.
Substitution thus extends to the fraction field and is injective.
\end{proof}

This proposition applies to the explicit positive Euler--Bessel images.
It removes joint independence from the construction of $\Phi_s$ and
from the polynomial-winding proof: the latter only uses genuine
embeddings, normality, subgroup intersections, and the absence of
parameter-only denominators, all supplied above.  Joint independence
is still used when the later integral-intersection arguments identify
the coordinate fields with rational function fields and compare
coefficients.  It cannot be inferred merely from the existence of the
phase embeddings.
\section{Introduction}

The manuscript \cite{EB} studies the integral closure $\Ebar$ of the ring $\E$
of $E$-values through the continued phase embeddings
$\Phi^{[j]}=\Phi^{[j]}_{-\pi/4}$, $j\in\Z$, obtained from the phase embedding
of Fischler and Rivoal \cite{FR} by continuing the canonical components of an
$E$-function $j$ times around the origin. An element of $\Frac(\E)$ whose
image has a finite pole is not integral over $\E$. The manuscript closes with
four questions \cite[\S``Scope'']{EB}:
\begin{enumerate}
\item[(Q1)] the arithmetic nature of $\gamma$, $\delta$, and of any specified
member of $K_0(2),K_0(4),K_0(8)$;
\item[(Q2)] algebraic independence of $\gamma,\delta,\kappa=K_0(2)/I_0(2)$ and
$C=I_1(2)/I_0(2)$;
\item[(Q3)] integral-intersection classifications with two or more harmonic
Bessel coordinates at dependent scales, for example
$\Qb(\beta(2),C(2),\beta(4),C(4))$;
\item[(Q4)] an effective criterion selecting the density-one set of integer
arguments of \cite[Theorem~jden:density]{EB}.
\end{enumerate}
The following sections develop the winding construction and its consequences.
The missing joint-independence input is retained explicitly, as specified in
Section~\ref{unified:scope}.

Throughout, $k=\Qb$, $p=i\pi$, $\tau=2\pi i$, and $\cR=\C[k_{\rm add}]$ is the
ring of finite sums $\sum c_\rho e^{\rho X}$ with $\rho\in k$. For positive
algebraic $t$ we use $I_0(t),I_1(t),K_0(t)$, the harmonic series
$Y(t)=\sum_{n\ge0}H_n(t/2)^{2n}/(n!)^2$, and
\[
 \beta(t)=\frac{Y(t)}{I_0(t)}=\gamma+\log\frac t2+\frac{K_0(t)}{I_0(t)},\qquad
 C(t)=\frac{I_1(t)}{I_0(t)},\qquad h(t)=\frac1{tI_0(t)^2}.
\]

\begin{thmA}[the winding line; Proposition~\ref{wlbase:wl:family}, Theorem~\ref{wlbase:wl:poly}]
On the Euler--Bessel fields $\F$ of Definition~\ref{wlbase:wl:def} the embeddings
$\Phi^{[j]}$ are the members $s\in\gamma+\tau\Z$ of a family of field
embeddings $\Phi_s$, $s\in\C$, whose defining formulas do not contain
$\gamma$. If $\xi\in\F$ is integral over $\E$, its image, written in
exponential coordinates, is a Laurent polynomial whose coefficients are
polynomials in $s$.
\end{thmA}

\begin{thmB}[Theorem~\ref{wlbase:tb:main}]
For positive algebraic $a$ and distinct positive algebraic $t_1,t_2$,
\[
 \Ebar\cap k\bigl(e^a,\Ein(a),\beta(t_1),C(t_1),\beta(t_2),C(t_2)\bigr)
 =k[e^a,e^{-a},\Ein(a)],
\]
and $\Ebar\cap k(\beta(t_1),C(t_1),\beta(t_2),C(t_2))=k$.
\end{thmB}

\begin{propC}[Proposition~\ref{wlbase:lim:main}]
On the fields $\F$, algebraic relations among $E$-values, joint independence
and the conclusion of Theorem~A remain true when $\gamma$ is replaced, in the
definitions of $\delta$ and $K_0(t)$, by an arbitrary complex number.
Arguments using only these inputs cannot settle \textup{(Q1)} or
\textup{(Q2)}.
\end{propC}

The pure two-block intersection in Theorem~B is already reported in the
current version of \cite{EB}. The retained exponential--Ein pair is the
additional feature of the mixed statement. For
rationally independent scales it recovers the case $m=2$ of
\cite[Theorem~mbmulti:main]{EB}. Remark~\ref{wlbase:lim:three} explains why three or
more harmonic coordinates at mutually dependent scales need a further idea.
Proposition~C extends the elementary model of \cite[\S``Scope'']{EB}, which
concerns algebraic input only, to the integrality information carried by the
phase embeddings on $\F$. Within that information, $\gamma$ enters only
through the position $s=\gamma$ of the genuine embedding on the line
(Remark~\ref{wlbase:lim:point}). Other inputs of \cite{EB} do see $\gamma$, for
example the positivity of $K_0(t)$ and the critical approximation limits;
Proposition~C says nothing about them.

\section{The winding line}

\subsection{Inputs}
We use the following results of \cite{EB}.
\begin{enumerate}
\item[(I1)] \cite[Lemma~mbc:continued]{EB} For $j\in\Z$,
$\Phi^{[j]}\colon\E\to\cR$ is an injective $k$-algebra homomorphism with
$\Phi^{[j]}(\xi)(1)=\xi$ and $\Phi^{[j]}(e^a)=e^{aX}$; it extends to
$\Frac(\E)$.
\item[(I2)] \cite[Lemma~mbc:pole]{EB} If $\xi\in\Frac(\E)\cap\Ebar$, then
$\Phi^{[j]}(\xi)\in\cR$ for every $j$.
\item[(I3)] \cite[(mbc:Iimage)--(mbc:Yimage), (dcr:eta-image)]{EB} the explicit
winding images, recalled in Lemma~\ref{wlbase:wl:closed}.
\item[(I4)] \cite[Lemma~mbmulti:cosets]{EB} For a torsion-free abelian group
$G$ and a subgroup $\Lambda$, $\C[G]\cap\Frac(\C[\Lambda])=\C[\Lambda]$.
\item[(I5)] \cite[Theorem~cf:joint]{EB} For positive algebraic $a$ and
distinct positive algebraic $t_1,\dots,t_m$, the numbers
$e^a,\Ein(a),I_0(t_i),C(t_i),\beta(t_i)$ are algebraically independent over $k$.
\end{enumerate}

\subsection{Closed form of the images}

\begin{lemma}\label{wlbase:wl:closed}
Let $t>0$ be algebraic, $j\in\Z$ and $s=\gamma+j\tau$. Put
\[
 v=\beta(t)-\log\frac t2-s=\frac{K_0(t)}{I_0(t)}-j\tau,\qquad
 q=e^{2t(1-X)},\qquad \mu=\frac{q-1}{p+v(q-1)}.
\]
Then
\begin{equation}\label{wlbase:wl:img}
 \Phi^{[j]}(\beta(t))=\beta(t)-v(v-p)\mu,\qquad
 \Phi^{[j]}(C(t))=C(t)-h(t)\mu .
\end{equation}
For algebraic $a>0$ and $v_a=\Ein(a)-\log a-s$,
\begin{equation}\label{wlbase:wl:ein}
 \Phi^{[j]}(\Ein(a))=\Ein(a)+v_a\bigl(e^{a(1-X)}-1\bigr).
\end{equation}
\end{lemma}

\begin{proof}
Write $I,J,K,L$ for $I_0(t),I_1(t),K_0(t),K_1(t)$, $\kappa_t=K/I$,
$g=\gamma+\log(t/2)$ and $\varepsilon=1/p$. By (I3),
$\Phi^{[j]}(I_0(t))=A_je^{tX}+B_je^{-tX}$,
$\Phi^{[j]}(I_1(t))=A'_je^{tX}+B'_je^{-tX}$ and
$\Phi^{[j]}(Y(t))=(g+j\tau)A_je^{tX}+(g+j\tau+p)B_je^{-tX}$, with
\[
 A_j=e^{-t}((1+2j)I-\varepsilon K),\ B_j=e^t(\varepsilon K-2jI),\
 A'_j=e^{-t}((1+2j)J+\varepsilon L),\ B'_j=e^t(-2jJ-\varepsilon L).
\]
Since $e^{-2tX}=e^{-2t}q$ and $2j-\varepsilon\kappa_t=(j\tau-\kappa_t)/p=-v/p$,
\[
 A_j+B_je^{-2t}q=e^{-t}I\Bigl(1+\frac vp(q-1)\Bigr),\qquad
 pB_je^{-2t}q=e^{-t}(K-j\tau I)q=e^{-t}Ivq .
\]
Hence $\Phi^{[j]}(\beta(t))=g+j\tau+pqv/(p+v(q-1))$. As $g+j\tau=\beta(t)-v$,
this equals $\beta(t)+v\bigl(pq-p-v(q-1)\bigr)/(p+v(q-1))=\beta(t)-v(v-p)\mu$.
Similarly $A'_j+B'_je^{-2t}q=e^{-t}(J+(2jJ+\varepsilon L)(1-q))$, and the
Wronskian $IL+JK=1/t$ gives
$\Phi^{[j]}(C(t))-C(t)=\varepsilon(1-q)/(tI^2(1+(v/p)(q-1)))=-h(t)\mu$.
Formula \eqref{wlbase:wl:ein} is \cite[(dcr:eta-image)]{EB}: there
$\eta_j=g_a+j\tau+d_je^{a(1-X)}$ with $g_a=\gamma+\log a$ and
$d_j=\Ein(a)-g_a-j\tau=v_a$.
\end{proof}

\subsection{The family}

\begin{definition}\label{wlbase:wl:def}
Fix positive algebraic $a$ and distinct positive algebraic $t_1,\dots,t_m$,
and write $\eta_a=\Ein(a)$, $\beta_i=\beta(t_i)$, $C_i=C(t_i)$, $h_i=h(t_i)$.
By (I5) the field
\[
 \F=k(e^a,\eta_a,\beta_1,C_1,\dots,\beta_m,C_m)
\]
is purely transcendental on these generators; identify it with
$k(E,U,x_1,y_1,\dots,x_m,y_m)$. For $s\in\C$ put
\[
 v_i(s)=\beta_i-\log\frac{t_i}2-s,\quad v_a(s)=\eta_a-\log a-s,\quad
 q_i=e^{2t_i(1-X)},\quad \mu_i(s)=\frac{q_i-1}{p+v_i(s)(q_i-1)},
\]
and let $\Phi_s$ send
\[
 E\mapsto e^{aX},\quad U\mapsto\eta_a+v_a(s)\bigl(e^{a(1-X)}-1\bigr),\quad
 x_i\mapsto\beta_i-v_i(s)\bigl(v_i(s)-p\bigr)\mu_i(s),\quad
 y_i\mapsto C_i-h_i\mu_i(s).
\]
\end{definition}

\begin{proposition}\label{wlbase:wl:family}
Let $\cM$ be the field of meromorphic functions of $X\in\C$.
\begin{enumerate}
\item For every $s\in\C$, $\Phi_s$ extends to a field embedding
$\F\to\cM$ with $\Phi_s(\xi)(1)=\xi$.
\item $\Phi_{\gamma+j\tau}=\Phi^{[j]}|_{\F}$ for every $j\in\Z$.
\item The data defining $\Phi_s$ are $e^a$, $\eta_a$, $\log a$, $\beta_i$,
$C_i$, $h_i$, $\log(t_i/2)$, $\pi$ and $s$. Euler's constant does not occur.
\end{enumerate}
\end{proposition}

\begin{proof}
At $X=1$ one has $q_i=1$, $\mu_i=0$ and $e^{a(1-X)}=1$, so each generator is
mapped to a function whose value at $X=1$ is the generator itself, for every
$s$. A nonzero polynomial $Q$ is therefore mapped to a function with value
$Q(e^a,\eta_a,\beta,C)\ne0$ at $X=1$, by (I5). Thus the substitution
homomorphism is injective and extends to $\F$, with $\Phi_s(\xi)(1)=\xi$.
Part (2) holds because both sides are field homomorphisms agreeing on the
generators (Lemma~\ref{wlbase:wl:closed}). Part (3) is read off from the formulas.
\end{proof}

\begin{remark}\label{wlbase:wl:general}
On $\F$ the images are rational functions of $s$, and they depend on the
winding number only through $s$ because the only monodromy comes from
$\log z$, carried by $K_0$ and $\Ein$. For a general $E$-value the
components in \cite{FR} may have non-integral rational exponents at $0$;
then $j\mapsto\Phi^{[j]}(\xi)$ is a quasi-polynomial, polynomial on each
residue class of $j$ modulo a common denominator.
\end{remark}

\subsection{Polynomial winding}
Let $\Lambda\subset\mathbb R$ be the group generated by $a,2t_1,\dots,2t_m$,
choose a basis $\lambda_1,\dots,\lambda_r$ and put
$z_k=e^{\lambda_k(1-X)}$. The $\lambda_k$ are $\Q$-linearly independent, so
$z_1,\dots,z_r$ are algebraically independent functions; $q_i$ and
$e^{a(1-X)}$ are Laurent monomials in $z$, and
$e^{aX}=e^ae^{-a(1-X)}$. Write $\Psi(z,s)$ for the tuple of images of the
generators, a tuple of rational functions of $(z,s)$. For
$\xi=P(e^a,\eta_a,\beta,C)/Q(e^a,\eta_a,\beta,C)$ with coprime
$P,Q\in k[E,U,x,y]$, put
\[
 F_\xi=P(\Psi)/Q(\Psi)\in\C(z_1,\dots,z_r,s),
\]
so that $\Phi_s(\xi)(X)=F_\xi(z(X),s)$; since the $z_k(X)$ are algebraically
independent, $F_\xi$ is determined by $\xi$.

\begin{theorem}[polynomial winding]\label{wlbase:wl:poly}
If $\xi\in\F\cap\Ebar$, then
$F_\xi\in\C[z_1^{\pm1},\dots,z_r^{\pm1},s]$. In particular, for every
one-parameter subgroup $w\mapsto(w^{u_1},\dots,w^{u_r})$, the restriction of
$F_\xi$ lies in $\C[w,w^{-1},s]$.
\end{theorem}

\begin{proof}
Work in the factorial ring $\C[z^{\pm1},s]$, whose units are the elements
$cz^b$ with $c\in\C^\times$. Write $F_\xi=N/D$ with $N,D\in\C[z^{\pm1},s]$
coprime, and let $s_j=\gamma+j\tau$.

\emph{Integer points.} By Proposition~\ref{wlbase:wl:family}(2) and (I2),
$\Phi_{s_j}(\xi)=\Phi^{[j]}(\xi)\in\cR$. It also lies in $\Frac(\C[\Lambda])$,
since the images of the generators do, so by (I4) it lies in
$\C[\Lambda]=\C[z^{\pm1}]$. For all but finitely many $j$ the specialization
$D(\cdot,s_j)$ is not identically zero, and then
$N(\cdot,s_j)/D(\cdot,s_j)\in\C[z^{\pm1}]$.

\emph{Denominators.} Let $D_0$ be an irreducible factor of $D$ and suppose
$D_0$ is not associate to an element of $\C[s]$. Then $D_0$ has at least two
monomials in $z$ with coefficients not identically zero, so for all but
finitely many $s_0$ the Laurent polynomial $D_0(\cdot,s_0)$ is not a monomial
and has zeros in the torus $(\C^\times)^r$; its zero set there has dimension
$r-1$. The common zero set of $D_0$ and $N$ in $(\C^\times)^r\times\C$ has
dimension at most $r-1$, so its fibre over all but finitely many $s_0$ has
dimension at most $r-2$. Hence for all but finitely many $j$ there is a point
of the torus where $D_0(\cdot,s_j)$ vanishes and $N(\cdot,s_j)$ does not,
and $N(\cdot,s_j)/D(\cdot,s_j)$ has a pole there. This contradicts the first
step. So $D=c(s)$ up to a unit, with $c\in\C[s]\setminus\{0\}$.

\emph{Every $s_0$.} Fix $s_0\in\C$. The tuple $\Psi(\cdot,s_0)$ consists of
well-defined rational functions of $z$, and $Q(\Psi(\cdot,s_0))\not\equiv0$
because its value at $z=(1,\dots,1)$ is $Q(e^a,\eta_a,\beta,C)\ne0$. Hence
$G_0=P(\Psi(\cdot,s_0))/Q(\Psi(\cdot,s_0))$ is a rational function of $z$,
and $F_\xi(z,s)\to G_0(z)$ as $s\to s_0$ for $z$ outside a proper Zariski-closed
set. If $c(s_0)=0$, coprimality gives $N(\cdot,s_0)\not\equiv0$, and
$|F_\xi(z,s)|\to\infty$ as $s\to s_0$ for $z$ outside a proper closed set,
which is impossible. Therefore $c$ has no zeros and
$F_\xi\in\C[z^{\pm1},s]$. Restriction to a one-parameter subgroup preserves
this form.
\end{proof}

Thus integrality is a polynomial identity in the winding variable. The pole
criterion can be applied at every complex $s$, and the asymptotic charts of
\cite{EB}, taken there along integer windings $j\to\infty$, become limits of
polynomials in $(w,s)$.

\section{Two harmonic coordinates at arbitrary scales}

\begin{theorem}\label{wlbase:tb:main}
Let $a$ be positive algebraic and $t_1\ne t_2$ positive algebraic. Then
\[
 \Ebar\cap k\bigl(e^a,\Ein(a),\beta(t_1),C(t_1),\beta(t_2),C(t_2)\bigr)
 =k[e^a,e^{-a},\Ein(a)].
\]
\end{theorem}

\begin{corollary}\label{wlbase:tb:cor}
For distinct positive algebraic $t_1,t_2$,
$\Ebar\cap k(\beta(t_1),C(t_1),\beta(t_2),C(t_2))=k$. In particular
$\Ebar\cap k(\beta(2),C(2),\beta(4),C(4))=k$.
\end{corollary}

\begin{proof}
Choose any $a$ and put $G=k(\beta(t_1),C(t_1),\beta(t_2),C(t_2))$. By the
theorem, $\Ebar\cap G$ lies in $k[e^{\pm a},\Ein(a)]\cap G$. This is $k$,
because $e^a$ and $\Ein(a)$ are algebraically independent over $G$ by (I5).
\end{proof}

The inclusion $\supseteq$ in Theorem~\ref{wlbase:tb:main} holds since $e^{\pm a}$ and
$\Ein(a)$ are $E$-values. The rest of this section proves $\subseteq$.

\subsection{Normalization on a one-parameter subgroup}
In the notation of \S2 with $m=2$, write $q_i=z^{n_i}$ and $e^{a(1-X)}=z^{m_0}$.
For an integer $L$ let $\mathbf u\in\Z^r$ have $k$-th entry the nearest
integer to $L\lambda_k$; then $\langle n_i,\mathbf u\rangle=2Lt_i+O(1)$ and
$\langle m_0,\mathbf u\rangle=La+O(1)$. For large $L$ the restriction to the
one-parameter subgroup $w\mapsto w^{\mathbf u}$ gives
\[
 q_i\mapsto w^{N_i},\qquad e^{a(1-X)}\mapsto w^M,\qquad
 N_1\ne N_2,\ N_1,N_2,M\in\Z_{>0}.
\]
Let $v=v_1(s)$; this is an affine bijection of $\C$. Put
$\Delta=v_2-v_1=K_0(t_2)/I_0(t_2)-K_0(t_1)/I_0(t_1)$ and
$\Delta_a=v_a-v_1$; both are constants. On the subgroup the images of
$(E,U,x_1,y_1,x_2,y_2)$ form
\[
 \Psi(w,v)=\bigl(e^aw^{-M},\ \eta_a+(v+\Delta_a)(w^M-1),\
 B_1,\Gamma_1,B_2,\Gamma_2\bigr),
\]
\[
 B_i=\beta_i-v_i(v_i-p)\mu_i,\quad \Gamma_i=C_i-h_i\mu_i,\quad
 \mu_i=\frac{w^{N_i}-1}{D_i},\quad D_i=p+v_i(w^{N_i}-1),
\]
with $v_1=v$ and $v_2=v+\Delta$. Let $\xi=R(e^a,\eta_a,\beta_1,C_1,\beta_2,C_2)$
lie in $\Ebar$, where $R=P/Q$ with coprime $P,Q\in
k[E,U,x_1,y_1,x_2,y_2]$. By Theorem~\ref{wlbase:wl:poly},
\[
 F:=R\circ\Psi\in\C[w,w^{-1},v].
\]
We use this in the following form.

\smallskip\noindent
($\ast$) \emph{There is no point $(w_0,v_0)\in\C^\times\times\C$ with all
$D_i(w_0,v_0)\ne0$, $Q(\Psi(w_0,v_0))=0$ and $P(\Psi(w_0,v_0))\ne0$.}

\smallskip
Indeed $F$ would have a pole there.

\subsection{Independence}

\begin{lemma}\label{wlbase:tb:indep}
\begin{enumerate}
\item The numbers $e^a,\eta_a,\beta_1,C_1,\beta_2,C_2,h_1,h_2$ are algebraically
independent over $k$.
\item For $i\in\{1,2\}$ and $u_i^*=\eta_a-Mp/N_i$, the numbers
$e^a,C_1,C_2,h_1,h_2,u_i^*$ are algebraically independent over $k$.
\end{enumerate}
\end{lemma}

\begin{proof}
(1) is (I5), since $h_i$ is a nonzero algebraic multiple of $I_0(t_i)^{-2}$.
For (2), let $S$ be a subset of $\{e^a,C_1,C_2,h_1,h_2\}$. The field $k(S)$ is
stable under complex conjugation. If $u_i^*$ were algebraic over $k(S)$, so
would be its conjugate $\eta_a+Mp/N_i$, hence $\eta_a$, contradicting (1).
If adjoining $u_i^*$ to the independent set $\{e^a,C_1,C_2,h_1,h_2\}$ produced a
relation, some element $x$ of the set would be algebraic over $k(S)(u_i^*)$,
with $S$ the other elements, but not over $k(S)$. Exchange would make $u_i^*$
algebraic over $k(S\cup\{x\})$, which was just excluded.
\end{proof}

\subsection{Asymptotic regimes}

\begin{lemma}\label{wlbase:tb:regimes}
Let $|v|\to\infty$. The following hold uniformly for $c$ in compact subsets
of the indicated sets.
\begin{enumerate}
\item[(A)] $w=1+c\,v^{-2}$, $c\in\C$:
$\Psi\to(e^a,\eta_a,\beta_1-N_1c/p,\,C_1,\,\beta_2-N_2c/p,\,C_2)$, and
$D_i\to p$.
\item[(B)] $w=1+c\,v^{-3/2}$ for a fixed branch of $v^{1/2}$, $c\in\C$:
$(E,U,\Gamma_1,\Gamma_2)\to(e^a,\eta_a,C_1,C_2)$, $D_i\to p$, and
\[
 B_i=\beta_i-N_i\frac cp\,v^{1/2}+N_i^2\Bigl(\frac cp\Bigr)^2+O(|v|^{-1/2}).
\]
\item[(C$_i$)] $w$ the branch near $1$ of
$(1-p/v_i+c/v_i^2)^{1/N_i}$, $c\in\C^\times$, and $l\ne i$:
$(E,U,\Gamma_i,\Gamma_l)\to(e^a,u_i^*,C_i+ph_i/c,C_l)$,
$D_i=c/v_i$, $D_l\to p(N_i-N_l)/N_i$, and
\[
 B_i=\frac pc\,v^2+O(|v|),\qquad B_l=\frac{N_l}{N_i-N_l}\,v+O(1).
\]
\item[(P$_i$)] On $\Pi_i=\{D_i=0\}$, parametrized by $w=1+\omega$ with
$\omega\to0$ (so $v_i=p/(1-(1+\omega)^{N_i})\to\infty$), and $l\ne i$:
$(E,U,\Gamma_l)\to(e^a,u_i^*,C_l)$, $D_l\to p(N_i-N_l)/N_i$,
\[
 v_i(v_i-p)=\frac{p^2}{N_i^2\omega^2}+O(|\omega|^{-1}),\qquad
 B_l=-\frac{pN_l}{N_i(N_i-N_l)}\,\omega^{-1}+O(1).
\]
\end{enumerate}
\end{lemma}

\begin{proof}
Put $\theta_i=w^{N_i}-1$, so $\mu_i=\theta_i/(p+v_i\theta_i)$ and
$B_i-\beta_i=-v_i(v_i-p)\mu_i$.

(A) Here $\theta_i=N_icv^{-2}+O(|v|^{-4})$, $v_i\theta_i=O(|v|^{-1})$,
$\mu_i=N_ic/(pv^2)+O(|v|^{-3})$ and $v_i(v_i-p)=v^2+O(|v|)$. Also
$(v+\Delta_a)(w^M-1)=O(|v|^{-1})$.

(B) Here $\theta_i=N_icv^{-3/2}+O(|v|^{-3})$ and
$v_i\theta_i=N_icv^{-1/2}+O(|v|^{-3/2})$. Expanding
$\mu_i=(\theta_i/p)\bigl(1-v_i\theta_i/p+(v_i\theta_i/p)^2-\cdots\bigr)$ gives
$v_i(v_i-p)\mu_i=N_i(c/p)v^{1/2}-N_i^2(c/p)^2+O(|v|^{-1/2})$. Moreover
$\Gamma_i-C_i=O(|v|^{-3/2})$ and $U-\eta_a=O(|v|^{-1/2})$.

(C$_i$) By construction $\theta_i=-p/v_i+c/v_i^2$, so $D_i=c/v_i$,
$\mu_i=-p/c+1/v_i$, $\Gamma_i\to C_i+ph_i/c$ and
$B_i=\beta_i+v_i(v_i-p)(p/c-1/v_i)=(p/c)v^2+O(|v|)$. For $l\ne i$,
$\theta_l=(N_l/N_i)\theta_i+O(|v|^{-2})$, hence
$v_l\theta_l=-(N_l/N_i)p+O(|v|^{-1})$, $D_l\to p(N_i-N_l)/N_i$,
$\mu_l=-N_l/((N_i-N_l)v)+O(|v|^{-2})$, $\Gamma_l\to C_l$ and
$B_l=N_lv/(N_i-N_l)+O(1)$. Finally $w^M-1=(M/N_i)\theta_i+O(|v|^{-2})$ gives
$U\to\eta_a-Mp/N_i=u_i^*$.

(P$_i$) Here $\theta_i=N_i\omega+O(\omega^2)$ and
$v_i=-p/\theta_i=-p/(N_i\omega)+O(1)$. For $l\ne i$,
$\theta_l=(N_l/N_i)\theta_i+O(\omega^2)$ and $v_l=v_i+O(1)$, so
$v_l\theta_l\to-(N_l/N_i)p$, $D_l\to p(N_i-N_l)/N_i$ and
$\mu_l=N_iN_l\omega/(p(N_i-N_l))+O(\omega^2)$. This gives the
formula for $B_l$ and $\Gamma_l\to C_l$. Also $(v+\Delta_a)(w^M-1)\to-pM/N_i$.
\end{proof}

\subsection{Proof of Theorem~\ref{wlbase:tb:main}}
Write $T=(e^a,\eta_a,\beta_1,C_1,\beta_2,C_2)$ and $N'=(0,0,N_1,0,N_2,0)$.
Throughout, $f$ denotes an irreducible factor of $Q$ in $k[E,U,x,y]$; since
$Q=f\,Q_1$, every zero of $f\circ\Psi$ at which $Q_1\circ\Psi$ is finite is a
zero of $Q\circ\Psi$.

\emph{Step 1: every $f$ is invariant under translation by $N'$.} Put
$\varphi(c)=f(T-(c/p)N')$, the limit of $f\circ\Psi$ in regime (A). Suppose
$\varphi$ is not constant and let $c^*$ be a zero. By Hurwitz's theorem
$f\circ\Psi(1+cv^{-2},v)$ has zeros $c_v\to c^*$ for large $|v|$. The set
$V=Z(P)\cap Z(f)$ has dimension at most $4$, and the Zariski closure of
$V+\C N'$ is a proper subvariety of $\C^6$ defined over $k$. By
Lemma~\ref{wlbase:tb:indep}(1) the point $T$ does not lie on it, so
$P(T-(c^*/p)N')\ne0$. Hence $P\circ\Psi\ne0$ at these zeros for large $|v|$,
while $Q\circ\Psi=0$ and $D_i\ne0$, contradicting ($\ast$). Thus
$\lambda\mapsto f(T+\lambda N')$ is constant, and since $T$ is generic over
$k$, $f(X+\lambda N')=f(X)$ identically.

\emph{Step 2: $Q$ does not depend on $x_1,x_2$.} Put $\varpi=N_2x_1-N_1x_2$ and
$\sigma=(N_1x_1+N_2x_2)/(N_1^2+N_2^2)$. By Step~1,
$Q\in k[E,U,\varpi,y_1,y_2]$. Let $f$ have $\deg_{\varpi}f>0$. In regime (B),
\[
 \varpi(B)\to \varpi_0+\kappa_N(c/p)^2,\qquad
 \varpi_0=N_2\beta_1-N_1\beta_2,\quad \kappa_N=N_1N_2(N_1-N_2)\ne0,
\]
and $\sigma(B)=-(c/p)v^{1/2}+O(1)$. The polynomial
$f(e^a,\eta_a,W,C_1,C_2)$ has positive degree in $W$, its leading coefficient
being nonzero at an independent point. Its roots $W^*$ are algebraic over
$k(e^a,\eta_a,C_1,C_2)$, while $\varpi_0$ is transcendental over that field.
So $\varphi(c)=f(e^a,\eta_a,\varpi_0+\kappa_N(c/p)^2,C_1,C_2)$ has a zero
$c^*\ne0$, and Hurwitz gives zeros $c_v\to c^*$ of $f\circ\Psi$. The other
factors of $Q$ converge, so $Q\circ\Psi=0$ at these zeros, and $D_i\ne0$
there. Write $P=\sum_nP_{(n)}(E,U,\varpi,y)\sigma^n$ and let $n_0$ be the
largest $n$ with $f\nmid P_{(n)}$; it exists because $f\nmid P$. At the zeros
the terms with $f\mid P_{(n)}$ vanish exactly. Since $f$ is irreducible of
positive degree in $\varpi$ and $f\nmid P_{(n_0)}$, the resultant in $\varpi$
is nonzero; evaluated at the independent point $(e^a,\eta_a,C_1,C_2)$, it shows
$P_{(n_0)}(e^a,\eta_a,W^*,C_1,C_2)\ne0$. Therefore
\[
 P\circ\Psi=P_{(n_0)}(e^a,\eta_a,W^*,C_1,C_2)\,(-c^*/p)^{n_0}\,v^{n_0/2}(1+o(1))\ne0,
\]
contradicting ($\ast$). Hence $Q\in k[E,U,y_1,y_2]$.

\emph{Step 3: $Q$ does not depend on $y_1,y_2$.} Let $f$ have
$\deg_{y_1}f>0$ and use regime (C$_1$). There $f\circ\Psi$ converges to
$\varphi(c)=f(e^a,u_1^*,C_1+ph_1/c,C_2)$. By Lemma~\ref{wlbase:tb:indep}(2),
$f(e^a,u_1^*,Y,C_2)$ has positive degree in $Y$, and each root $Y^*$ is
algebraic over $k(e^a,u_1^*,C_2)$, hence $Y^*\ne C_1$. So $\varphi$ has the zero
$c^*=ph_1/(Y^*-C_1)\in\C^\times$, and Hurwitz gives zeros $c_v\to c^*$ of
$f\circ\Psi$; there $Q\circ\Psi=0$, $D_1=c_v/v_1\ne0$ and $D_2\ne0$. Write
$P=\sum_{a',b'}P_{(a',b')}(E,U,y_1,y_2)x_1^{a'}x_2^{b'}$ and let $W_0$ be the
largest value of $2a'+b'$ with $f\nmid P_{(a',b')}$. Terms with
$f\mid P_{(a',b')}$ vanish exactly at the zeros, and
$B_1^{a'}B_2^{b'}=(p/c)^{a'}\varrho^{b'}v^{2a'+b'}(1+O(|v|^{-1}))$ with
$\varrho=N_2/(N_1-N_2)$. Hence at the zeros
\[
 v^{-W_0}P\circ\Psi\longrightarrow
 \sum_{\substack{2a'+b'=W_0\\ f\nmid P_{(a',b')}}}
 P_{(a',b')}(e^a,u_1^*,Y^*,C_2)\Bigl(\frac{Y^*-C_1}{h_1}\Bigr)^{a'}\varrho^{b'},
\]
because $p/c^*=(Y^*-C_1)/h_1$. Each displayed coefficient
$P_{(a',b')}(e^a,u_1^*,Y^*,C_2)$ is nonzero, by the resultant argument at the
independent point $(e^a,u_1^*,C_2)$. The limit is a polynomial in $h_1^{-1}$
with coefficients in $\overline{k(e^a,u_1^*,C_1,C_2)}$, not all zero, and
$h_1$ is transcendental over that field by Lemma~\ref{wlbase:tb:indep}(2). So the
limit is nonzero, contradicting ($\ast$). Regime (C$_2$) treats $y_2$ in the
same way. Hence $Q\in k[E,U]$.

\emph{Step 4: $P$ does not depend on $x_1,y_1,x_2,y_2$.} Now
$Q\circ\Psi=Q(e^aw^{-M},\eta_a+(v+\Delta_a)(w^M-1))$ lies in $\C[w^{\pm1},v]$,
so $P\circ\Psi=F\cdot Q\circ\Psi\in\C[w^{\pm1},v]$ has no pole along
$\Pi_1$. Write $P=\sum_{n=0}^dP_n$ with $P_n$ homogeneous of degree $n$ in
$(x_1,y_1)$, coefficients in $k[E,U,x_2,y_2]$, and $P_d\ne0$. Near a point of
$\Pi_1$ off $\Pi_2$, all coordinates of $\Psi$ other than $(B_1,\Gamma_1)$
are regular, and $(B_1,\Gamma_1)=(\beta_1,C_1)-\mu_1\,(v_1(v_1-p),h_1)$,
where $\mu_1$ has a simple pole along $\Pi_1$ ($D_1$ is linear in $v$ with
coefficient $w^{N_1}-1\ne0$ there). Thus $P\circ\Psi=\sum_ng_n\mu_1^n$ with
regular $g_n$ and $g_d=(-1)^dP_d(v_1(v_1-p),h_1;B_2,\Gamma_2,E,U)$. If $d>0$,
regularity forces $g_d|_{\Pi_1}\equiv0$. Write
$P_d=\sum_\nu c_\nu(x_2,y_2,E,U)Z_1^\nu Z_2^{d-\nu}$ in the variables
$(Z_1,Z_2)$ standing for $(x_1,y_1)$, let $e_\nu=\deg_{x_2}c_\nu$ and
$\widehat E=\max\{e_\nu+2\nu:c_\nu\ne0\}$. In regime (P$_1$), with
$\tilde b=-N_2/(N_1(N_1-N_2))$ and $\ell c_\nu$ the leading coefficient of
$c_\nu$ in $x_2$ (a polynomial in $y_2,E,U$),
\[
 \omega^{\widehat E}g_d\longrightarrow
 (-1)^dp^{\widehat E}\sum_{e_\nu+2\nu=\widehat E}
 \ell c_\nu(C_2,e^a,u_1^*)\,\tilde b^{\,e_\nu}N_1^{-2\nu}h_1^{d-\nu}.
\]
The powers of $p$ combine to $p^{\widehat E}$ because each term contributes
$p^{2\nu}$ from $v_1(v_1-p)$ and $p^{e_\nu}$ from $B_2$. The sum is a
polynomial in $h_1$ whose coefficients $\ell c_\nu(C_2,e^a,u_1^*)$ are nonzero,
by Lemma~\ref{wlbase:tb:indep}(2), and $h_1$ is transcendental over
$k(C_2,e^a,u_1^*)$. So the limit is nonzero and $g_d|_{\Pi_1}\not\equiv0$, a
contradiction. Hence $P$ does not depend on $(x_1,y_1)$; regime (P$_2$)
removes $(x_2,y_2)$. Thus $R\in k(E,U)$.

\emph{Step 5: $R\in k[E^{\pm1},U]$.} Now
$F(w,v)=R(e^aw^{-M},\eta_a+(v+\Delta_a)(w^M-1))\in\C[w^{\pm1},v]$. Let $f$ be
an irreducible factor of $Q$ not associate to $E$. Its zero set is a curve not
contained in $\{E=0\}$, nor in $\{E=e^a\}$ since $e^a\notin k$. Choose a point
$(E_0,U_0)$ on it with $E_0\notin\{0,e^a\}$ and $P(E_0,U_0)\ne0$; only finitely
many points are excluded. Take $w_0$ with $w_0^M=e^a/E_0\ne1$ and
$v_0=(U_0-\eta_a)/(w_0^M-1)-\Delta_a$. The images of $E$ and $U$ are regular
everywhere and take the values $(E_0,U_0)$ at $(w_0,v_0)$, so $F$ has a pole
there. Hence $Q$ is a monomial in $E$ and $R\in k[E^{\pm1},U]$, which proves
the theorem.\qed

\begin{remark}
Step~5 proves, modulo Theorem~\ref{wlbase:wl:poly} and (I1)--(I4), the
one-point case of \cite[Theorem~ner:main]{EB} for a positive argument: on the
winding line the coordinate $U=\Ein(a)$ moves freely, so only Laurent
polynomials survive. Theorem~\ref{wlbase:tb:main} also gives the units:
$(\Ebar\cap\F)^\times=k^\times e^{a\Z}$ for $m=2$.
\end{remark}

\section{What the winding line cannot see}

\begin{proposition}[$\gamma$-blindness on $\F$]\label{wlbase:lim:main}
For $\gamma'\in\C$ define the formal counterparts
\[
 \delta'=e\bigl(\Ein(1)-\gamma'\bigr),\qquad
 K_0'(t)=Y(t)-\Bigl(\gamma'+\log\frac t2\Bigr)I_0(t)\quad(t>0).
\]
\begin{enumerate}
\item Algebraic relations over $k$ among $E$-values, the joint independence
(I5), the family $(\Phi_s)_{s\in\C}$ on the fields $\F$, and
Theorem~\ref{wlbase:wl:poly} for every $\xi\in\F\cap\Ebar$ do not involve $\gamma$.
\item For $\gamma'=1/2$: $\gamma'\in k$ and
$\trdeg_kk(\gamma',\delta',K_0'(2)/I_0(2),C(2))=3$. For
$\gamma'=\Ein(1)-d/e$ with $d\in k^\times$: $\delta'=d$. For
$\gamma'=\beta(n)-\log(n/2)-b/I_0(n)$ with $b\in k$: $K_0'(n)=b$.
\end{enumerate}
Consider an argument whose only information about $\gamma$ consists of the
identities
\[
 \delta=e\bigl(\Ein(1)-\gamma\bigr),\qquad
 K_0(t)=Y(t)-\Bigl(\gamma+\log\frac t2\Bigr)I_0(t),
\]
and whose other inputs are among those in \textup{(1)}. It proves its
conclusion for every $\gamma'\in\C$. By \textup{(2)}, no such argument can
show that $\gamma$, $\delta$ or a specified $K_0(n)$ is transcendental, or
that $\trdeg_kk(\gamma,\delta,\kappa,C)=4$.
\end{proposition}

\begin{proof}
(1) is Proposition~\ref{wlbase:wl:family}(3) together with the fact that the other
inputs concern $E$-values only. (2) For $\gamma'=1/2$ the field
$k(e,\Ein(1),I_0(2),C(2),\beta(2))$ has transcendence degree $5$ by (I5) and
contains $\gamma',\delta',\beta(2)-\gamma',C(2)$; the first is algebraic and
the other three are independent, as $\delta'=e(\Ein(1)-1/2)$,
$\beta(2)-1/2$ and $C(2)$ are nonconstant functions of disjoint sets of
independent coordinates. The other two choices are direct substitutions.
\end{proof}

\begin{remark}[where $\gamma$ enters]\label{wlbase:lim:point}
On $\F$ the genuine embeddings are $\Phi_{\gamma+j\tau}=\Phi^{[j]}$; the
analytic construction singles out these points of the line. An argument
settling (Q1) or (Q2) through phase embeddings must use a property of the
point $s=\gamma$ that generic points of the line lack. An example is the
rigidity
\begin{equation}\label{wlbase:lim:R}
 \xi\in\E,\ \ \Phi^{[0]}(\xi)\in\C\ \Longrightarrow\ \xi\in k .
\end{equation}
It implies $\gamma\notin k$. Indeed, if $\gamma=g\in k$, then
$\delta=e(\Ein(1)-g)$ lies in $\E$ and in $\F$ for $a=1$, and by
Definition~\ref{wlbase:wl:def}
\[
 \Phi_s(\delta)=(s-g)e^X+e\bigl(\Ein(1)-s\bigr),
\]
which is the constant $\delta$ exactly at $s=g=\gamma$, that is, for
$\Phi^{[0]}$. Since $\delta$ is then transcendental by (I5), \eqref{wlbase:lim:R}
fails. Property \eqref{wlbase:lim:R} is of the same nature as the conjecture
$\mathbf D\cap\E=k$ formulated by Fischler and Rivoal \cite{FRG}, which
already excludes $\gamma\in k$: if $\gamma\in k$, then $\delta$, a value of
the Euler series, lies in $\mathbf D\cap\E$ and is transcendental.

Other inputs of \cite{EB} are $\gamma$-sensitive and lie outside
Proposition~\ref{wlbase:lim:main}: the positivity $K_0(t)>0$, used for instance in
\cite[Theorem~jden:rank]{EB}, and the critical approximation limits of
\cite[\S limits]{EB}. The latter identify $\gamma+2\kappa$ and $\gamma-2\kappa$
with limits of explicit rational sequences, with errors $O(n^{-2})$ and worse,
and $B-\gamma$ with the limit of $R_N(1/N)+\log N$. They would become decisive
only with arithmetic control of these sequences, which is not available.
\end{remark}

\begin{remark}[on (Q4)]
An effective selection of the density-one set of
\cite[Theorem~jden:density]{EB} would exhibit an explicit integer $n$ with
$K_0(n)$ transcendental, which is a statement of type (Q1). By
Proposition~\ref{wlbase:lim:main}(2) (third choice) such a statement does not follow
from the inputs listed in Proposition~\ref{wlbase:lim:main}(1) alone.
\end{remark}

\begin{remark}[three or more harmonic coordinates]\label{wlbase:lim:three}
The proof of Theorem~\ref{wlbase:tb:main} uses $m=2$ in two places. In Step~2 one
intermediate regime suffices because the quotient of the $x$-plane by $N$ is
one-dimensional; for $m$ scales one needs the regimes $w=1+cv^{-1-1/K}$,
$K=2,\dots,m$, in which the coordinates along $N,\dots,N^{\circ(K-1)}$ (with
$N^{\circ k}=(N_i^k)_i$) diverge at different rates, and terms of equal weight
can cancel. In regimes (C$_i$) and (P$_i$) the coordinates $B_l$, $l\ne i$,
diverge at a common rate with rational ratios $N_l/(N_i-N_l)$. For $m=2$ only
one such coordinate occurs, so the leading coefficients cannot cancel. For
$m\ge3$ the leading forms can vanish along these rational directions, and the
next-order terms contain the shifts
$v_l-v_i=K_0(t_l)/I_0(t_l)-K_0(t_i)/I_0(t_i)
=\beta(t_l)-\beta(t_i)-\log(t_l/t_i)$. Separating them from the harmonic
coordinates is a problem about $E$-values and logarithms of algebraic
numbers. For rationally independent scales \cite[Theorem~mbmulti:main]{EB}
already covers every $m$.
The three-harmonic case is treated in Theorem~\ref{ext:th3:main};
the obstruction discussed here remains relevant to larger full blocks.
\end{remark}

\begin{remark}[routes outside the phase embeddings]
(i) $p$-adic: $\delta$ is the Borel sum of the archimedeanly divergent series $\sum(-1)^nn!$, which
converges in every $\Q_p$. Global-relation theorems for Euler's factorial
series \cite{MZ} exclude one rational value common to all places of a large
set of primes, but say nothing at the archimedean place, and
the defining exponential and Ein series do not converge at $1$ in
the usual $p$-adic topology, so this does not supply a $p$-adic evaluation
of the displayed identity for $\gamma$.
(ii) Logarithmic forms: $\gamma$ also occurs as $-\psi(1)$ and in Euler--Lehmer
constants, as $\gamma$ plus linear forms in logarithms; Baker's theorem gives
``at most one algebraic'' statements \cite{MS}. Combining them with $E$-values
requires separating $E$-values from logarithms and $\pi$, which is open
already for $e$ and $\pi$. (iii) Mahler-type: Catalan's integral
$\gamma=\int_0^1(1+x)^{-1}\sum_{n\ge1}x^{2^n-1}\,dx$ expresses $\gamma$ as
an integral of a Mahler function rather than a value at an algebraic point.
\end{remark}

\section{Numerical checks}

The following numerical checks and table are retained from the supplied
winding-line manuscript. They used \texttt{mpmath} at 40--80 digits;
the original randomized runs were not independently replayed for this
edition. They are illustrations and are not used as proofs.
\begin{enumerate}
\item The closed forms \eqref{wlbase:wl:img} and \eqref{wlbase:wl:ein} agree with the images
computed from the coefficients $A_j,B_j,A'_j,B'_j$ of \cite[\S mbc]{EB}
(themselves the continued components of $I_0(tz)-\varepsilon K_0(tz)$ and
$\varepsilon K_0(tz)$), for $t\in\{2,4,\sqrt5,1/3\}$, $a\in\{1,3/7\}$, random
$j\in[-40,40]$ and random complex $X$, to better than $10^{-77}$. At $X=1$ the
formulas return $\beta(t),C(t)$ for complex $s$.
\item Symbolically, after $j=(s-\gamma)/\tau$ the coefficients $A_j,B_j$ and
the two coefficients of the $Y$-image are free of $\gamma$.
\item The limits of Lemma~\ref{wlbase:tb:regimes} hold numerically. For the actual
constants at $(t_1,t_2)=(2,4)$ with $(N_1,N_2)=(1,2)$ and the model exponent
$M=3$ (the asymptotics depend only on the exponents), the discrepancies are:
\end{enumerate}
\begin{center}
\begin{tabular}{@{}lcccc@{}}
\toprule
$|v|$ & (A) & (B) & (C$_1$), (C$_2$) & (P$_1$), (P$_2$)\\
\midrule
$10^6$ & $2.7\cdot10^{-6}$ & $6.5\cdot10^{-4}$ & $2.8\cdot10^{-5}$, $1.1\cdot10^{-5}$ & $9.9\cdot10^{-6}$, $2.5\cdot10^{-6}$\\
$10^9$ & $2.7\cdot10^{-9}$ & $2.1\cdot10^{-5}$ & $2.8\cdot10^{-8}$, $1.1\cdot10^{-8}$ & $9.9\cdot10^{-9}$, $2.5\cdot10^{-9}$\\
\bottomrule
\end{tabular}
\end{center}
The decay is $O(|v|^{-1})$, and $O(|v|^{-1/2})$ in regime (B), as predicted.
The same decay holds for the exponent pairs $(3,5),(7,2),(1,6),(2,3),(5,11)$
up to $|v|=10^{18}$. The leading coefficients used in Steps~3 and~4 were
checked symbolically and numerically on test polynomials, and for random
denominators the zeros of $f\circ\Psi$ appear at the points $c^*$ predicted
in Steps~2 and~3.

% Additional section. Conditional on the winding hypotheses and the
% joint-independence input explicitly recorded in the containing paper.
% Cross-reference wlbase:tb:main is the namespaced mixed two-block theorem.
\section{A finite winding criterion in bounded degree}
\label{finitewind:section}

The polynomial-winding argument can be made finite when the degrees of
its input rational function are bounded. The resulting criterion concerns
exact identities of rational functions; it is not a numerical algorithm
for deciding equality of arbitrary special-function constants.

\begin{lemma}[Exceptional specializations]\label{finitewind:lemma}
Let $A,B\in\mathbb C[w,v]$ satisfy
\[
 \deg_w A,\deg_w B\le W,\qquad
 \deg_v A,\deg_v B\le S,
 \qquad B(1,v)=b\in\mathbb C^\times.
\]
Put $F=A/B$. Either
$F\in\mathbb C[w,w^{-1},v]$, or there are at most
$2WS+2S$ distinct numbers $v_0\in\mathbb C$ for which
$F(w,v_0)\in\mathbb C[w,w^{-1}]$.
\end{lemma}

\begin{proof}
The case $A=0$ is immediate. Otherwise cancel the greatest common
polynomial divisor and write $F=N/D$ with $N,D\in\mathbb C[w,v]$
coprime. Removing a power of $w$ from $D$ does not affect whether a
specialization is Laurent; denote the remaining denominator by $D_0$.
The separate degree bounds remain at most $W,S$ after cancellation.
Because $B(1,v)$ is a nonzero constant, every polynomial factor of $B$
has nonzero constant value at $w=1$. In particular $D_0$ has no
nonconstant factor depending on $v$ alone, and $D_0(w,v_0)$ is not
identically zero at any $v_0$.

If $\deg_wD_0=0$, it follows that $D_0$ is constant, and the first
alternative holds. Suppose instead that $q=\deg_wD_0>0$. Write
$d_0(v)=D_0(0,v)$ and let $d_q(v)$ be the leading coefficient in $w$.
Both are nonzero polynomials. Since $N$ and $D_0$ are coprime,
\[
 H(v)=d_0(v)d_q(v)\operatorname{Res}_w(N,D_0)
\]
is a nonzero polynomial. The standard determinant formula for the
resultant gives
\[
 \deg_v H\le 2S+
  (\deg_wN)(\deg_vD_0)+(\deg_wD_0)(\deg_vN)
 \le 2WS+2S.
\]
If $H(v_0)\ne0$, then $D_0(w,v_0)$ has positive degree, nonzero
constant coefficient, and no common root with $N(w,v_0)$.
It therefore has a nonzero complex root at which $F(w,v_0)$ has a
pole. Consequently this specialization is not Laurent. Every Laurent
specialization must be a zero of $H$, proving the bound.
\end{proof}

Fix the positive-scale setting of the winding family, with $m\ge1$
Bessel blocks and one exponential--Ein pair. Choose a one-parameter
subgroup for which
\[
 e^{a(1-X)}=w^M,\qquad q_i=w^{N_i},\qquad M,N_i\in\mathbb Z_{>0},
\]
with the $N_i$ distinct, as in the proof of
Theorem~\ref{wlbase:tb:main}. Put $v=v_1$, $v_i=v+\Delta_i$,
$v_a=v+\Delta_a$, and write the restricted images as
\[
 E\longmapsto e^aw^{-M},\quad
 U\longmapsto\eta_a+(v+\Delta_a)(w^M-1),\quad
 x_i\longmapsto B_i,\quad y_i\longmapsto\Gamma_i,
\]
where
\begin{align*}
 D_i&=p+(v+\Delta_i)(w^{N_i}-1),\\
 B_i&=\frac{\beta_iD_i-(v+\Delta_i)(v+\Delta_i-p)
                         (w^{N_i}-1)}{D_i},\\
 \Gamma_i&=\frac{C_iD_i-h_i(w^{N_i}-1)}{D_i}.
\end{align*}
Denote this tuple by $\Psi(w,v)$.

\begin{proposition}[Finite recognition of polynomial winding]
\label{finitewind:general}
Let $R=P/Q$ with $P,Q\in k[E,U,x_1,y_1,\ldots,x_m,y_m]$,
$Q\ne0$, and total degrees at most $d\ge1$. Set
\[
 W=d\left(2M+\sum_{i=1}^mN_i\right),\qquad S=(m+1)d,
 \qquad B_d=2S(W+1).
\]
If $R(\Psi(w,v_j))$ is Laurent in $w$ at more than $B_d$ distinct
values $v_j$, then
\[
 R(\Psi(w,v))\in\mathbb C[w,w^{-1},v].
\]
Otherwise, if this polynomial-winding conclusion fails, at most $B_d$
values $v_j\in\mathbb C$ can give Laurent specializations.
\end{proposition}

\begin{proof}
Multiply both substituted polynomials by
\[
 L(w,v)=w^{Md}\prod_{i=1}^mD_i(w,v)^d.
\]
Then $A=LP(\Psi)$ and $B=LQ(\Psi)$ belong to
$\mathbb C[w,v]$. To check the bounds, consider a monomial of $P$ or
$Q$ having exponents $e,u,r_i,c_i$ on $E,U,x_i,y_i$ respectively.
The power $w^{Md}$ clears $w^{-Me}$, since $e\le d$.
The total $w$-degree is bounded by
\[
 Md+Mu+d\sum_iN_i\le W.
\]
For the $v$-degree, the $i$th denominator contributes
$d-r_i-c_i$, the numerator of $B_i^{r_i}$ contributes at most $2r_i$,
and that of $\Gamma_i^{c_i}$ at most $c_i$. Thus the bound is
\[
 md+u+\sum_ir_i\le(m+1)d=S.
\]
Finally, at $w=1$ the image tuple equals the original tuple $T$,
and every $D_i$ equals $p$. Joint independence gives
\[
 B(1,v)=p^{md}Q(T)\ne0,
\]
independently of $v$. Lemma~\ref{finitewind:lemma} applies.
\end{proof}

\begin{theorem}[Finite exact test for the mixed two-block intersection]
\label{finitewind:twoblock}
Assume the phase and joint-independence inputs used in
Theorem~\ref{wlbase:tb:main}. Let
\[
 \xi=R(e^a,\eta_a,\beta_1,C_1,\beta_2,C_2),
 \qquad \deg P,\deg Q\le d,\qquad d\ge1,
\]
and choose positive subgroup exponents $M,N_1,N_2$ with $N_1\ne N_2$.
Set
\[
 B_d=6d\bigl(d(2M+N_1+N_2)+1\bigr).
\]
For any $B_d+1$ distinct integers $j$, the following are equivalent:
\begin{enumerate}
\item $\xi$ is integral over $\mathbf E$;
\item $\xi\in k[e^a,e^{-a},\eta_a]$;
\item every restricted image $R(\Psi(w,v_j))$ in the selected finite
      list is Laurent in $w$, where
      $v_j=\beta_1-\log(t_1/2)-\gamma-2\pi i j$.
\end{enumerate}
If $\xi$ lies outside the displayed ring, at most $B_d$ windings
in the entire integer sequence can pass this Laurent test.
\end{theorem}

\begin{proof}
The implication (2)$\Rightarrow$(1) follows because $e^{\pm a}$ and
$\eta_a$ are $E$-values. Integrality implies polynomial winding, hence
(1)$\Rightarrow$(3). The parameters $v_j$ are distinct, so
Proposition~\ref{finitewind:general} gives
$R\circ\Psi\in\mathbb C[w,w^{-1},v]$ from (3).
Steps 1--5 of the proof of Theorem~\ref{wlbase:tb:main} establish
precisely the algebraic implication
\[
 R\circ\Psi\in\mathbb C[w,w^{-1},v]
       \quad\Longrightarrow\quad R\in k[E,E^{-1},U].
\]
Those steps use joint independence and the displayed subgroup formulas;
after obtaining polynomial winding they use no further integrality
hypothesis. Thus (3)$\Rightarrow$(2). The last assertion follows from
the contrapositive and the exceptional-specialization bound.
\end{proof}

\begin{remark}
For instance, $a=1$, $t_1=1/2$, $t_2=1$ admit
$(M,N_1,N_2)=(1,1,2)$. For rational functions with numerator and
denominator of total degree at most one, $B_1=36$: any $37$ distinct
windings give a sufficient exact test. This conservative bound is not
claimed to be optimal. The theorem neither singles out Euler's constant
arithmetically nor turns finite-precision pole checks into exact proofs.
\end{remark}
\section{Inputs, scope, and comparison with the preceding notes}
Put $k=\overline{\mathbb Q}$. Let $\mathbf E$ denote the ring of values of
E-functions at algebraic points, and let $\overline{\mathbf E}$ be its
integral closure in $\mathbb C$. Here ``integral'' means satisfying a monic
polynomial with coefficients in $\mathbf E$; it does not mean being a
rational integer. For $t\in k^\times$ write
\[
 Y(t)=\sum_{n\ge0}\frac{H_n(t/2)^{2n}}{(n!)^2},\qquad
 \beta(t)=\frac{Y(t)}{I_0(t)},\qquad C(t)=\frac{I_1(t)}{I_0(t)}.
\]
Here $H_0=0$ and $H_n=\sum_{j=1}^n1/j$. The entire integral is
$\Ein(z)=\int_0^z(1-e^{-u})\,du/u$.

The intersection results below use the continued phase embeddings, their
pole criterion, the exact Bessel component formulas, and joint
independence from \cite{EB}. These are assumptions imported from that
research manuscript, not consequences of numerical tests. The
polynomial-winding principle in \cite{WL} is the additional organizing
input. Baker's theorem is the extra published arithmetic input used for
three harmonic coordinates. The resultant result from \cite{BR} enters
only the fractional-order derivative-ratio proposition in the
derivative-coordinate section.

There is a necessary correction to the novelty description in \cite{WL}.
The two-block field $k(\beta(2),C(2),\beta(4),C(4))$ is already classified
in the current version of \cite[Theorems \texttt{tbr:main},
\texttt{ctb:main}]{EB}. The additional intersection proved in \cite{WL}
retains $e^a$ and $\Ein(a)$ together with two full positive-scale Bessel
blocks. The three-harmonic theorem and the mixed affine-projection theorem
below extend in different directions; neither implies the full mixed
three-block classification.

\paragraph{Polynomial winding.} We use Theorem~\ref{wlbase:wl:poly}.

The proof in \cite{WL} uses the genuine integer winding values to rule
out every denominator depending on the torus variables. Evaluation at
the identity of the torus returns the original algebraically independent
generators for every $s$, and rules out the remaining denominators in
$s$. This is a necessary consequence of integrality, not a proved
converse. The complex-scale version is detailed in the affine-projection
section below.

\paragraph{Mixed two-block intersection.} We use Theorem~\ref{wlbase:tb:main}.

All conclusions here retain the original arithmetic coefficient field
$k$, except for the expressly displayed exponential--Ein coefficients.
No general separation theorem for E-values and logarithms is assumed.

% Research fragment.  Uses the winding-line theorem wl:poly and the
% joint-independence input (I5) of upload/winding_line.tex.
% Labels have the prefix th3:.  Bibliography key BakerIII is supplied below.
\section{Three harmonic coordinates at arbitrary positive scales}

The third harmonic coordinate can be admitted without any rational
independence assumption on the scales.  Here the coefficient field remains
$k=\overline{\mathbb Q}$; no derivative quotient or additional $E$-value is
adjoined to the field in the following theorem.

\begin{theorem}\label{ext:th3:main}
Let $t_1,t_2,t_3$ be distinct positive algebraic numbers.  Then
\[
 \overline{\mathbf E}\cap
 k\bigl(\beta(t_1),\beta(t_2),\beta(t_3)\bigr)=k.
\]
In particular, this holds for the three dependent scales $2,4,8$.
\end{theorem}

We use precisely two inputs from the winding-line construction: the
algebraic independence of the three numbers $\beta(t_i)$ over $k$, and
Theorem~\ref{wlbase:wl:poly}, which makes the image of an integral element a
Laurent polynomial in the phase coordinates and a polynomial in the
winding coordinate.  The additional arithmetic input is Baker's theorem
on linear forms in logarithms \cite{BakerIII}.

\subsection{Phase restriction and the three charts}
Put $p=i\pi$, $\beta_i=\beta(t_i)$, $\ell_i=\log(t_i/2)$ and
$v_i=\beta_i-\ell_i-s$.  The harmonic image is
\begin{equation}\label{ext:th3:B}
 B_i=\beta_i-\frac{v_i(v_i-p)(q_i-1)}{p+v_i(q_i-1)}
     =s+\ell_i+\frac{p v_iq_i}{p+v_i(q_i-1)}.
\end{equation}
As in the one-parameter subgroup restriction in the proof of
Theorem~\ref{wlbase:tb:main}, choose distinct positive integers $N_i$ such that
$q_i=w^{N_i}$.  This is possible even when the original scales are not
commensurable: if $2t_i=\sum_j n_{ij}\lambda_j$ in an integral basis of
their additive group, choose integers $u_j$ nearest to $L\lambda_j$.
Then $N_i=\sum_jn_{ij}u_j=2Lt_i+O(1)$ are positive and distinct for large
$L$.  Multiplying all the $u_j$ by $2$, if necessary, also ensures
$N_i\ge2$.

Set $v=v_1$ and $\Delta_i=v_i-v_1$, so $v_i=v+\Delta_i$ with constant
$\Delta_i$.  Write
\[
 \mathcal D_i=p+(v+\Delta_i)(w^{N_i}-1).
\]
For an integral element $\xi=R(\beta_1,\beta_2,\beta_3)$, the restriction
of the polynomial winding theorem gives
\begin{equation}\label{ext:th3:Laurent}
 R(B_1,B_2,B_3)\in\mathbb C[w,w^{-1},v].
\end{equation}
In particular, if $R=P/Q$ is reduced over $k$, there is no point
$(w,v)\in\mathbb C^\times\times\mathbb C$ at which every $\mathcal D_i$
is nonzero and $Q(B)=0$ but $P(B)\ne0$.

Choose rational coefficients $a_{ri}$, $1\le r,i\le3$, such that
\[
 \sum_{i=1}^3 a_{ri}N_i^j=\delta_{rj}\qquad(1\le r,j\le3),
\]
and put $D_r(x)=\sum_i a_{ri}x_i$, $d_r=D_r(\beta)$ and
$\widetilde D_r=D_r(B)$.  This invertible change of coordinates is defined
over $k$, and $d_1,d_2,d_3$ are algebraically independent over $k$.

\begin{lemma}\label{ext:th3:charts}
For $K\in\{1,2,3\}$ put $v=H^K$ and
$w=1+pzH^{-K-1}$, with $H\to+\infty$.  Uniformly for $z$ in a compact
subset of $\mathbb C$,
\[
 B_i=\beta_i+\sum_{j=1}^K(-1)^jN_i^jz^jH^{K-j}+O(H^{-1}),
 \qquad \mathcal D_i\longrightarrow p.
\]
Consequently
\[
 \widetilde D_r=
 \begin{cases}
 d_r+(-1)^rz^rH^{K-r}+O(H^{-1}),&r\le K,\\
 d_r+O(H^{-1}),&r>K.
 \end{cases}
\]
\end{lemma}

\begin{proof}
Let $\theta_i=w^{N_i}-1$.  Then
$\theta_i=N_ipzH^{-K-1}+O(H^{-2K-2})$ and
$(v+\Delta_i)\theta_i/p=N_izH^{-1}+O(H^{-K-1})$.
In the first expression in \eqref{ext:th3:B}, expand
$(1+(v+\Delta_i)\theta_i/p)^{-1}$ through degree $K-1$ in $H^{-1}$.
The factor $(v+\Delta_i)(v+\Delta_i-p)\theta_i/p$ equals
$N_izH^{K-1}+O(H^{-1})$.  Terms of order $K$ or higher in the geometric
expansion contribute $O(H^{-1})$; the other errors have the same bound.
This proves the displayed expansion and then the assertions about the
dual Vandermonde coordinates.
\end{proof}

\subsection{Eliminating every denominator}
\begin{lemma}\label{ext:th3:denominator}
If $R\in k(x_1,x_2,x_3)$ satisfies \eqref{ext:th3:Laurent}, then
$R\in k[x_1,x_2,x_3]$.
\end{lemma}

\begin{proof}
Write $R=P/Q$ in reduced form, now using the coordinates $D_1,D_2,D_3$.
We may suppose $P\ne0$.  Every irreducible factor $f$ of $Q$ is coprime
to $P$.

\emph{First chart.}
Suppose that $f$ depends on $D_1$.  The $K=1$ chart has limit
$(d_1-z,d_2,d_3)$, and
$f(d_1-z,d_2,d_3)$ is a nonconstant polynomial in $z$.
Let $z_*$ be a root.  The algebraic set $Z(P)\cap Z(f)\subset\mathbb A^3$
has dimension at most $1$.  Its translates in the $D_1$ direction have
Zariski closure of dimension at most $2$, defined over $k$.
Since $d=(d_1,d_2,d_3)$ is generic over $k$, that closure does not contain
$d$.  Therefore $P(d_1-z_*,d_2,d_3)\ne0$.
Uniform convergence in Lemma~\ref{ext:th3:charts} and Hurwitz's theorem
give zeros $z_H\to z_*$ of $f(\widetilde D)$; at them $P(\widetilde D)\ne0$
and all $\mathcal D_i\ne0$ for large $H$.  This contradicts
\eqref{ext:th3:Laurent}.  Thus $Q\in k[D_2,D_3]$.

\emph{Second chart.}
Suppose that an irreducible factor $f\in k[D_2,D_3]$ of $Q$ has positive
degree in $D_2$.  Its specialization $f(W,d_3)$ has a root $W_*$ algebraic
over $k(d_3)$.  As $d_2$ is transcendental over $k(d_3)$, choose a nonzero
$z_*$ with $z_*^2=W_*-d_2$.  The $K=2$ chart has
\[
 \widetilde D_1=-zH+O(1),\qquad
 \widetilde D_2\longrightarrow d_2+z^2,\qquad
 \widetilde D_3\longrightarrow d_3.
\]
Again Hurwitz gives zeros $z_H\to z_*$ of $f(\widetilde D_2,\widetilde D_3)$.
Write $P=\sum_n P_n(D_2,D_3)D_1^n$ and take the largest $n_0$ such that
$f\nmid P_{n_0}$.  At these zeros all terms with $f\mid P_n$ vanish
exactly.  The resultant of $f$ and $P_{n_0}$ in $D_2$ is a nonzero
polynomial in $D_3$: irreducibility and Gauss's lemma ensure that the
two polynomials are coprime over $k(D_3)$.  Its value at $d_3$ is nonzero,
so $P_{n_0}(W_*,d_3)\ne0$.  Hence
\[
 H^{-n_0}P(\widetilde D)\longrightarrow
 (-z_*)^{n_0}P_{n_0}(W_*,d_3)\ne0
\]
along the zeros, a contradiction.  It follows that $Q\in k[D_3]$.

\emph{Third chart.}
Because $k$ is algebraically closed, an irreducible nonconstant factor
of $Q$ has the form $f=D_3-\alpha$, with $\alpha\in k$.
Choose $z_*$ with $z_*^3=d_3-\alpha$.  This number is nonzero and
transcendental over $k$.  The $K=3$ chart gives
\[
 \widetilde D_1=-zH^2+O(1),\qquad
 \widetilde D_2=z^2H+O(1),\qquad
 \widetilde D_3\longrightarrow d_3-z^3.
\]
Hurwitz gives zeros $z_H\to z_*$ with $\widetilde D_3=\alpha$ exactly.
The nonzero reduction $P(D_1,D_2,\alpha)=\sum_{a,b}c_{ab}D_1^aD_2^b$
has coefficients $c_{ab}\in k$.  Let
$W=\max\{2a+b:c_{ab}\ne0\}$.  Then along those zeros
\[
 H^{-W}P(\widetilde D)\longrightarrow
 G(z_*),\qquad
 G(Z)=\sum_{2a+b=W}c_{ab}(-1)^aZ^{a+2b}.
\]
For fixed $2a+b=W$, the exponent $a+2b=2W-3a$ determines $a$ and $b$.
Thus the displayed monomials have distinct exponents and $G$ is a nonzero
polynomial over $k$.  Transcendence of $z_*$ implies $G(z_*)\ne0$,
again a contradiction.  Therefore $Q$ is constant.
\end{proof}

\subsection{The primitive-root pole and Baker's theorem}
\begin{lemma}\label{ext:th3:Baker}
If $r>0$ is algebraic and $r\ne1$, then
$\log r+i\pi A$ is transcendental for every $A\in k$.
\end{lemma}

\begin{proof}
The real number $\log r\ne0$ and the purely imaginary logarithm
$\log(-1)=i\pi$ are linearly independent over $\mathbb Q$.
Baker's theorem makes $1,\log r,\log(-1)$ linearly independent over $k$;
equivalently, a nonzero algebraic linear form in these two logarithms
cannot be algebraic \cite{BakerIII}.  The proposed form has coefficient
$1$ on $\log r$, so it is transcendental.
\end{proof}

\begin{lemma}\label{ext:th3:numerator}
If $P\in k[x_1,x_2,x_3]$ and
$P(B_1,B_2,B_3)\in\mathbb C[w,w^{-1},v]$, then $P$ is constant.
\end{lemma}

\begin{proof}
Suppose otherwise.  Among the variables on which $P$ depends choose
$x_i$ with maximal exponent $N_i$, and write
\[
 P=\sum_{d=0}^{d_0}c_d(x_j,x_k)x_i^d,
 \qquad d_0>0,\quad c_{d_0}\ne0,
\]
omitting from $c_d$ any variable on which $P$ does not depend.
Consider the irreducible pole divisor
\[
 \Pi_i:\quad \mathcal D_i=0,\qquad
 v_i=-\frac{p}{w^{N_i}-1},\qquad
 s=\beta_i-\ell_i+\frac{p}{w^{N_i}-1}.
\]
Its coordinate ring is the localization of $\mathbb C[w,w^{-1}]$ at
$w^{N_i}-1$.  At its generic point $B_i$ has a simple pole, with nonzero
coefficient, while the other coordinates occurring in $P$ are regular.
We prove that $c_{d_0}(B_j,B_k)$ does not vanish identically on $\Pi_i$.

Choose a primitive $N_i$-th root of unity $\zeta$.  Since $N_i\ge2$,
$\zeta\ne1$.  Each other exponent occurring in $P$ is strictly smaller
than $N_i$, so $\zeta^{N_l}\ne1$ for these indices $l$.
On $\Pi_i$, as $w\to\zeta$, the variable $s$ tends to infinity and
\eqref{ext:th3:B} gives
\begin{equation}\label{ext:th3:boundary}
 B_l=s+\ell_l+
 p\frac{\zeta^{N_l}}{\zeta^{N_l}-1}+o(1).
\end{equation}
In particular, the other denominators are not identically zero on
$\Pi_i$.  If both other variables occur, their difference tends to
\[
 L_{jk}=\log\frac{t_j}{t_k}
  +p\left(\frac{\zeta^{N_j}}{\zeta^{N_j}-1}
           -\frac{\zeta^{N_k}}{\zeta^{N_k}-1}\right).
\]
The parenthesized number is algebraic.  Since $t_j/t_k\ne1$,
Lemma~\ref{ext:th3:Baker} implies that $L_{jk}$ is transcendental.
Write the nonzero polynomial $c_{d_0}$ uniquely as
\[
 c_{d_0}(x_j,x_k)=\sum_{b=0}^{b_*}H_b(x_j-x_k)x_k^b,
 \qquad H_{b_*}\ne0.
\]
By \eqref{ext:th3:boundary},
\[
 B_k^{-b_*}c_{d_0}(B_j,B_k)\longrightarrow H_{b_*}(L_{jk})\ne0.
\]
Thus this leading coefficient is not identically zero on $\Pi_i$.
If only one other variable occurs, \eqref{ext:th3:boundary} makes it tend
to infinity, which gives the same conclusion for a nonzero univariate
polynomial.  If neither occurs, the conclusion is immediate.

At the generic point of $\Pi_i$, the term
$c_{d_0}(B_j,B_k)B_i^{d_0}$ therefore has a pole of order $d_0$,
whereas every lower term has smaller pole order.  This contradicts
$P(B)\in\mathbb C[w,w^{-1},v]$, since $\Pi_i$ lies in
$\mathbb C^\times\times\mathbb C$.  Consequently $P$ is constant.
\end{proof}

\begin{proof}[Proof of Theorem~\ref{ext:th3:main}]
An element of the displayed intersection has a unique expression
$R(\beta_1,\beta_2,\beta_3)$ over $k$ by joint independence.
The polynomial winding theorem yields \eqref{ext:th3:Laurent}.
Lemma~\ref{ext:th3:denominator} makes $R$ a polynomial and
Lemma~\ref{ext:th3:numerator} makes it constant.  The reverse inclusion holds
because algebraic numbers are $E$-values.
\end{proof}

\begin{remark}\label{ext:th3:scope}
The theorem settles the three-coordinate harmonic subfield of the
dependent-scale integral-intersection question.  It does not classify
the field obtained by adjoining all three derivative quotients $C(t_i)$,
nor does it extend the coefficient field to arbitrary $E$-values.
For four harmonic coordinates the present proof has two distinct
obstacles: the fourth chart permits equal leading monomials such as
$D_1D_3$ and $D_2^2$, and the primitive-root boundary presents two
logarithmic differences whose algebraic independence does not follow
from Baker's theorem.  No claim about the arithmetic nature of
$\gamma$, $\delta$, or an individual $K_0(t)$ is made.
\end{remark}

% Add to the containing bibliography:
% \bibitem{BakerIII} A.~Baker, Linear forms in the logarithms of algebraic
% numbers (III), Mathematika \textbf{14} (1967), 220--228.
% \href{https://doi.org/10.1112/S0025579300003843}
% {doi:10.1112/S0025579300003843}.

% Internal proof fragment, 22 October 2026.
% New synthesis of winding_line.tex and affine_rational_closure.tex.
% Not an edit to either attachment or the authoritative manuscript.
% Dependencies: cf:joint, continued phase/pole lemmas, cru:coefficient,
% and the one-point exponential--Ein exact intersection.

\section{An exponential--Ein field with an arbitrary affine Bessel projection}
\label{ext:mae:section}

The following result permits arbitrarily many dependent complex scales,
but retains only one affine projection of their Bessel coordinates.  It
does not assert the corresponding full multivariate intersection.

\begin{theorem}\label{ext:mae:main}
Let $k=\overline{\mathbb Q}$, let $a>0$ be algebraic, and set
$E=e^a$ and $\eta=\operatorname{Ein}(a)$.  Let
$t_1,\ldots,t_m\in k^\times$ be distinct modulo sign.  For
$\ell_0,u_i,v_i\in k$ put
\[
 L=\ell_0+\sum_{i=1}^m
             \bigl(u_i\beta(t_i)+v_iC(t_i)\bigr).
\]
Assume that at least one pair $(u_i,v_i)$ is nonzero.  Then
\[
 \boxed{\quad
 \overline{\mathbf E}\cap k(E,\eta,L)
                  =k[E,E^{-1},\eta].\quad}
\]
There is no rational-independence condition on $a,t_1,\ldots,t_m$.
\end{theorem}

We first record why polynomial winding remains available for complex
scales.  Choose a common nonexceptional summation direction.  The exact
complex component normalization gives constants $p_i\in\{\pi i,-\pi i\}$.
With $q_i=e^{2t_i(1-X)}$ the continued images have the form
\begin{align}
 B_i&=\beta_i-
       \frac{b_i(b_i-p_i)(q_i-1)}{p_i+b_i(q_i-1)},
 &\Gamma_i&=C_i-
       \frac{h_i(q_i-1)}{p_i+b_i(q_i-1)},
 \label{ext:mae:complex-images}\\
 E&\longmapsto E e^{-a(1-X)},
 &\eta&\longmapsto\eta+b_a\bigl(e^{a(1-X)}-1\bigr),
 \label{ext:mae:ein-images}
\end{align}
where $h_i=1/(t_iI_0(t_i)^2)\ne0$.  On the genuine winding sequence,
all $b_i$ and $b_a$ are a common affine parameter plus a fixed constant:
$b_i=v+\Delta_i$, $b_a=v+\Delta_a$, with $v=-2\pi i j+O(1)$.
In particular, the signs $p_i$ do not change the slope of the winding
parameter.  These are the normalized complex component formulas; the
constants $\Delta_i$ are not assumed to be algebraically independent.
A change of the common summation direction can add a fixed branch
constant to $\Delta_a$; it does not change the common winding slope.

Let $H=\mathbb Za+\sum_i\mathbb Z(2t_i)$, and choose a basis of this
finite-rank free abelian group.  Write its normalized characters as
$z_1,\ldots,z_d$.  Substitution in \eqref{ext:mae:complex-images}--
\eqref{ext:mae:ein-images} defines a rational function $F_\xi(z,v)$ for every
$\xi\in k(E,\eta,\beta_i,C_i:1\le i\le m)$.  Evaluation at $z=(1,\ldots,1)$
returns the original generators, independently of $v$.  Joint independence
therefore makes this a field embedding for every fixed complex $v$.
If $\xi$ is integral over $\mathbf E$, the polynomial-winding argument
gives
\begin{equation}
             F_\xi\in\mathbb C[z_1^{\pm1},\ldots,z_d^{\pm1},v].
 \label{ext:mae:poly-winding}
\end{equation}
For completeness, each genuine winding specialization lies in
$\mathbb C[H]$: the phase-image pole lemma gives membership in the full
group algebra, and separation by $H$-cosets gives its intersection with
$\operatorname{Frac}\mathbb C[H]$.  A reduced denominator of $F_\xi$ with
a nonconstant torus factor would remain an uncancelled denominator on
all but finitely many winding slices, a contradiction.  Its remaining
factors depend only on $v$.  None can vanish at $v_0$, because the
original rational substitution at $v_0$ is well defined generically in
$z$, its denominator being nonzero at $z=(1,\ldots,1)$.  This proves
\eqref{ext:mae:poly-winding}.  This argument uses no reality assumption on
the phase lattice.

\begin{proof}[Proof of Theorem~\ref{ext:mae:main}]
Discard inactive blocks.  The numbers $E,\eta,L$ are algebraically
independent by the joint independence theorem, since $L$ is a nonconstant
algebraic affine combination of the Bessel coordinates.  Write an integral element of their field
as $R(E,\eta,L)$ with $R=P/Q$ and coprime
$P,Q\in k[Z,U,T]$.

\emph{A suitable one-dimensional phase restriction.}
Choose a homomorphism $\ell:H\to\mathbb Z$ avoiding the finitely many
hyperplanes defined by $a$, $2t_i$, and $2t_i\pm2t_l$.
Replace $\ell$ by $-\ell$ if needed so that $M=\ell(a)>0$.
Replace each $t_i$ by $-t_i$ when needed, using
$\beta(-t_i)=\beta(t_i)$ and $C(-t_i)=-C(t_i)$ and changing the
corresponding $v_i$.  We may then assume
\[
        N_i=\ell(2t_i)>0,\qquad N_i\ne N_l\quad(i\ne l).
\]
Choose the common nonexceptional direction after these replacements.
Restrict the phase characters by $e^{h(1-X)}\mapsto w^{\ell(h)}$.
In the resulting formulas put
\begin{align*}
 D_i(w,v)&=p_i+(v+\Delta_i)(w^{N_i}-1),\\
 B_i(w,v)&=\beta_i-
 \frac{(v+\Delta_i)(v+\Delta_i-p_i)(w^{N_i}-1)}{D_i(w,v)},\\
 \Gamma_i(w,v)&=C_i-\frac{h_i(w^{N_i}-1)}{D_i(w,v)},\\
 \mathcal E(w)&=Ew^{-M},\qquad
 \mathcal U(w,v)=\eta+(v+\Delta_a)(w^M-1),\\
 \mathcal L(w,v)&=\ell_0+\sum_i(u_iB_i+v_i\Gamma_i).
\end{align*}
The rational substitution remains injective, since its value at $w=1$
is the original tuple for every $v$.  By \eqref{ext:mae:poly-winding},
\begin{equation}
 R(\mathcal E,\mathcal U,\mathcal L)
                         \in\mathbb C[w,w^{-1},v].
 \label{ext:mae:regular}
\end{equation}
It consequently has no pole on $\mathbb C^\times\times\mathbb C$.

\emph{A bounded affine-coordinate chart when a harmonic coefficient occurs.}
First suppose that some $u_i$ is nonzero.
The numbers $N_i/p_i$ are nonzero and pairwise distinct: the $N_i$ are
distinct positive integers and $p_i/p_l\in\{1,-1\}$.
Vandermonde invertibility shows that some integer $r$, $1\le r\le m$,
satisfies
\[
              A_r:=\sum_i u_i(N_i/p_i)^r\ne0.
\]
Choose the least such $r$.  Let $v$ tend to positive real infinity and
put $w=1+c v^{-1-1/r}$.  Uniformly for $c$ in a compact set,
\begin{align}
 B_i&=\beta_i+
       \sum_{k=1}^{r}(-1)^k(N_i c/p_i)^k v^{1-k/r}+o(1),
 &\Gamma_i&\longrightarrow C_i.
 \label{ext:mae:chart}
\end{align}
Indeed, if $z_i=(v+\Delta_i)(w^{N_i}-1)/p_i$, then
$z_i=(N_i c/p_i)v^{-1/r}+O(v^{-1-1/r})$ and
$B_i-\beta_i=-(v+\Delta_i-p_i)z_i/(1+z_i)$.
The geometric expansion through degree $r$ has error $O(v^{-1/r})$;
the fixed shifts contribute only $o(1)$.  The denominators $D_i$ tend
uniformly to $p_i$.  The minimality of $r$ now gives
\begin{equation}
 \mathcal E\longrightarrow E,\qquad
 \mathcal U\longrightarrow\eta,\qquad
 \mathcal L\longrightarrow L+(-1)^r A_r c^r.
 \label{ext:mae:chart-limit}
\end{equation}

Suppose that $Q$ has positive degree in $T$.  Coprimality over
$k(Z,U)$ makes $\operatorname{Res}_T(P,Q)$ a nonzero rational function
of $Z,U$.  Joint independence of $E,\eta$ shows that the specialized
polynomials $P(E,\eta,T)$ and $Q(E,\eta,T)$ remain coprime and that
$Q(E,\eta,T)$ retains its positive degree.  Choose a root $\alpha$
of the latter.  Then $P(E,\eta,\alpha)\ne0$, and $\alpha\ne L$
because $L$ is transcendental over $k(E,\eta)$.
Choose $c_0\ne0$ such that
$L+(-1)^r A_r c_0^r=\alpha$.  On a sufficiently small disc about
$c_0$, the limit denominator in \eqref{ext:mae:chart-limit} has an isolated
zero, while the limit numerator is nonzero.  Uniform convergence and
Rouch\'e's theorem produce, for all sufficiently large $v$, a zero
of the substituted denominator in this disc at which its numerator
and all $D_i$ are nonzero.  This is a pole at a point with $w\ne0$,
contrary to \eqref{ext:mae:regular}.  Consequently $Q$ has degree zero in $T$
in this case.

\emph{The pure derivative-coordinate case.}
Now suppose that every $u_i=0$, and choose $i$ with $v_i\ne0$.
Set $u_i^*=\eta-Mp_i/N_i$.  The three numbers $E,u_i^*,L$ are
algebraically independent.  To prove this without assuming that the
complex derivative coordinates are real, enlarge the coefficient field to
\[
 F_0=k\bigl(E,C(t_l),C(\overline{t_l}):1\le l\le m\bigr).
\]
The coefficients of the Bessel functions are real, so this field is
stable under complex conjugation.  Repeated sign classes in the displayed
family are identified using $C(-t)=-C(t)$.  Joint independence makes
$\eta$ transcendental over $F_0$.  If $u_i^*$ were algebraic over
$F_0$, its conjugate would be algebraic as well, and their average
$\eta$ would be algebraic over $F_0$, a contradiction.  Thus $u_i^*$
is transcendental over $F_0$, and in particular over $k(E,L)$.
The numbers $E,L$ are independent, proving the assertion.

For $v\to+\infty$, put $b_i=v+\Delta_i$ and choose the branch near $1$ of
\[
 w=\left(1-\frac{p_i}{b_i}+\frac{c}{b_i^2}\right)^{1/N_i}.
\]
Uniformly for $c$ in a compact subset of $\mathbb C^\times$,
\begin{align*}
 D_i&=c/b_i,&
 \Gamma_i&\longrightarrow C_i+\frac{p_i h_i}{c},\\
 D_l&\longrightarrow p_l-\frac{N_lp_i}{N_i}\ne0,
 &\Gamma_l&\longrightarrow C_l\quad(l\ne i),\\
 \mathcal E&\longrightarrow E,&
 \mathcal U&\longrightarrow u_i^*,\qquad
 \mathcal L\longrightarrow L+\frac{v_i p_i h_i}{c}.
\end{align*}
If $Q$ has positive degree in $T$, specialize its resultant with $P$
at the independent pair $(E,u_i^*)$.  A root $\alpha$ of
$Q(E,u_i^*,T)$ is not a root of $P(E,u_i^*,T)$, and $\alpha\ne L$
by the independence just proved.  Take
$c_0=v_i p_i h_i/(\alpha-L)\ne0$.  On a small disc about $c_0$,
the same uniform-convergence and Rouch\'e argument produces a zero
of the image denominator with nonzero numerator.  All $D_l$ are
nonzero there for large $v$, including $D_i=c/b_i$, and $w\ne0$.
This contradicts \eqref{ext:mae:regular}.  Thus, in both cases,
\begin{equation}
                              Q\in k[Z,U].
 \label{ext:mae:no-T-denominator}
\end{equation}

\emph{Eliminating positive degree in the affine coordinate.}
Suppose that $P$ has degree $d>0$ in $T$, with nonzero leading
coefficient $P_d(Z,U)$.  Choose an active index $i$, that is,
$(u_i,v_i)\ne(0,0)$.  On the divisor
$D_i=0$ take $w\to1$, $w\ne1$, and set
\[
                 v=-\Delta_i-\frac{p_i}{w^{N_i}-1}.
\]
Along this divisor,
\begin{align*}
 \mathcal E&\longrightarrow E,&
 \mathcal U&\longrightarrow u_i^*:=\eta-\frac{Mp_i}{N_i},\\
 D_l&\longrightarrow p_l-\frac{N_lp_i}{N_i}\ne0
           &&(l\ne i).
\end{align*}
The numbers $E,u_i^*$ are algebraically independent over $k$.
Indeed, $E,\eta$ are real and independent; if $u_i^*$ were algebraic
over $k(E)$, complex conjugation would make its conjugate algebraic
over the same field, and their average $\eta$ would be algebraic over
$k(E)$.  Here $Mp_i/N_i$ is purely imaginary, which is the only
reality input in this proof.  Hence
\[
                         P_d(E,u_i^*)Q(E,u_i^*)\ne0.
\]
For $w$ sufficiently close to but different from $1$, both factors
remain nonzero.  At such a point $\mathcal L$ has a genuine simple
pole along $D_i=0$: its singular numerator is
\[
 -\bigl(u_i b_i(b_i-p_i)+v_i h_i\bigr)(w^{N_i}-1),
 \qquad b_i=-\frac{p_i}{w^{N_i}-1},
\]
which is nonzero for $w$ sufficiently close to $1$: if $u_i\ne0$,
its quadratic term dominates as $|b_i|\to\infty$; otherwise
$v_i h_i\ne0$.
Other blocks are regular there.  Therefore
$P(\mathcal E,\mathcal U,\mathcal L)/Q(\mathcal E,\mathcal U)$
has a pole of order $d$, again contradicting \eqref{ext:mae:regular}.
It follows that $P\in k[Z,U]$ as well.

We have proved that the integral element belongs to $k(E,\eta)$.
The one-point exponential--Ein exact intersection gives
\[
              \overline{\mathbf E}\cap k(E,\eta)
                              =k[E,E^{-1},\eta].
\]
Conversely, all three displayed generators of this ring are E-values,
so every element of it is integral over $\mathbf E$.  This completes
the proof.
\end{proof}

\begin{corollary}[A mixed three-scale logarithm circuit]\label{ext:mae:circuit}
For
\[
 \Delta=\beta(2)-2\beta(4)+\beta(8)
 =\frac{K_0(2)}{I_0(2)}-2\frac{K_0(4)}{I_0(4)}
                         +\frac{K_0(8)}{I_0(8)},
\]
and every positive algebraic $a$, one has
\[
 \overline{\mathbf E}\cap k(e^a,\operatorname{Ein}(a),\Delta)
                      =k[e^a,e^{-a},\operatorname{Ein}(a)].
\]
In particular, a rational expression in these three independent
numbers can be integral only if its reduced expression loses all
dependence on $\Delta$ and is Laurent polynomial in $e^a$ and
polynomial in $\operatorname{Ein}(a)$.
\end{corollary}

\begin{remark}[Scope]
The coefficients of the affine Bessel combination are algebraic.
The rational function subsequently applied to it may have arbitrary
coefficients in $k(e^a,\operatorname{Ein}(a))$, as the theorem states.
This is stronger than exclusion in $k(L)$ alone and covers any finite
number of harmonic coordinates, including mutually dependent complex
scales.  It retains just one affine coordinate; it does not classify
$k(e^a,\operatorname{Ein}(a),\beta(t_i),C(t_i):1\le i\le m)$ for
$m\ge3$.  The pure derivative-coordinate case is included, even at
complex scales, by using the conjugation-stable enlarged field in the
proof.  For positive derivative scales that case was already implied
by the existing arbitrary-derivative-scale theorem.
\end{remark}

% Internal proof fragment, 22 October 2026.
% Relative strengthening of mixed_affine_ein_closure.tex.
% Does NOT classify the integral intersection of the coefficient field.

\section{Affine Bessel exclusion over the positive exponential-integral field}
\label{ext:ape:section}

\begin{lemma}[A common imaginary translation of real independent numbers]
\label{ext:ape:imaginary-translation}
Let $F\subseteq\mathbb C$ be a field of characteristic zero that is stable
under complex conjugation.  Let $x_1,\ldots,x_n$ be real and algebraically
independent over $F$.  Let $T$ be any real number, and let
$c_1,\ldots,c_n\in F\cap\mathbb R$.  Then
\[
                     x_1+iTc_1,\ldots,x_n+iTc_n
\]
are algebraically independent over $F$.
\end{lemma}

\begin{proof}
The case $c=0$ is immediate.  Otherwise choose an invertible matrix
$M$ over $F\cap\mathbb R$ sending the vector $c$ to the first standard
basis vector, and put $y=Mx$.  The $y_i$ are real and independent over
$F$, while
\[
                    M(x+iTc)=(y_1+iT,y_2,\ldots,y_n).
\]
The field $F(y_2,\ldots,y_n)$ is stable under conjugation.  A dependence
of the displayed tuple would make $y_1+iT$ algebraic over that field,
because $y_2,\ldots,y_n$ are independent.  Conjugation would also make
$y_1-iT$ algebraic over it.  Their average $y_1$ would then be algebraic,
contrary to the independence of the $y_i$.  An invertible linear change
of coordinates preserves algebraic independence, completing the proof.
\end{proof}

Put
\[
 K_{\exp}=\overline{\mathbb Q}(e^\rho:\rho\in\overline{\mathbb Q}),
 \qquad
 \mathcal F_+=K_{\exp}
       \bigl(\operatorname{Ein}(a):a\in\overline{\mathbb Q},\ a>0\bigr).
\]

\begin{theorem}[Relative affine exclusion]\label{ext:ape:main}
Let $t_1,\ldots,t_m\in k^\times$, $k=\overline{\mathbb Q}$, be distinct
modulo sign, and let
\[
 L=\ell_0+\sum_{i=1}^m
              \bigl(u_i\beta(t_i)+v_iC(t_i)\bigr),
 \qquad \ell_0,u_i,v_i\in k,
\]
be nonconstant.  Then
\[
 \boxed{\quad
 \overline{\mathbf E}\cap\mathcal F_+(L)
                     =\overline{\mathbf E}\cap\mathcal F_+.
 \quad}
\]
Equivalently, if a rational function of $L$ with coefficients in
$\mathcal F_+$ is integral over $\mathbf E$, its reduced rational
expression is constant in $L$.
\end{theorem}

\begin{proof}
One inclusion is immediate.  An element of $\mathcal F_+(L)$ uses only
finitely many distinct positive arguments $a_1,\ldots,a_n$ and finitely
many exponential values.  Enlarge their exponential character lattice
to contain the $a_l$ and all $2t_i$, and choose a basis
$\rho_1,\ldots,\rho_d$.  Put
\[
                 E_j=e^{\rho_j},\qquad
                 \eta_l=\operatorname{Ein}(a_l).
\]
The $E_j$ are algebraically independent by the exponential theorem,
and the joint theorem makes the $\eta_l$ and $L$ independent over
$K_{\exp}$.  Write the proposed integral element as
$R(E,\eta,L)=P(E,\eta,L)/Q(E,\eta,L)$, with coprime
\[
 P,Q\in k[Z_1^{\pm1},\ldots,Z_d^{\pm1},U_1,\ldots,U_n,V].
\]

The complex polynomial-winding argument applies to this larger finite
field.  It adds the exponential image $E_j e^{-\rho_j(1-X)}$ and, for
each $l$, the image
\[
       \eta_l+(v+\Delta_l)\bigl(e^{a_l(1-X)}-1\bigr).
\]
All fixed branch constants are absorbed into the $\Delta_l$.  Choose
and orient a phase-lattice homomorphism as in Theorem~\ref{ext:mae:main},
so that the Bessel exponents $N_i>0$ are pairwise distinct.  Write
\[
 M_l=\ell(a_l),\qquad R_j=\ell(\rho_j).
\]
Neither positivity nor nonvanishing of the $M_l,R_j$ is needed.  The
folded images are
\[
 \mathcal E_j=E_jw^{-R_j},\qquad
 \mathcal U_l=\eta_l+(v+\Delta_l)(w^{M_l}-1),
\]
and the Bessel images $B_i,\Gamma_i$ are exactly those in
Theorem~\ref{ext:mae:main}.  In particular,
\begin{equation}
             R(\mathcal E,\mathcal U,\mathcal L)
                         \in\mathbb C[w,w^{-1},v].
 \label{ext:ape:regular}
\end{equation}

Suppose first that some $u_i\ne0$.  Let $r$ be the first nonzero
moment $\sum_i u_i(N_i/p_i)^r$, as in that theorem.  On
$w=1+c v^{-1-1/r}$ one has uniformly on compact $c$-sets
\[
 \mathcal E\longrightarrow E,\qquad
 \mathcal U\longrightarrow\eta,\qquad
 \mathcal L\longrightarrow L+A c^r,
 \qquad A\ne0.
\]
If $Q$ has positive degree in $V$, its nonzero resultant with $P$
specializes nontrivially at $(E,\eta)$.  A root $\alpha$ of the
specialized $Q$ satisfies $P(E,\eta,\alpha)\ne0$ and $\alpha\ne L$.
A nonzero solution of $L+Ac^r=\alpha$, followed by Rouch\'e's theorem,
then gives a genuine pole in \eqref{ext:ape:regular}.  All original Bessel
denominators tend to their nonzero $p_i$.  Hence $Q$ is independent
of $V$.

For the pure derivative case, choose $i$ with $v_i\ne0$.  The finite
$C_i$ chart of Theorem~\ref{ext:mae:main} gives
\[
 \mathcal E\longrightarrow E,\qquad
 \mathcal U_l\longrightarrow
       \eta_l^*:=\eta_l-\frac{M_lp_i}{N_i},\qquad
 \mathcal L\longrightarrow L+\frac{v_i p_i h_i}{c},
 \quad c\ne0.
\]
Consider the conjugation-stable field
\[
 F_0=K_{\exp}\bigl(C(t_l),C(\overline{t_l}):1\le l\le m\bigr),
\]
with repeated sign classes identified.  The real numbers
$\eta_1,\ldots,\eta_n$ are independent over $F_0$.  Apply
Lemma~\ref{ext:ape:imaginary-translation} with $T=\pi$ and
$c_l=\mp M_l/N_i$, according as $p_i=\pm\pi i$.  It follows that
the $\eta_l^*$ are independent over $F_0$.  In particular the tuple
$(E,\eta^*,L)$ is algebraically independent over $k$.  The same
resultant argument applies at $(E,\eta^*)$.  For a root $\alpha$
one has $\alpha\ne L$, so the nonzero choice
$c=v_i p_i h_i/(\alpha-L)$ and Rouch\'e's theorem again contradict
\eqref{ext:ape:regular}.  Thus in every case $Q$ has degree zero in $V$.

It remains to eliminate positive degree of $P$ in $V$.  Let $P_d(Z,U)$
be its nonzero leading coefficient, and choose any active block $i$.
On its pole divisor, parameterized by
\[
             v=-\Delta_i-\frac{p_i}{w^{N_i}-1},\qquad w\to1,
\]
the other Bessel denominators tend to
$p_l-N_lp_i/N_i\ne0$, while
\[
                   \mathcal E\longrightarrow E,
       \qquad \mathcal U_l\longrightarrow\eta_l^*.
\]
Apply Lemma~\ref{ext:ape:imaginary-translation} over $K_{\exp}$.
It makes $\eta_1^*,\ldots,\eta_n^*$ independent over that field,
and hence
$P_d(E,\eta^*)Q(E,\eta^*)\ne0$.
At a point of this divisor sufficiently close to $w=1$, the two
factors remain nonzero.  The affine coordinate $\mathcal L$ has a
genuine simple pole there: its singular numerator is
\[
 -\bigl(u_i b_i(b_i-p_i)+v_i h_i\bigr)(w^{N_i}-1),
 \qquad b_i=-p_i/(w^{N_i}-1),
\]
which cannot vanish near $w=1$ for an active index.  Thus a positive
$d$ produces a pole of order $d$ in \eqref{ext:ape:regular}, impossible.
Both $P$ and $Q$ are therefore independent of $V$, proving that the
original element belongs to $\mathcal F_+$.
\end{proof}

\begin{remark}[What this relative result leaves open]
The theorem allows arbitrary coefficients from all algebraic exponential
values and all positive-argument first exponential-integral values.
It does not assert that
$\overline{\mathbf E}\cap\mathcal F_+$ equals a specified polynomial
ring, and it does not classify the full field of three or more
independent harmonic Bessel coordinates at dependent scales.  The
restriction to positive Ein arguments supplies the reality used in
Lemma~\ref{ext:ape:imaginary-translation}; no analogous assertion for
arbitrary complex Ein arguments is made here.
\end{remark}


% Research bridge note only. No attachment or authoritative manuscript edited.
\section{Derivative coordinates: a valid transfer and a limitation}

Put $k=\overline{\mathbb Q}$ and retain the order-zero coordinates
$\beta(t)=Y(t)/I_0(t)$ and $C(t)=I_1(t)/I_0(t)$.

\begin{proposition}[All integer orders and derivative orders]
For every integer $m\ge0$, every $n\ge0$, and every positive algebraic
$t$, put
\[
 C_{m,n}(t)=\frac{I_m^{(n+1)}(t)}{I_m^{(n)}(t)}.
\]
The denominator is positive. There is a nondegenerate fractional-linear
transformation $M_{m,n,t}\in\operatorname{PGL}_2(k)$ such that
\[
 C_{m,n}(t)=M_{m,n,t}(C(t)).
\]
In particular, $k(C_{m,n}(t))=k(C(t))$.
\end{proposition}

\begin{proof}
The modified Bessel recurrences
\[
 I_{m+1}=I_{m-1}-\frac{2m}{t}I_m,\qquad
 K_{m+1}=K_{m-1}+\frac{2m}{t}K_m
\]
show that $I_m$ and $\widetilde K_m=(-1)^mK_m$ are obtained from
the respective pairs $(I_0,I_0')$ and $(K_0,K_0')$ by the same
rational coefficient row. The initial pairs agree with this convention
because $I_0'=I_1$ and $K_0'=-K_1$. Differentiation and the common
order-zero equation preserve this property. Consequently there is a
matrix $T_{m,n}(t)\in\operatorname{Mat}_2(\mathbb Q(t))$ with
\[
 \begin{pmatrix}I_m^{(n)}&\widetilde K_m^{(n)}\\
 I_m^{(n+1)}&\widetilde K_m^{(n+1)}\end{pmatrix}
 =T_{m,n}(t)
 \begin{pmatrix}I_0&K_0\\I_0'&K_0'\end{pmatrix}.
\]
These matrices have no poles away from $t=0$.

The power series of $I_m$ has positive coefficients and nonnegative
integer exponents, so $I_m^{(r)}(t)>0$ for every $r\ge0$ and $t>0$.
Moreover,
\[
 K_m(t)=\int_0^\infty e^{-t\cosh u}\cosh(mu)\,du
\]
implies $(-1)^rK_m^{(r)}(t)>0$. Thus the determinant on the left
of the matrix identity has sign $(-1)^{m+n+1}$ and is nonzero.
The order-zero Wronskian is also nonzero. Hence $T_{m,n}(t)$ is
invertible for every positive $t$. At algebraic $t$, its entries are
algebraic, and division of the second row value by the first gives
the asserted nondegenerate fractional-linear transformation.
\end{proof}

\begin{corollary}[Transfer of the attached winding theorem]
Use Theorem~\ref{wlbase:tb:main}, proved in \cite{WL}. For positive algebraic
$a$ and distinct positive algebraic $t_1,t_2$, and arbitrary
nonnegative integers $m_1,m_2,n_1,n_2$, one has
\[
 \overline{\mathbf E}\cap
 k(e^a,\Ein(a),\beta(t_1),C_{m_1,n_1}(t_1),
                    \beta(t_2),C_{m_2,n_2}(t_2))
 =k[e^a,e^{-a},\Ein(a)].
\]
Without the exponential-integral pair, the intersection equals $k$.
All additional ratios $C_{m,n}(t_i)$ may be adjoined simultaneously
without enlarging either field.
\end{corollary}

This is a coordinate-change consequence of the winding theorem. The
new resultant theorem is not required for this positive-scale,
integer-order statement. It supplies no additional algebraically
independent coordinates at a fixed scale.

\subsection{What the resultant theorem does supply}

Let $\nu\in\mathbb Q\setminus\frac12\mathbb Z$,
$f=x^{-\nu}J_\nu(x)$, $c=2\nu+1$, and let $A_n,B_n$ be the
coefficient polynomials from \cite{BR} at $d=0$:
\[
 f^{(n)}=x^{-n}(A_nf+xB_nf').
\]
For $n\ge1$, define
\[
 D_n(x)=A_n(x)B_{n+1}(x)-A_{n+1}(x)B_n(x).
\]

\begin{proposition}[Exact algebraicity locus for adjacent derivative ratios]
The polynomial $D_n\in\mathbb Q[x]$ is even and monic of degree $2n$.
For every nonzero algebraic $t$, the ratio
$L_n(t)=f^{(n+1)}(t)/f^{(n)}(t)$ is defined, and
\[
 L_n(t)\in k\quad\Longleftrightarrow\quad D_n(t)=0.
\]
If $D_n(t)\ne0$, then
$k(L_n(t))=k(f'(t)/f(t))$.
\end{proposition}

\begin{proof}
The resultant theorem gives $f^{(n)}(t)\ne0$ for every $n$.
Writing $L=f'/f$, the ratio is
\[
 L_n(t)=\frac{A_{n+1}(t)+tB_{n+1}(t)L(t)}
 {t(A_n(t)+tB_n(t)L(t))}.
\]
Siegel's theorem makes $L(t)$ transcendental. The displayed
fractional-linear transformation is nondegenerate exactly when
$D_n(t)\ne0$, in which case it is invertible over $k$ and its value
is transcendental. If $D_n(t)=0$, its coefficient rows are
proportional over $k$; the denominator row is nonzero, since
$f^{(n)}(t)\ne0$, and its value is algebraic.

For even $n=2r$, the $x$-leading coefficients of $A_n$ and $B_{n+1}$
are both $(-1)^r$, while the other determinant product has smaller
$x$-degree. For odd $n=2r+1$, the leading coefficients of $B_n$ and
$A_{n+1}$ are $(-1)^r$ and $(-1)^{r+1}$, respectively. In both cases
$D_n$ has leading term $x^{2n}$. Evenness follows from the recurrence.
\end{proof}

For example,
\[
 D_2(x)=x^2(x^2-c).
\]
At $\nu=1/3$, $c=5/3$ and $t=\sqrt{5/3}$, the nonvanishing theorem
applies, yet
\[
 \frac{f'''(t)}{f''(t)}=-\sqrt{5/3}.
\]
This also follows immediately by differentiating
$f''+(c/x)f'+f=0$. It is a concrete counterexample to the proposed
implication ``all derivative denominators are nonzero, therefore all
derivative ratios are transcendental or independent.''

\subsection{Why nonintegral orders do not inherit the winding line}

For $\nu\notin\mathbb Z$, a fundamental pair for the normalized
Bessel equation is
\[
 x^{-\nu}J_\nu(x),\qquad x^{-\nu}J_{-\nu}(x).
\]
The first member is single valued at zero and the second has exponent
$-2\nu$. A counterclockwise loop at zero therefore acts diagonally
with eigenvalues $1$ and $e^{-4\pi i\nu}$. For rational $\nu$, this
monodromy has finite order. The same statement holds for the modified
Bessel equation after rotating the argument.

In contrast, at order zero the second solution has logarithmic
monodromy $K_0\mapsto K_0-2\pi i I_0$. Its powers produce the
unbounded additive parameter used by the winding-line construction.
For rational nonintegral order, the finitely many values of the
diagonal monodromy cannot justify interpolation to a continuous
parameter, the large-winding parabolic limits, or the infinitely many
line slopes in the integral-intersection argument.

The resultant theorem addresses zeros of derivative coefficients; it
does not change this monodromy. A nonintegral-order analogue of the
exact integral-intersection theorem would require a new arithmetic or
geometric input. In particular, no statement about Euler's constant,
Gompertz's constant, or the full four-constant independence target
follows from combining these two mechanisms alone.


\section{What remains open}
The three-harmonic theorem gives the full rational field on three
harmonic ratios, while the affine-projection theorems permit any number
of blocks but retain just one affine coordinate. The relative theorem
also enlarges the coefficient field to $\mathcal F_+$, without
classifying $\overline{\mathbf E}\cap\mathcal F_+$ itself. These results do not give
\[
 \overline{\mathbf E}\cap
 k(\beta(t_1),C(t_1),\beta(t_2),C(t_2),\beta(t_3),C(t_3))=k
\]
for arbitrary dependent scales. Adjoining the exponential--Ein pair to
all three harmonic coordinates is also not justified by combining the
two theorems separately. For four harmonic coordinates the present
higher-chart and logarithmic-boundary arguments both have additional
cancellation problems.

The formal replacement argument of \cite{WL} concerns the unmarked
family of phase embeddings and the stated E-value inputs. It is not a
proof that every E-function method must fail to determine Euler's
constant. Extra properties distinguishing the canonical parameter
$s=\gamma$ are outside that formal replacement argument. None is
established in this supplement.

Finally, the factorial series $\sum_{n\ge0}(-1)^nn!$ diverges
archimedeanly: the Euler--Gompertz constant is its Borel/Hardy sum,
$\int_0^\infty e^{-u}/(1+u)\,du$. The existence of its separate p-adic
sums does not identify those values with the complex sum. This
clarification does not affect the phase-intersection proofs.

% Independent additions to the supplied Extensions manuscript.
% No EB or BR input is used in Theorem det:closed or Proposition det:modified.
% The final arithmetic interpretation of the exceptional locus retains the
% explicitly stated nonvanishing/Siegel inputs of the source manuscript.
\section{A closed derivative determinant and its exceptional set}
\label{det:section}

The coefficient determinant in the preceding derivative discussion admits
an explicit formula.  This calculation is independent of every phase
embedding and of the unpublished resultant input.  It is an addition to
the supplied manuscripts; no assertion of historical priority is intended.

Let $c$ be a complex parameter.  Define $A_0=1$, $B_0=0$ and
\begin{align}
 A_{n+1}&=xA_n'-nA_n-x^2B_n,\label{det:AB-a}\\
 B_{n+1}&=A_n+xB_n'+(1-c-n)B_n.\label{det:AB-b}
\end{align}
These recurrences give
$f^{(n)}=x^{-n}(A_nf+xB_nf')$ for every solution of
$f''+(c/x)f'+f=0$ away from zero.  Write
$D_n=A_nB_{n+1}-A_{n+1}B_n$, and let $(a)_j$ denote the rising factorial,
with $(a)_0=1$.

\begin{theorem}[Closed determinant formula]\label{det:closed}
For every integer $n\ge0$,
\begin{equation}\label{det:formula}
 D_n(x)=\sum_{j=0}^{\lfloor n/2\rfloor}
 \frac{(-1)^j n!(c/2)_j}{(n-2j)!\,j!}\,x^{2n-2j}.
\end{equation}
In particular, this identity holds in $\mathbb Q[c,x]$.
\end{theorem}

\begin{proof}
Fix $c$ and choose independent local solutions $f,g$ on a disc avoiding
zero.  Put $v_n=(f^{(n)},g^{(n)})$ and
$W_n=\det(v_n,v_{n+1})$.  Differentiating
$xf''+cf'+xf=0$ exactly $n$ times gives, for $n\ge1$,
\[
 xv_{n+2}+(c+n)v_{n+1}+xv_n+nv_{n-1}=0.
\]
Taking the determinant with $v_{n+1}$ shows that
\[
 W_{n+1}=W_n+\frac n x\det(v_{n-1},v_{n+1})
         =W_n+\frac n x W_{n-1}'.
\]
Abel's identity gives $W_0'/W_0=-c/x$, and the original equation gives
$W_1=W_0$.  Therefore $r_n=W_n/W_0$ satisfies
\[
 r_0=r_1=1,\qquad
 r_{n+1}=r_n+\frac n x r_{n-1}'-\frac{nc}{x^2}r_{n-1}.
\]
An induction now gives
\[
 r_n(x)=\sum_{j=0}^{\lfloor n/2\rfloor}
 \frac{(-1)^j n!(c/2)_j}{(n-2j)!\,j!}\,x^{-2j}.
\]
For completeness, if
$a_{n,j}=n!(c/2)_j/((n-2j)!j!)$ in its indicated range and is zero
otherwise, the required coefficient identity is
\[
 a_{n+1,j}=a_{n,j}+n(c+2j-2)a_{n-1,j-1}.
\]
It follows directly from $(c/2)_j=(c/2+j-1)(c/2)_{j-1}$.
Finally the two derivative rows in the coefficient representation have
determinant $x^{-2n}D_n$, so $W_n=x^{-2n}D_nW_0$.
This proves the formula for each $c$ and hence the polynomial identity.
\end{proof}

\begin{corollary}[The exact nonzero degree]\label{det:degree}
Suppose that $c\notin\{0,-2,-4,\ldots\}$.  The multiplicity of zero
in $D_n$ is exactly $2\lceil n/2\rceil$.  Thus $D_n$ has exactly
$2\lfloor n/2\rfloor$ nonzero roots counted with multiplicity.
If $c>0$, none of those roots lies on the imaginary axis.
\end{corollary}

\begin{proof}
The coefficient of the lowest displayed power in
\eqref{det:formula} is nonzero under the hypothesis on $c$, which gives
the multiplicity and the remaining degree.  For real $y\ne0$ and $c>0$,
\[
 \frac{D_n(iy)}{(iy)^{2n}}
 =\sum_{j=0}^{\lfloor n/2\rfloor}
 \frac{n!(c/2)_j}{(n-2j)!\,j!}\,y^{-2j}>0.
\]
\end{proof}

For example,
\begin{align*}
 D_2&=x^2(x^2-c),\\
 D_3&=x^4(x^2-3c),\\
 D_4&=x^4\bigl(x^4-6cx^2+3c(c+2)\bigr),\\
 D_5&=x^6\bigl(x^4-10cx^2+15c(c+2)\bigr).
\end{align*}
Under the nonvanishing and Siegel inputs used for the adjacent-ratio
algebraicity criterion, these are precisely its exceptional loci.
For $\nu\in\mathbb Q\setminus\tfrac12\mathbb Z$ and $c=2\nu+1$,
the exceptional algebraic arguments have total multiplicity
$2\lfloor n/2\rfloor$, rather than the degree bound $2n$ obtained
from monicity alone.  This is a multiplicity count, not a claim that
all roots are distinct or real.

The polynomials $P_n(x;c)=x^{-n}D_n(x)$ also satisfy the formal generating
identity
\[
 \sum_{n=0}^\infty P_n(x;c)\frac{z^n}{n!}
 =e^{xz}(1+z^2)^{-c/2}.
\]
In particular $P_n'=nP_{n-1}$.  This is an Appell polynomial identity;
it does not imply that the zeros are all real.  For instance, the
quadratic in $x^2$ in $D_4$ has discriminant $24c(c-1)$, so for
$0<c<1$ its nonzero roots are not real or purely imaginary.

\subsection{No degeneration at positive modified-Bessel arguments}

\begin{proposition}[All adjacent normalized modified derivatives]
\label{det:modified}
Let $\nu\in\mathbb Q$, $\nu>-\tfrac12$, and put
$F_\nu(t)=t^{-\nu}I_\nu(t)$ for $t>0$.  For every integer $n\ge0$
and every positive algebraic $t$, the ratio
\[
 L_{\nu,n}(t)=\frac{F_\nu^{(n+1)}(t)}{F_\nu^{(n)}(t)}
\]
is defined and is a nondegenerate fractional-linear transformation over
$\overline{\mathbb Q}$ of $F_\nu'(t)/F_\nu(t)$.  Consequently
\[
 \overline{\mathbb Q}(L_{\nu,n}(t))
 =\overline{\mathbb Q}\!\left(\frac{F_\nu'(t)}{F_\nu(t)}\right)
 =\overline{\mathbb Q}\!\left(\frac{I_{\nu+1}(t)}{I_\nu(t)}\right).
\]
This proposition uses neither the phase embeddings nor the derivative
resultant manuscript.
\end{proposition}

\begin{proof}
The function $F_\nu$ satisfies
$F_\nu''+(c/t)F_\nu'-F_\nu=0$, where $c=2\nu+1>0$.
Its series is a positive constant times
\[
 \sum_{m=0}^\infty\frac{t^{2m}}{4^m m!(\nu+1)_m}.
\]
Every derivative is strictly positive on $(0,\infty)$, since infinitely
many positive terms survive differentiation.  Repeated use of the
equation expresses $(F_\nu^{(n)},F_\nu^{(n+1)})$ by a rational
coefficient matrix applied to $(F_\nu,F_\nu')$.

Rotating the ordinary equation by $x=it$ in
Theorem~\ref{det:closed}, or repeating its Wronskian recurrence with the
sign of the last term reversed, shows that the determinant of this
matrix is
\[
 (-1)^n\sum_{j=0}^{\lfloor n/2\rfloor}
 \frac{n!(c/2)_j}{(n-2j)!\,j!}\,t^{-2j}.
\]
It is nonzero for positive $t$.  At algebraic $t$ the matrix has
algebraic entries, so its two row values give an invertible
fractional-linear transformation.  Finally the standard derivative
recurrence gives $F_\nu'/F_\nu=I_{\nu+1}/I_\nu$.
\end{proof}

For $\nu\in\mathbb Q\setminus\tfrac12\mathbb Z$ with $\nu>-\tfrac12$, the published
Siegel transcendence theorem for the base normalized Bessel logarithmic
derivative (as stated and used in \cite{LM95}) additionally makes every
$L_{\nu,n}(t)$ transcendental.
This arithmetic corollary requires only that classical theorem and
the preceding proof; it does not require the derivative-nonvanishing
resultant input.  The field equality is a statement at one fixed order
and scale.  It does not identify fields at different nonintegral orders
or supply the additive monodromy needed by the winding method.

% Add to bibliography:
% \bibitem{LM95} L.~Lorch and M.~E.~Muldoon,
% \emph{Transcendentality of zeros of higher derivatives of functions
% involving Bessel functions}, Internat. J. Math. Math. Sci. \textbf{18}
% (1995), no.~3, 551--560.
% \href{https://doi.org/10.1155/S0161171295000706}
% {doi:10.1155/S0161171295000706}.

\clearpage
\part{Pisot floors, prime clocks, and local obstructions}
\section{The problem, the source boundary, and the new results}
Let
\[
 A=\min\{x>1:\lfloor x^{3^k}\rfloor\text{ is prime for every }k\ge1\}
\]
be Mills' constant. Saito proved its irrationality in
\cite{SaitoIrrational}. His later reduction \cite[Theorem~1.7]{SaitoVariants} states that if $A$ is
algebraic, some $\beta=A^{3^m}$ is a totally real cubic Pisot number
whose two small conjugates satisfy
\begin{equation}\label{pis:eq:saitobounds}
 |\beta_3|<-\beta_2\le
 \min\{|\beta_3|^{17/23},\beta^{-17/40}\}.
\end{equation}
The trace is then the floor at all sufficiently large cubic-power
levels. Excluding every such algebraic candidate would prove
unconditional transcendence of $A$. Saito's Problem~1.1 in
\cite{SaitoIntervals} asks the stronger question: does every cubic
Pisot number have infinitely many composite cubic-power floors?

The word ``question'' matters here. The source papers contain
explicit problems and alternatives, and an answer to a modified
sequence does not automatically settle the original least constant.
All source theorem numbers below refer to the specified arXiv
versions, since journal numbering may differ.

\subsection{What is already resolved and what remains}
\begin{center}\small
\begin{tabular}{p{.25\textwidth}p{.29\textwidth}p{.35\textwidth}}
\toprule
Source and item & Mathematical target & Status used in this paper\\
\midrule
\cite{SaitoIrrational}, Q1.6 & Irrationality of the least constant with exponent $2^k$ & Not proved here.\\
\cite{SaitoIrrational}, Q1.7 & The cubic Pisot power alternative for Mills' constant & Not excluded in full.\\
\cite{SaitoIrrational}, Q5.3 & RH implies Mills transcendence & Already answered by Saito; his later Theorem~1.8 uses the Density Hypothesis.\\
\cite{SaitoIntervals}, P1.1 & All cubic Pisot floors at $3^n$ & Restricted closures and quantitative exceptions below; full problem remains.\\
\cite{SaitoIntervals}, P1.7--1.8 & Nonreversible clocks and shifted dyadic Fibonacci values & Already treated in the supplied older manuscripts [008] and [007].\\
\cite{SaitoDimension}, Q1.4--1.5 & Unbounded multiplier and Wright-set dimensions & Already treated in supplied older [031], Theorems~14.3 and~15.1.\\
\bottomrule
\end{tabular}
\end{center}
The last two rows report the source manuscript's descriptions of earlier
research notes. Those notes are not supplied with this edition and their
claimed resolutions are not independently certified here. They are not counted again as new resolutions.
The RH implication is a published external result of Saito, not a
consequence newly obtained from the present work.

\subsection{Contribution relative to the supplied manuscripts}
The source Pisot manuscript reports a prior inventory of 244 TeX
manuscripts, including 38 in the Mills--Pisot topic. That archive is not
part of the three-manuscript source set for this edition. References to
it below record the source manuscript's provenance, not a new audit of
those 244 files.
The core additions and their limitations are as follows.
\begin{enumerate}
\item Theorem~\ref{pis:thm:windows} synchronizes finitely many shifts on
one progression. The old single-shift positive-density statements
alone do not supply this synchronization.
\item Theorem~\ref{pis:thm:shiftedsieve} extends the unshifted cubic
coefficient sieve of round S37 to every shift and every odd degree,
uniformly in the shift, with an improved explicit range of constants.
The older [031] already contains qualitative density-one results;
the quantitative rate and the stated extensions are the additions.
\item Theorem~\ref{pis:thm:pell-new} adds three exact witnesses to the
Pell cover of rounds S35--S36. It removes two further valuation
layers, but leaves an infinite suborbit to be analyzed.
\item Theorem~\ref{pis:fac:valuationpatterns} refines the finite patterns
of S28 by specifying exact selected prime valuations and their
periodic recurrence. Prime-designated positions still require
Schinzel's hypothesis H. The all-index irreducibility of S27 is
reproduced as prior input, not counted as new.
\item Theorem~\ref{pis:transfer:thm:ray} quantifies the persistence of
Saito-balanced units in a fixed congruence coset, explaining why
finite local exclusions and the real size conditions can coexist.
\item Theorems~\ref{pis:res:thm:ambient} and~\ref{pis:res:thm:galois} turn
an inherited cubic local obstruction into exclusions for all rational
powers of Mills' constant in larger number fields. Theorem~
\ref{pis:rescyc:thm:norm-one-mod-three} records the existing cyclic
norm restriction in a positive-density formulation for both floor branches.
\item Theorem~\ref{pis:res:thm:equivalence} removes the seed assumption
from finite-prime survival by allowing a fixed nonzero integral norm.
It applies to every remaining totally real cubic field and every
prescribed finite prime set. Corollary~\ref{pis:res:cor:normalized}
places such survivors in any real interval after normalization.
\end{enumerate}
The manuscript contains mathematical proofs and exact finite
certificates; external priority remains unestablished. In particular,
none of the remaining numbered global questions in the first four
rows is claimed as a new complete resolution.

\subsection{Conventions}
A composite integer is positive and larger than one, with a proper
prime divisor. Finite initial exceptions are harmless only after
growth has been proved. Lower density always refers to the index
$n$, not to the enormous integer value of the trace. A Pisot number
is a real algebraic integer greater than one with every other
conjugate strictly inside the unit circle. Traces count all conjugates
with multiplicity when polynomial specializations are used.


\section{Simultaneous composite windows from a three-adic limit}\label{pis:sec:windows}
For a monic integral cubic $X^3-uX^2+vX-N$, let $U_n$ be its
$3^n$-th root power sum, let $V_n$ be the corresponding second
symmetric function, and let $E_n=N^{3^n}$. Newton's identities give
\begin{equation}\label{pis:eq:cubicmap}
 (U_{n+1},V_{n+1},E_{n+1})=
 (U_n^3-3U_nV_n+3E_n,\,
 V_n^3-3E_nU_nV_n+3E_n^2,\,E_n^3).
\end{equation}
The norm coordinate must be advanced, except when $N=0,1,-1$.

\begin{lemma}[Three-adic contraction]\label{pis:lem:contraction}
The sequence $U_n$ has a limit $\lambda\in\Z_3$, with
$\lambda-U_n\in3^{n+1}\Z_3$. This limit depends only on
$(u,v,N)\pmod3$.
\end{lemma}
\begin{proof}
Let $F$ denote the polynomial map in \eqref{pis:eq:cubicmap}. It is the
identity modulo three. If $\mathbf x\equiv\mathbf y\pmod{3^r}$,
$r\ge1$, expansion gives
$F(\mathbf x)\equiv F(\mathbf y)\pmod{3^{r+1}}$; every first
derivative has coefficients divisible by three, and higher terms
have the requisite divisibility. Consecutive iterates therefore
differ by a vector divisible by $3^{n+1}$ at step $n$. The same
argument for two initial points congruent modulo three proves
equality of their limits.
\end{proof}

\begin{theorem}[A common progression for finitely many shifts]
\label{pis:thm:windows}
Suppose $U_n\to+\infty$ in the real topology. Let
$\mathcal H\subset\Z$ be finite and assume
\[
 \lambda+h\notin\{1,-1\}\qquad(h\in\mathcal H).
\]
There are integers $n_0,D\ge1$ such that all
$U_{n_0+jD}+h$, $h\in\mathcal H$, are composite for every
sufficiently large $j$.
\end{theorem}
\begin{proof}
Choose a common large $n_0$ so that $Q_h=U_{n_0}+h>1$ for all
shifts. When $3\nmid\lambda+h$, also arrange
\[
 v_3(Q_h-1)<n_0,\qquad v_3(Q_h+1)<n_0.
\]
The two limiting valuations are finite by hypothesis, so
Lemma~\ref{pis:lem:contraction} permits this choice.
If $3\mid\lambda+h$, the prime three divides $U_n+h$ at every
level and supplies a recurring divisor.

For any other shift factor $Q_h$ into primes. If every prime factor
$q$ satisfied $v_3(q^2-1)\ge n_0$, then each such $q$ would be
$1$ or $-1$ modulo $3^{n_0}$. Their product, including
multiplicities, would make $Q_h\equiv\pm1\pmod{3^{n_0}}$,
a contradiction. Hence there is a prime $q_h\mid Q_h$ with
\[
 v_3(q_h^2-1)\le n_0-1.
\]
This prime is not three. Every nonzero root of the cubic modulo
$q_h$ has order dividing
$M_h=\operatorname{lcm}(q_h-1,q_h^2-1,q_h^3-1)$.
Elementary lifting of valuations gives
\[
 v_3(M_h)\le v_3(q_h^2-1)+1\le n_0.
\]
After powering to $3^{n_0}$, all remaining root orders are prime
to three. A period $D_h$ for cubing on these groups therefore gives
$q_h\mid U_{n_0+jD_h}+h$ for every $j\ge0$. Zero roots and repeated
roots cause no difficulty; the power sum counts multiplicities.

Take $D$ to be the least common multiple of the finitely many
$D_h$, with $D=1$ when there are none. All chosen prime divisors recur on the same progression.
Real growth makes them proper divisors for all sufficiently large
$j$. Unlike a prime-seed argument, this construction works even
when every initial value selected is already composite.
\end{proof}

\begin{corollary}[Windows around cubic Pisot floors]\label{pis:cor:pisotwindows}
Let $\beta$ be a cubic Pisot number and let $\lambda$ be its
three-adic trace limit. If $\lambda$ is not an ordinary integer,
then for every $H\ge0$ there is an arithmetic progression of indices
on which, eventually, all
\[
 \lfloor\beta^{3^n}\rfloor+h,\qquad |h|\le H,
\]
are composite simultaneously. Nonreal conjugates are permitted.
\end{corollary}
\begin{proof}
Apply Theorem~\ref{pis:thm:windows} to all integer shifts in
$[-H-1,H]$. None can give $\lambda+h=\pm1$. The sum of the two
small conjugate powers tends to zero, so the floor is eventually
either $U_n$ or $U_n-1$. The stated floor window lies inside the
chosen trace window, regardless of changes in the sign of the
small-conjugate sum.
\end{proof}

Separate positive-density statements for different shifts cannot
be intersected without an argument. The common initial index and
the prime-divisor construction above provide that missing step.
The single-shift order-budget mechanism was already present in the
older manuscript \cite{Keio009}; the simultaneous progression is
the strengthening here.

\begin{proposition}[Exact residue-class scope]\label{pis:prop:fourteen}
The limit $\lambda$ is an ordinary integer precisely when the
residual cubic over $\F_3$ either splits completely or is
$(X-r)(X^2+1)$ with $r\in\F_3$. Thus
Corollary~\ref{pis:cor:pisotwindows} applies to fourteen of the
twenty-seven coefficient triples modulo three.
\end{proposition}
\begin{proof}
For each nonzero residual root, take its Teichm\"uller lift, the
unique prime-to-three root of unity with that residue. In a
three-adic splitting field an integral root of positive valuation
vanishes under repeated cubing; a unit root's principal-unit part
tends to one. Frobenius permutes the multiset of residual roots,
so $\lambda$ is the sum of the Teichm\"uller lifts, with
multiplicity.

A linear root contributes $0,1$, or $-1$. The irreducible quadratic
$X^2+1$ contributes two fourth roots of unity with sum zero. The
other two irreducible quadratics are $X^2+X-1$ and $X^2-X-1$.
Their lifted pairs have sum with square $-2$, so adding one linear
contribution cannot give an integer.

For an irreducible cubic, the nonzero roots have order thirteen or
twenty-six. The lifted trace is
$\pm(\zeta^a+\zeta^{3a}+\zeta^{9a})$, where $\zeta$ is a
primitive thirteenth root and $a\ne0\pmod{13}$. This is not
rational: otherwise a polynomial of degree at most twelve with
at most four nonzero coefficients would vanish at $\zeta$, in
contradiction to its minimal polynomial
$1+X+\cdots+X^{12}$. There are ten completely split monic cubics
(multisets of three roots from a three-element set), three of the
second integer-limit type, six with the other quadratic factors,
and eight irreducible cubics. These exhaust the twenty-seven cases.
\end{proof}

\begin{center}
\begin{tabular}{c l}
\toprule
Integer limit $\lambda$ & $(u,v,N)\pmod3$\\
\midrule
$-3$ & $(0,0,2)$\\
$-2$ & $(1,1,0)$\\
$-1$ & $(2,0,0),\ (2,1,2),\ (2,2,1)$\\
$0$ & $(0,0,0),\ (0,1,0),\ (0,2,0)$\\
$1$ & $(1,0,0),\ (1,1,1),\ (1,2,2)$\\
$2$ & $(2,1,0)$\\
$3$ & $(0,0,1)$\\
\bottomrule
\end{tabular}
\end{center}
The complement of this table is the fourteen-class scope. This is
a count of residue classes, not a density assertion about Pisot
numbers under an unspecified ordering. For example, the plastic
number, the real root of $X^3-X-1$, has irreducible reduction modulo
three and satisfies the simultaneous-window conclusion.

For the unshifted trace problem only the six classes with
$\lambda=\pm1$ can escape Theorem~\ref{pis:thm:windows}. These include
the surviving Mills alternatives. For the predecessor trace
$U_n-1$, only the four classes with $\lambda=0$ or $2$ can escape;
all four have $N\equiv0\pmod3$. In particular, whenever
$3\nmid N$ and the floor is eventually $U_n-1$, its composite
indices have positive lower density. This last single-shift
consequence was already available in the earlier local obstruction.

% This section follows windows_section.tex and uses its labels and notation.
\section{Field criteria and a uniform density bound}
\label{pis:local:sec:fields}
The preceding simultaneous-return theorem has a useful field-theoretic
consequence. We first record an elementary bound that also handles
changing signs of nonreal conjugates.

\begin{proposition}[A parity and trace criterion]
\label{pis:local:prop:half}
Let $\beta$ be a cubic Pisot number of odd norm. If
$\operatorname{Tr}(\beta)\equiv0\pmod3$, then
\[
 \liminf_{M\to\infty}\frac{1}{M}
 \#\{0\le n<M:\lfloor\beta^{3^n}\rfloor
                  \text{ is composite}\}\ \ge\ \frac12.
\]
The two small conjugates may be nonreal.
\end{proposition}
\begin{proof}
Let $U_n,V_n,E_n$ be the coefficient sequences in
\eqref{pis:eq:cubicmap}. Since the norm is odd, $E_n=1$ modulo two.
The map on $(U_n,V_n)\pmod2$ is
\[
 (a,b)\longmapsto(a+ab+1,\ b+ab+1).
\]
Its four states satisfy
\[
 00\longmapsto11,\qquad11\longmapsto11,
 \qquad01\longleftrightarrow10.
\]
Consequently $U_n$ is odd either at every sufficiently large index
or at exactly one parity of the indices. At least half of all
indices therefore have $U_n$ odd. At the same time,
\eqref{pis:eq:cubicmap} modulo three gives $U_n\equiv U_0\equiv0\pmod3$
for every $n$.

The sum of the two small conjugate powers is real and tends to zero.
Thus, for every sufficiently large $n$,
$\lfloor\beta^{3^n}\rfloor$ belongs to $\{U_n,U_n-1\}$.
When $U_n$ is odd, the first possibility is divisible by three and
the second is divisible by two. Growth makes these divisors proper.
This proves the density bound without any assumption on the sign
pattern of the small-conjugate sum.
\end{proof}

\begin{theorem}[Cubic fields without a three-adic embedding]
\label{pis:local:thm:noembedding}
Let $K$ be a cubic number field admitting no field embedding into
$\mathbb Q_3$. Let $\beta\in\mathcal O_K$ be a cubic Pisot number
with $3\nmid N_{K/\mathbb Q}(\beta)$. Then the composite indices of
$\lfloor\beta^{3^n}\rfloor$ contain, up to finitely many terms,
an arithmetic progression. In particular they have positive lower
density.

More precisely, write $f$ for the minimal polynomial of $\beta$.
If $f\bmod3$ is irreducible, every finite window about these floors
is simultaneously composite on a progression, as in
Corollary~\ref{pis:cor:pisotwindows}. Otherwise $f\bmod3$ is a cube
of a nonzero linear factor. In this second case, if the norm is
also odd, the lower density is at least $1/2$.
\end{theorem}
\begin{proof}
A simple root of $f\bmod3$ would lift by Hensel's lemma to a root
in $\mathbb Q_3$. Since $K=\mathbb Q(\beta)$, that root would define
an embedding $K\hookrightarrow\mathbb Q_3$, contrary to hypothesis.
A monic cubic over $\mathbb F_3$ without a simple linear factor
is either irreducible or $(X-a)^3$. The norm assumption makes
$a\ne0$ in the latter case.

In the irreducible case Proposition~\ref{pis:prop:fourteen} gives a
noninteger three-adic trace limit, and
Corollary~\ref{pis:cor:pisotwindows} proves the asserted window statement.
In the remaining case the three Teichm\"uller lifts of the residual
roots are all $1$ or all $-1$. Thus the trace limit is $\lambda=3$
or $\lambda=-3$, respectively. Neither $\lambda$ nor $\lambda-1$
equals $1$ or $-1$. Apply Theorem~\ref{pis:thm:windows} to the two
shifts $\{0,-1\}$. Both $U_n$ and $U_n-1$ are composite on one
eventual progression; the Pisot estimate again places the floor
among these two integers. Finally $(X-a)^3$ has zero quadratic
coefficient modulo three, so the odd-norm case satisfies
Proposition~\ref{pis:local:prop:half}.
\end{proof}

\begin{corollary}[Totally ramified cubic fields]
\label{pis:local:cor:ramified}
If three is totally ramified in a cubic number field $K$, then
every cubic Pisot unit of $K$ has composite floors at a set of
indices of lower density at least $1/2$.
\end{corollary}
\begin{proof}
The algebra $K\otimes_{\mathbb Q}\mathbb Q_3$ is a field of
degree three, so $K$ has no embedding into $\mathbb Q_3$.
For any $\beta\in\mathcal O_K$, its three local conjugates have
the same residue in the residue field $\mathbb F_3$. Its minimal
polynomial therefore reduces to $(X-a)^3$ modulo three. A unit
has nonzero residue and norm $\pm1$. The last assertion of
Theorem~\ref{pis:local:thm:noembedding} applies.
\end{proof}

For example, $K=\mathbb Q(r)$ with $r^3-3r-1=0$ is a totally
real cyclic cubic field: the polynomial is irreducible, has
discriminant $81$, and its translate
$X^3+3X^2-3$ is Eisenstein at three. Hence the corollary covers
every Pisot unit in this field. It also covers every Pisot unit in
$\mathbb Q(\sqrt[3]{2})$, since
$(X-1)^3-2=X^3-3X^2+3X-3$ is Eisenstein at three.
An explicit example in the latter field is
$\beta=1+\sqrt[3]{2}+\sqrt[3]{4}$, whose minimal polynomial is
$X^3-3X^2-3X-1$. It has norm one and a unique real root greater
than one; its conjugate pair has modulus $\beta^{-1/2}<1$.
Thus the density bound applies even though the floor correction
need not have an eventually constant sign.

\paragraph{Completion of the norm-positive case in the field of
discriminant $1524$.}
Let $r$ be the largest root of $X^3-7X^2-X+1$ and let
$K=\mathbb Q(r)$. Showa round 20 proves that every norm-one unit
$\beta\in K$ with $\operatorname{Tr}(\beta^{3^n})\to+\infty$
has a positive lower density of composite trace terms. This is a
prior result of the supplied archive, not a new trace theorem here.
It also yields the following precise floor statement.

\begin{corollary}[Conditional on the cited round-20 trace theorem]
\label{pis:local:cor:1524}
Every cubic Pisot unit of norm one in the field of discriminant
$1524$ has a positive lower density of composite floors
$\lfloor\beta^{3^n}\rfloor$.
\end{corollary}
\begin{proof}
The field is totally real. Since $\beta>1$ and its norm is one,
its two small conjugates have the same sign. If both are negative,
the floor equals the trace eventually, and the cited round-20
theorem applies. If both are positive, the floor equals $U_n-1$
eventually. The four-state parity calculation in the proof of
Proposition~\ref{pis:local:prop:half} shows that $U_n$ is odd at a
set of indices of density at least $1/2$. At those sufficiently
large indices $U_n-1$ is a positive even composite integer.
\end{proof}

\paragraph{Scope and comparison with the supplied manuscripts.}
The source manuscript reports that \cite{Keio009} proves infinitely many
composite floors for every nonreal cubic Pisot number, and gives
positive-density results for fixed real shifts. Its projective
return theorem already includes the cubic-residue norm refinement;
that refinement is not new here. The absence of a three-adic
embedding was already a necessary obstruction for a hypothetical
prime trace tail in \cite{Keio031}. The new use here is the
simultaneous treatment of the two possible floor corrections,
together with the explicit density-$1/2$ criterion. In particular
the statements above do not depend on an unjustified intersection
of separate positive-density sets.

The discriminant-$1524$ corollary is a completion from traces to
all norm-one Pisot floors using a prior field-specific theorem.
It makes no assertion about norm-minus-one units in that field:
round 17 exhibits such units that survive several local tests.
Nor do these criteria exclude every totally real cubic Pisot
candidate in Saito's Problem~1.1 or establish Mills transcendence.

\section{Shifted traces: exact local counts and sparse exceptions}\label{pis:sec:density}
This section extends the unshifted cubic sieve of submitted round
S37. The extensions are uniformity in the integer shift, both floor
branches in totally real cubic fields, and arbitrary odd degree. We
include the proof, since changing a trace to a shifted trace is not
automatically legitimate in a fixed-norm argument.

For an odd integer $d\ge3$, write
\[
 f_{\mathbf a}(X)=X^d-a_1X^{d-1}+a_2X^{d-2}-\cdots-a_d,
 \qquad P_n(\mathbf a)=\sum_{j=1}^d\alpha_j^{3^n}.
\]
Here the roots are counted with multiplicity, and Newton's identities
define $P_n\in\Z[a_1,\ldots,a_d]$ even at reducible specializations.
Its total degree is $3^n$, and its unique term of that degree is
$a_1^{3^n}$. For a sequence of integers, a \emph{recurring-divisor
certificate} is a prime dividing all terms on an infinite arithmetic
progression of indices. If the terms tend to positive infinity, such
a certificate gives positive lower density of composite terms.

\begin{theorem}[Uniform shifted sieve]\label{pis:thm:shiftedsieve}
Put $c_0=1/(e\sqrt{2\log3})$. Fix an odd $d\ge3$ and $0<c<c_0$.
Let $\mathcal B\subset\Z^d$ be any box with side lengths
$X_1,\ldots,X_d\ge3$, and put $m=\min X_i$ and $V=\prod X_i$.
For every integer $h$, all but
\begin{equation}\label{pis:eq:exceptionbound}
 O_{d,c}\bigl(V\exp(-c\sqrt{\log m})\bigr)
\end{equation}
coefficient vectors in $\mathcal B$ have a recurring-divisor
certificate for $P_n(\mathbf a)+h$. The constant is independent of
$h$ and of the location of the box.

There is a constant $C_{d,c}>0$ such that the certificate can be
chosen with prime $p\le\exp(C_{d,c}\sqrt{\log m})$ and progression
period at most $\exp(dC_{d,c}\sqrt{\log m})$. Consequently, whenever
$P_n(\mathbf a)\to+\infty$, the lower density of composite shifted
terms is positive, and is at least the reciprocal of that period.
\end{theorem}

The theorem gives an upper bound on a set of possible exceptions; it
does not prove that this set is empty. The assertion is about the
full coefficient space. A thin fixed field, fixed norm, or selected
one-parameter family requires a separate argument.

\begin{lemma}[Exact shifted local count]\label{pis:lem:shiftedcount}
For every prime $p\equiv2\pmod3$, $n\ge1$, and $t\in\F_p$,
\[
 \#\{\mathbf a\in\F_p^d:P_n(\mathbf a)=t\}=p^{d-1}.
\]
Repeated roots and zero constant terms are allowed.
\end{lemma}
\begin{proof}
The coefficient vectors are Frobenius-stable orbits of ordered root
tuples under $S_d$. Twisted orbit counting gives the number in question
as $d!^{-1}\sum_{\sigma\in S_d}\#\operatorname{Fix}(F\sigma)$,
where the fixed tuples also satisfy the power-sum condition. This is
an identity of integer counts, including orbits with stabilizers.
All tuples under consideration lie in a common finite extension, so
ordinary finite-set counting applies.

A cycle of length $j$ in $\sigma$ corresponds to an element
$b\in\F_{p^j}$ and contributes
$\operatorname{Tr}_{\F_{p^j}/\F_p}(b^{3^n})$ to the power sum.
Every permutation of an odd number of letters has an odd cycle.
For such a cycle $j$, $3\nmid p^j-1$, so $b\mapsto b^{3^n}$ is a
bijection of $\F_{p^j}$. The field trace is surjective, with fibers
of size $p^{j-1}$. Fixing all other cycles leaves exactly that many
choices for this one, giving $p^{d-1}$ fixed tuples for each
permutation. Averaging proves the claim. Surjectivity of field trace
does not require $p\nmid j$.
\end{proof}

\begin{lemma}[Geometric separation]\label{pis:lem:shiftedgeometry}
Over an algebraically closed field of characteristic different from
three, $P_n-t$ is irreducible for every $n\ge1$ and every $t$.
For distinct indices $i,j$ and any $s,t\in\F_p$, $p\ne3$,
\[
 \#\{P_i=s,\ P_j=t\}\le 3^{i+j}p^{d-2}.
\]
\end{lemma}
\begin{proof}
Under the finite, surjective Vieta map the first polynomial pulls
back to $Z_1^r+\cdots+Z_d^r-t$, $r=3^n$. Regard this as a polynomial
in $Z_d$ over $k[Z_1,\ldots,Z_{d-1}]$. The nonconstant polynomial
$Z_1^r+\cdots+Z_{d-1}^r-t$ has a prime factor with multiplicity one:
it is squarefree, since any repeated factor would divide all its
partial derivatives, which are multiples of the relatively prime
monomials $Z_1^{r-1},\ldots,Z_{d-1}^{r-1}$. Eisenstein at such a
factor proves irreducibility of the pullback, and hence of $P_n-t$.

Different indices have different degrees, so the two hypersurfaces
have no common component. Their intersection has pure dimension
$d-2$ and degree at most $3^{i+j}$, by B\'ezout. An affine variety of
dimension $r$ and degree $D$ has at most $Dp^r$ rational points.
One proof is induction on dimension: on each irreducible component
choose a nonconstant coordinate, slice at its $p$ rational values,
and apply the degree bound for hyperplane sections. Constant
coordinates are discarded. This argument also handles components
defined only over an algebraic closure, since it bounds their
$\F_p$-points individually. These facts give the displayed bound.
\end{proof}

\begin{lemma}[A uniform return budget]\label{pis:lem:returnbudget}
Fix $B\ge1$ and put $b_d=\lfloor\log_3(d/2)\rfloor$.
Suppose $p\equiv2\pmod3$ and $v_3(p+1)\le B$. For every coefficient
vector modulo $p$, its power-sum sequence is periodic from index
$B+b_d$, with period less than $p^d$.
\end{lemma}
\begin{proof}
For an irreducible root factor of degree $j$, each nonzero root has
order dividing $p^j-1$. For odd $j$ this order is prime to three;
for even $j=2k$, the lifting-the-exponent identity gives
\[
 v_3(p^{2k}-1)=v_3(p+1)+v_3(k)\le B+b_d.
\]
After that many cubings all nonzero root powers have orders prime
to three. Cubing is then a permutation of their finite groups.
Zero roots stay zero and multiplicities do not change this argument.
The least common multiple $M$ of the remaining root orders is at
most the product of $p^j-1$ over distinct nonzero irreducible factors,
which is less than $p^d$. A common period is the order of $3$ modulo
$M$, with period one when $M=1$.
\end{proof}

\begin{proof}[Proof of Theorem~\ref{pis:thm:shiftedsieve}]
The eligible primes of Lemma~\ref{pis:lem:returnbudget} form fixed
reduced residue classes modulo $3^{B+1}$ of density
\[
 \delta=\delta_B=\frac{1-3^{-B}}2.
\]
Let $k_0=B+b_d$ and use the $L$ consecutive levels
$k_0,\ldots,k_0+L-1$. At an eligible prime let $\Omega_p$ be the
union of the loci $P_n+h=0$ at these levels, and put
$\omega(p)=|\Omega_p|$. Exact counts and the first two
inclusion--exclusion inequalities give
\[
 \frac Lp-\frac{D_L}{p^2}\le\frac{\omega(p)}{p^d}\le\frac Lp,
 \qquad
 D_L=\sum_{k_0\le i<j<k_0+L}3^{i+j}<9^{k_0+L}.
\]
For $p>y:=9^{k_0+L+1}$ we therefore have
\begin{equation}\label{pis:eq:uniformforbidden}
 \frac{L-1}{p}\le\frac{\omega(p)}{p^d}\le\frac Lp.
\end{equation}
These estimates are independent of the shift, even when it varies
with the box.

We use the classical multidimensional large sieve in the form
\begin{equation}\label{pis:eq:largesieve}
 |S|\ll_d\frac{(m+Q^2)^d}{H(Q)},\qquad
 H(Q)=\sum_{\substack{q\le Q\\q\ {\text{squarefree}}}}
 \prod_{p\mid q}\frac{\omega(p)}{p^d-\omega(p)},
\end{equation}
for a set $S$ in a cube of side $m$ avoiding the forbidden classes.
For completeness, its normalization can be seen directly. On
$(\Z/p\Z)^d$ take a mean-zero function equal to $1$ on the allowed
classes and to $-(p^d-\omega(p))/\omega(p)$ on the forbidden classes.
Its normalized squared norm is $1/g(p)$, where
$g(p)=\omega(p)/(p^d-\omega(p))$. The Chinese-remainder product,
Parseval, and Cauchy--Schwarz imply
\[
 g(q)|S|^2\le
 \sum_{\substack{\mathbf a\bmod q\\
 \gcd(a_1,\ldots,a_d,q)=1}}
 \left|\sum_{\mathbf x\in S}e^{2\pi i\mathbf a\cdot\mathbf x/q}\right|^2.
\]
Primitive frequencies with denominators at most $Q$ are
$Q^{-2}$-separated. The analytic large sieve bounds their combined
energy by $O_d((m+Q^2)^d)|S|$, proving
\eqref{pis:eq:largesieve}. This is a classical input, not a new sieve
inequality; see \cite{Kowalski}.

Set $Q=\lfloor\sqrt m\rfloor$, $s=\sqrt{\log Q}$,
$L=\lfloor a s\rfloor$, and $z=\exp(bs)$, with fixed positive
$a,b$ such that
\[
 b>a\log9,\qquad ab\delta<1.
\]
Only eligible primes in $(y,z]$ are used. The prime number theorem
in these fixed residue classes and partial summation give
\[
 \sum_{y<p\le z}^{*}\frac1p
 =\delta\log\frac{b}{a\log9}+o(1),\qquad
 \sum_{p\le z}^{*}\frac{\log p}{p}=(\delta+o(1))\log z.
\]
The star restricts to eligible primes. The moving endpoints cause
no uniform-modulus issue: $B$ is fixed first and both endpoints
tend to infinity.

Expand the finite Euler product $W=\prod_{y<p\le z}^{*}(1+g(p))$
with its nonnegative squarefree weights. Equation
\eqref{pis:eq:uniformforbidden} implies
\[
 \log W\ge
 \left(a\delta\log\frac b{a\log9}+o(1)\right)s.
\]
The weighted mean of $\log q$ in this expansion is exactly
\[
 \sum_{y<p\le z}^{*}\frac{g(p)}{1+g(p)}\log p
 \le (ab\delta+o(1))\log Q.
\]
Since $ab\delta<1$, Markov's inequality leaves a fixed positive
fraction of the weight on $q\le Q$. Thus $H(Q)\gg W$.
Take
\[
 a=\frac1{e\sqrt{\delta\log9}},\qquad
 b=(1-\varepsilon)e\sqrt{\frac{\log9}{\delta}}
\]
with $\varepsilon>0$ sufficiently small. As $\varepsilon\downarrow0$,
the coefficient in the lower bound for $\log H(Q)$, when expressed
in units of $\sqrt{\log m}$, approaches
$e^{-1}\sqrt{\delta/\log3}$. Finally choose a fixed large $B$ so
that this exceeds the prescribed $c<c_0$. We obtain
$H(Q)\gg_{d,c}\exp(c\sqrt{\log m})$.

The cube sieve proves \eqref{pis:eq:exceptionbound} for a cube. A
rectangular box is covered by at most $2^dV/m^d$ such cubes. A point
not in the sifted set has a zero at one of the chosen levels modulo
some $p\le z$. Lemma~\ref{pis:lem:returnbudget} supplies the claimed
certificate and period bound. Bounded values of $m$ are absorbed
in the constants.
\end{proof}

\begin{corollary}[Both totally real floor branches]\label{pis:cor:floorbox}
Among coefficient triples in any box of minimum side $m$, the
totally real cubic Pisot polynomials for which
$\lfloor\beta^{3^n}\rfloor$ fails to be composite on a set of
positive lower density number at most
\[
 O_c\bigl(X_1X_2X_3e^{-c\sqrt{\log m}}\bigr),\qquad c<c_0.
\]
\end{corollary}
\begin{proof}
If the two small real conjugates are $\gamma,\eta$, their odd-power
sum has a constant nonzero sign for all sufficiently large powers.
The only potential exception is $\gamma=-\eta$, which would make
$\beta$ the integer trace, contradicting cubic irreducibility.
Thus the floor is eventually $P_n$ or $P_n-1$. Apply the theorem
to both shifts and take a union bound. The traces tend to infinity.
\end{proof}

\begin{corollary}[Many shifts]\label{pis:cor:manyshifts}
For any finite set $\mathcal H\subset\Z$, the exceptional count for
having a certificate for every $h\in\mathcal H$ is at most
$O_{d,c}(|\mathcal H|V e^{-c\sqrt{\log m}})$.
In particular, for $0<\eta<c<c_0$, one may take
$|h|\le\exp(\eta\sqrt{\log m})$ and retain a saving
$e^{-(c-\eta)\sqrt{\log m}}$.
The progressions for different shifts in this corollary need not
coincide. A genuinely simultaneous progression is proved separately
in Section~\ref{pis:sec:windows}.
\end{corollary}

\section{Two further valuation layers on the twice-square Pell locus}
\label{pis:sec:pell-new}
Consider
\[
 g_u(X)=X^3-uX^2-X+1,\qquad
 4u^2+5u+23=2w^2,\quad u\ge2,\quad w>0.
\]
Let $P_n(u)$ be the sum of the $3^n$-th powers of the roots of $g_u$.
Signs at $-1,0,1,u,u+1$ give roots $-x,y,\beta$ with
$0<x,y<1<u<\beta<u+1$. Neither rational candidate $\pm1$ is a root.
The norm is $-1$, and the second symmetric function $-1$ gives
$\beta(y-x)=-1+1/\beta<0$. Thus $x>y$ and
$P_n(u)=\lfloor\beta^{3^n}\rfloor$ at every level.
Write
\[
 T(u,w)=(17u+12w+10,24u+17w+15),\qquad
 (a_j,r_j)=T^j(9,14),\quad (b_j,s_j)=T^j(1,4).
\]
The Pell reduction in the supplied round~35 identifies the entire
locus with $a_j$ for $j\ge0$ and $b_j$ for $j\ge1$.
Round~36 excludes the first orbit and all second-orbit indices except
$j\equiv0,4374\pmod{13122}$. We have independently replayed the
complete finite certificate for that assertion. The following
three additional witnesses strengthen its conclusion.

\begin{theorem}\label{pis:thm:pell-new}
If $(P_n(u))_{n\ge0}$ is eventually prime on this Pell locus, then
\[
 u=b_j,\qquad39366\mid j,\qquad
 v_3(u-1)\ge10,\qquad u>10^{59049}.
\]
Every parameter excluded by this condition has a positive lower
density of composite trace indices.
\end{theorem}
\begin{proof}
Put $\eta=17+12\sqrt2$ and write $j=2k$. The identity
\[
 (8b_j+5)+4s_j\sqrt2=(13+16\sqrt2)\eta^j
\]
shows that $\eta^{2m}=1\pmod p$ makes $b_{2k}$ periodic modulo $p$
with period dividing $m$. The following table consists of exact
certificates; its column $c$ is $b_{2a}\bmod p$.
\[
\begin{array}{r|r|r|r|r|r}
p&m&a&c&n&v_3((p-1)(p^2-1)(p^3-1))\\\hline
8279983&6561&2187&7051991&4863973&25\\
27398737&19683&6561&15563564&4929815&31\\
15037811&19683&13122&1777082&549606&9
\end{array}
\]
In every row, $p$ is prime, $\eta^{2m}=1\pmod p$, and
$P_n(c)=0\pmod p$. The root-order return lemma makes each such
zero recurrent, because $n$ exceeds the last column. Thus the entire
parameter-index class $k\equiv a\pmod m$ has infinitely many
composite traces, along an arithmetic progression of trace indices.

For reproducibility, put $M_p=(p-1)(p^2-1)(p^3-1)$ and
$E=3^n\bmod M_p$. Binary polynomial powering modulo $(g_c,p)$ gives
$X^E=A+BX+CX^2$, with coefficient triples in the order of the table
\[
\begin{gathered}
(7459560,4961071,5552814),\qquad (21974350,24684017,2338784),\\
(13021067,3973052,8155287).\end{gathered}
\]
Each satisfies $3A+cB+(c^2+2)C=0\pmod p$. This verifies a large
trace-index witness without constructing the integer trace.
Primality of the three modest witness primes is certified by trial
division up to their square roots.

The round~36 cover leaves $k\equiv0,2187\pmod{6561}$.
The first new row removes the latter class. The two remaining rows
remove $k\equiv6561,13122\pmod{19683}$, leaving $19683\mid k$.
Therefore $39366\mid j$.
The exact identity $v_3(b_j-1)=1+v_3(j)$ for positive even $j$
now gives $v_3(u-1)\ge10$. On the Pell curve, $w>4u/3$, so
$T(u,w)_1>33u$. Hence
\[
 u=b_j>33^j\ge33^{39366}=1089^{19683}>10^{59049}.
\]
\end{proof}

\begin{corollary}[Relative count on the complete Pell locus]
Let $\mathcal P(X)$ count the parameters in the displayed Pell locus
with $u\le X$, and let $\mathcal E(X)$ count those with an eventually
prime trace tower. Then
\[
 \mathcal P(X)=\frac{2\log X}{\log(17+12\sqrt2)}+O(1),\qquad
 \mathcal E(X)\le
 \frac{\log X}{39366\log(17+12\sqrt2)}+O(1).
\]
In particular $\limsup_{X\to\infty}\mathcal E(X)/\mathcal P(X)
\le1/78732$.
\end{corollary}
\begin{proof}
Each affine Pell orbit grows as a positive constant times
$(17+12\sqrt2)^j$. The theorem retains at most every $39366$-th
index of the second orbit and none of the first.
\end{proof}

\begin{remark}[The remaining uniform step]
The strengthened numerical and valuation bounds concern this special
parameter family, not the numerical value of Mills' constant.
The exceptional suborbit is not proved empty. Indeed, for any fixed
modulus $M$, let $L$ be the period of the orbit of $(1,4)$ under $T$
modulo $M$. Such a period exists because the linear part of $T$ has
determinant one. At every positive index divisible by $39366L$,
$b_j\equiv1\pmod M$. The trace pair $(1,-1)$ is an exact fixed point
of the cubing recurrence, so $P_n(b_j)\equiv1\pmod M$ for every $n$.
Thus an arbitrary prescribed finite divisibility sieve still leaves
unboundedly many parameters in the surviving branch.
\end{remark}

% Audited synthesis and additional deductions by local_factors agent.
% Sources: Topic 02 Showa R12,R22,R23,R26,R27,R28; 21 Oct 2026.
% R27 and R28 are already in the supplied archive: do not label them new.
\section{Uniform irreducibility and exact finite divisibility patterns}
\label{pis:sec:factor-patterns}

\paragraph{Provenance and improvement.}
The supplied Showa round~27 already proves uniform irreducibility in
an explicit cyclic cubic unit family. It supersedes the earlier
six-term conditional prefix result as an existence theorem, although
the earlier six polynomials belonged to a different monomial family.
Round~28 already realizes every finite prime/composite pattern under
Schinzel's hypothesis H. We reproduce the decisive proofs and add two
refinements: a general signed-specialization criterion with an exact
three-block limitation, and arbitrary prescribed local multiplicities
which recur unchanged on infinite progressions of trace indices.
These refinements do not remove hypothesis H from any assertion of
simultaneous prime values.

\subsection{The balanced unit and its all-index irreducibility}
Put
\[
\begin{gathered}
A=a^3+2a^2+3a+3,\qquad Q=a^2+a+1,\\
f_a(Z)=Z^3-AZ^2-a^2Z+1,\qquad
d(a)=(a+1)(a^2+3)(a^2+3a+3).\end{gathered}
\]
The discriminant is $d(a)^2$. In the cyclic root ordering set
\[
 s=\frac{ar+1}{1-Qr},\qquad t=\frac{r-1}{Qr+a},
 \qquad\beta_a=-r^5s.
\]
For integer $a\ge8$, take $r$ to be the largest root. The rational-root
test and square discriminant show that the field is cyclic cubic.
The trace seed is
\begin{align*}
 U_0={}&a^{14}+9a^{13}+45a^{12}+161a^{11}+450a^{10}
 +1020a^9+1918a^8\\
 &+3021a^7+3987a^6+4379a^5+3942a^4+2820a^3
 +1521a^2+558a+105,\\
 V_0={}&-3a^6-15a^5-40a^4-75a^3-105a^2-90a-33,
\end{align*}
and
\begin{equation}\label{pis:fac:rec}
 U_{n+1}=U_n^3-3U_nV_n-3,\qquad
 V_{n+1}=V_n^3+3U_nV_n+3.
\end{equation}
These identities follow from $rst=-1$ and Newton's identities:
$U_0=-T_{5,1}$ and $V_0=-T_{4,5}$, where
$T_{ij}=r^is^j+s^it^j+t^ir^j$. Thus
$U_n(a)=\operatorname{Tr}(\beta_a^{3^n})$.

\begin{theorem}[Showa round 27, independently audited]
\label{pis:fac:uniform}
Every $U_n$ is monic and irreducible over $\mathbb Q$, of degree
$14\cdot3^n$, and $U_n(-1)=1$. For every integer $a\ge8$ the number
$\beta_a$ is a cubic Pisot unit of norm $-1$ whose conjugates are
$\beta_a,-X,Y$, with
\[
 0<Y<X<1,\quad X^{40}\beta_a^{17}<1,
 \quad X^{23}<Y^{17},\qquad
 U_n(a)=\lfloor\beta_a^{3^n}\rfloor.
\]
\end{theorem}
\begin{proof}
First every odd trace is positive at every real parameter. For real
$a\ne-1$, the three roots of $f_a$ are distinct and real. Their product
is $-1$ and their pairwise sum is $-a^2\le0$, so they have the form
$p>q>0,-w<0$ with $pqw=1$. If $a>-1$, then $f_a(1)<0$ gives
$q<1<p$. The cyclic fractional-linear map sends $p$ to $-w$:
its value cannot vanish or have a pole on this connected parameter
interval, since its determinant is $(a+1)^2\ne0$, and it is negative
at $a=0$. The conjugates of $\beta_a$ are therefore
\[
 P=p^5w,\qquad R=w^5q,\qquad -B=-q^5p,
 \qquad B/P=q^6/p^3<1.
\]
Their odd-power sum is positive. If $a<-1$, then $f_a(1)>0$ and
$A(a)<A(-1)=1$ show that both positive roots lie below one: if both
exceeded one, $A=p+q-1/(pq)>1$. In either cyclic orientation the
negative conjugate has magnitude $B<1$, while $PRB=1$ forces a positive
conjugate above one. The odd trace is again positive. At $a=-1$ the
specialized conjugate multiset is $(1,1,-1)$, and every odd trace is one.

Direct reduction gives
\begin{equation}\label{pis:fac:res}
 \overline U_0=a^5(a-1)^3(a^2+1)^3\quad\text{over }\mathbb F_3.
\end{equation}
For $n\ge1$ the three Hensel blocks of $U_n$ are irreducible over
$\mathbb Q_3$ and have degrees $5\cdot3^n,3\cdot3^n,6\cdot3^n$.
Here is the valuation argument, including repeated residual factors.
At a root $\xi$ of $U_n$, monicity gives integrality and
\[
 v_3(U_j(\xi))=3^{j-n}\quad(0\le j<n).
\]
The last step follows from $U_{n-1}^3=3(1+U_{n-1}V_{n-1})$;
descending steps compare valuations less than one with terms divisible
by three. If a residual irreducible factor $h$ has multiplicity $m$
and degree $e$, write $U_0=H^m B+3C$ with $B(\xi)$ a unit.
Then $v_3(H(\xi))=1/(m3^n)$. The ramification index is divisible by
$m3^n$ and the residue degree by $e$. Their product equals the degree
of the Hensel block, proving its irreducibility.

A monic rational factor $F$ of $U_n$ is integral and has no real root,
so it has even degree and is positive on the real line. The only proper
even subset sums of $5,3,6$ are $6$ and $8$. The corresponding residual
factors are
\[
 (a^2+1)^{3^{n+1}},\qquad
 a^{5\cdot3^n}(a-1)^{3^{n+1}}.
\]
Both have value $-1$ modulo three at $a=-1$. On the other hand $F$ and
its monic complement have positive integer values there multiplying
to $U_n(-1)=1$, so $F(-1)=1$. This contradiction proves irreducibility
for every $n\ge1$.

For the seed, the exact factorizations modulo two have degrees $6,8$,
and those modulo thirteen have degrees $5,9$. For reproducibility,
the factors are specified by ascending coefficient arrays:
\[
\begin{array}{c|l}
2&(1,0,1,1,0,1,1),\quad(1,0,0,1,0,1,1,0,1)\\
13&(12,7,8,8,4,1),\quad(12,7,2,11,12,9,6,4,5,1).
\end{array}
\]
Their products and irreducibility are certified by the Frobenius
criterion $X^{p^d}=X\bmod F$ and
$\gcd(F,X^{p^{d/q}}-X)=1$ for primes $q\mid d$. The disjoint proper
factor-degree sets prove rational irreducibility of $U_0$.

Finally, for $a\ge8$ and $w=-s$,
\[
 a^3<r<2a^3,\quad a/Q<w<1/a,\quad 0<t<1/a^2,
 \quad rwt=1,
\]
so
\[
 (4/5)a^{14}<\beta_a<32a^{14},\quad
 X=t^5r<2a^{-7},\quad Y=w^5t<a^{-7}.
\]
Since $V_0(a)<-1$ and $\beta_aXY=1$,
$\beta_a(Y-X)=V_0(a)+1/\beta_a<0$, giving $X>Y$.
Also $X^{40}\beta_a^{17}<2^{125}a^{-42}\le1/2$, and
$X^{23}/Y^{17}=X^{40}\beta_a^{17}$. Every odd-power trace is its
positive dominant power minus a number strictly between zero and one,
which proves the floor identities.
\end{proof}

\subsection{A signed-specialization criterion and its exact limitation}
\begin{proposition}\label{pis:fac:signcriterion}
Suppose integer polynomials $U_0,V_0$ satisfy \eqref{pis:fac:rec}, $U_0$ is
monic, $\deg V_0<2\deg U_0$, and for some integer $a_*$ one has
$(U_0(a_*),V_0(a_*))=(1,-1)$. Suppose each $U_n$ is strictly positive
on $\mathbb R$. Write
$\overline U_0=\prod_{i=1}^r h_i^{m_i}$ into distinct monic irreducible
factors over $\mathbb F_3$, with $\deg h_i=d_i$. Give block $i$ the label
\[
 \lambda_i=(m_id_i\bmod2,\ h_i(a_*)^{m_i})
 \in(\mathbb Z/2\mathbb Z)\times\mathbb F_3^\times.
\]
If no nonempty proper subcollection of these labels has product
$(0,1)$, then $U_n$ is irreducible over $\mathbb Q$ for every $n\ge1$.
Moreover, this label criterion is possible only for $r\le3$. With
$r=3$, it holds exactly when the labels are the three different
nonidentity elements of this group.
\end{proposition}
\begin{proof}
The root-valuation proof above gives one irreducible local block of
degree $m_id_i3^n$ and reduction $h_i^{m_i3^n}$ for each $i$.
Since $3^n$ is odd, its parity and value at $a_*$ retain precisely the
label $\lambda_i$. A proper rational factor must have even degree and
positive integer value one at $a_*$, contradicting the label condition.
The product of all labels is $(0,1)$, because $U_0$ has even degree
and value one at $a_*$. The group has three nonidentity elements,
each of order two. A minimal product-one sequence cannot contain the
identity (unless it has length one), or a repeated element (unless
it has length two). Thus it has length at most three; the three-element
case must consist of all three nonidentity elements. Conversely those
three labels have product one and no proper product-one subcollection.
\end{proof}
This is a limitation of this particular factor obstruction, not a
bound on the number of local factors of a globally irreducible polynomial.
For the balanced seed the labels are $(1,-1),(1,+1),(0,-1)$, so the
criterion is extremal. The four-block degree-$31$ seed of the earlier
six-prefix argument cannot be handled by this exact criterion alone.

\subsection{Simple forcing roots and exact local multiplicities}
We next state a strengthening of round~28. All prime-value assertions
below explicitly assume the usual integer-prime hypothesis H.

\begin{lemma}[Forcing roots, from round 28]\label{pis:fac:forcing}
For every finite window $I=\{0,\ldots,L\}$ and $j\in I$, a positive-
Dirichlet-density set of primes $p\equiv2\pmod3$ has a residue $b$ such
that
\[
 U_j(b)=0,\quad U'_j(b)\ne0,\quad d(b)\ne0,
 \qquad U_i(b)\ne0\quad(i\in I\setminus\{j\}).
\]
Any prescribed finite set of primes may be excluded.
\end{lemma}
\begin{proof}
Let $K_j=\mathbb Q(\xi)$ for a root of $U_j$. It has an odd-degree
completion above three. For $j\ge1$ this follows from the local blocks
of degrees $5\cdot3^j$ and $3\cdot3^j$. For $j=0$, Hensel lifts the
residual factor $a^5$ to a degree-five block; at least one of its
irreducible constituents has odd degree. Thus
$\mathbb Q(\sqrt{-3})$ is not a subfield of $K_j$, since its completion
at three has degree two and would embed in every completion of $K_j$.

In a Galois extension containing the splitting field of $U_j$ and
$\mathbb Q(\sqrt{-3})$, choose an automorphism fixing $\xi$ and acting
nontrivially on the quadratic field. It exists by the preceding
noncontainment. Its conjugacy class has a fixed root and the same
nontrivial quadratic character. Chebotarev gives a positive-density
set of primes with a root of $U_j$ and $p\equiv2\pmod3$.
Exclude the finitely many divisors of $\operatorname{disc}(U_j)$,
$\operatorname{Res}(U_j,d)$, and
$\operatorname{Res}(U_j,U_i)$ for $i\ne j$. All these integers are
nonzero by irreducibility and unequal degrees. The asserted simple
root and all other nonvanishing conditions follow.
\end{proof}

\begin{theorem}[Prescribed exact multiplicities with periodic recurrence]
\label{pis:fac:valuationpatterns}
Fix a finite window $I=P\sqcup C$, a positive integer $M$, and for each
$j\in C$ a nonempty finite list of positive integers
$e_{j,1},\ldots,e_{j,r_j}$. There are pairwise distinct primes
$p_{j,k}\ge5$ with $p_{j,k}\equiv2\pmod3$, all coprime to $M$, and an arithmetic progression
$\mathcal A$ of integer parameters with the following unconditional
properties for every sufficiently large $a\in\mathcal A$:
\begin{enumerate}
\item $a\equiv-1\pmod M$, and every odd-power trace of $\beta_a$ is
$1$ modulo $M$;
\item for $i\in I$,
\[
 v_{p_{j,k}}(U_i(a))=
 \begin{cases}e_{j,k},&i=j,\\0,&i\ne j;\end{cases}
\]
\item with
\[
 T_{j,k}=\operatorname{ord}_{\,p_{j,k}^{e_{j,k}}
                      (p_{j,k}^3-1)}(3),
\]
one has
\[
 v_{p_{j,k}}(U_{j+hT_{j,k}}(a))=e_{j,k}\qquad(h\ge0).
\]
\end{enumerate}
Assuming hypothesis H, infinitely many $a\in\mathcal A$ additionally
make every $U_i(a)$, $i\in P$, prime. Every prescribed composite position
$j$ then has at least $r_j$ different prime divisors, with exactly the
specified multiplicities for the chosen divisors. If $P$ is empty,
no conjecture is required.
\end{theorem}
\begin{proof}
For each pair $(j,k)$ choose a fresh prime $p\ge5$ and simple forcing root
$(p,b)$ from Lemma~\ref{pis:fac:forcing}. Hensel's lemma lifts $b$ uniquely
to a root modulo $p^{e+1}$, where $e=e_{j,k}$. Among the $p$ lifts of
the root modulo $p^e$, choose any of the other $p-1$ classes modulo
$p^{e+1}$. Call it $c$. Taylor expansion and $U'_j(b)\ne0$ give
$v_p(U_j(c))=e$; every other designated $U_i(c)$ is a unit modulo $p$.
CRT now imposes $a\equiv-1\pmod M$ and
$a\equiv c_{j,k}\pmod{p_{j,k}^{e_{j,k}+1}}$ simultaneously. The resulting
progression has modulus
\[
 D=M\prod_{j\in C}\prod_{k=1}^{r_j}p_{j,k}^{e_{j,k}+1}.
\]
This proves the exact multiplicities within the finite window.

We verify recurrence at the same precision. Fix a chosen $p,e$ and
consider
\[
 \mathcal B=(\mathbb Z/p^{e+1}\mathbb Z)[Z]/(f_a(Z)).
\]
Because $p\nmid d(a)$, its reduction is etale. The square discriminant
forces splitting type $1+1+1$ or $3$, so every residue-field unit is
killed by $p^3-1$. The kernel of reduction on $\mathcal B^\times$ is
$1+p\mathcal B$ and is killed by $p^e$, by induction with the binomial
theorem. Hence every unit of $\mathcal B$ is killed by
$E=p^e(p^3-1)$. The root transformation defining $s$ has invertible
denominator: a simultaneous zero of $1-QZ$ and $aZ+1$ would force
$(a+1)^2=0\bmod p$, excluded by $d(a)\ne0$. Thus $\beta=-r^5s$ is a unit
in this algebra, and its traces agree with the integer polynomials.
Since $p\equiv2\pmod3$, $\gcd(E,3)=1$. For $T=\operatorname{ord}_E(3)$,
\[
 U_{n+T}(a)\equiv U_n(a)\pmod{p^{e+1}}\qquad(n\ge0).
\]
The nonzero residue of $U_j(a)/p^e$ modulo $p$ is therefore retained
at every $j+hT$, proving the exact recurring valuations.

Write the CRT progression as $a=A_0+Dz$. For each $i\in P$, the
polynomial $U_i(A_0+Dz)$ is irreducible, primitive, and has positive
leading coefficient. Their product has no fixed prime divisor:
primes dividing $M$ give values one; each forcing prime gives nonzero
values at every prime-designated index; at a prime not dividing $D$
choose $z$ so that $A_0+Dz\equiv-1$. Hypothesis H consequently supplies
infinitely many simultaneous prime values. All other claims are
unconditional congruences. Finally, every odd trace is an integer
polynomial in $a$ and has value one at $a=-1$, proving the modulo-$M$
assertion. Taking $a$ large ensures the Pisot bounds and makes each
forced value larger than its chosen prime divisors.
\end{proof}

The theorem prescribes valuations of selected divisors, not a complete
prime factorization: unselected prime factors may also occur. It
strengthens the finite-pattern statement in the archive while retaining
its essential finite-window quantifier and its explicit conjectural
prime-value assumption.

\subsection{Audit boundaries and remaining arithmetic question}
The single-block norm-minus-one family of round~22 also has all-index
local irreducibility, but its index-zero trace is a fourth power; its
prime-prefix statement deliberately begins at the next index. The
cyclic unit family above is stronger for conditional prime-prefix
existence and starts at index zero. Round~23 is compatible with the
present recurrence theorem: its isolated divisors are chosen at primes
congruent to one modulo three, while the recurring divisors above lie
at primes congruent to two modulo three. In round~23 the prime and unit
may depend on the isolated index; the result does not say that every
trace has such a divisor in one fixed tower.

Round~26 proves a separate polynomial-family obstruction: fixed nonzero
norm and square discriminant force the polynomial coefficient functions
to be algebraically dependent and, over $\mathbb C$, to factor through
one common polynomial. Its logarithmic two-form argument and the
normalization-to-an-affine-line argument are valid; they do not prevent
the one-parameter construction used here. We do not use a conjectural
cyclic lift of the round~22 family, the multivariate Hilbert construction
of round~24, or an unverified primality certificate from round~25 as
inputs to the proofs above.

Uniform irreducibility is an unconditional result already in the
archive. Arbitrarily long prime prefixes remain conditional on H;
exact prescribed local divisibility is unconditional. None implies
one fixed algebraic number whose entire infinite trace tower is prime.
Nor does any identify the least-prime path defining Mills' constant.
The full cubic Pisot compositeness problem and unconditional Mills
transcendence remain open within these arguments.

% Transfer audit and a quantitative ray-unit barrier.
% Written 21 October 2026. This is a section fragment, not a standalone file.
% New relative to the audited older [031] and supplied non-Topic02 notes;
% no literature-wide priority assertion is intended.

\section{Finite local conditions and logarithmic unit lattices}
\label{pis:transfer:sec:rayunits}

The new non-Mills material does not supply an additional proved route to
Mills transcendence. A useful transfer is nevertheless possible: the
unit-lattice viewpoint gives a quantitative explanation of why finitely
many local conditions and Saito's strict conjugate bounds can coexist.
It identifies a genuine limitation of that combination of tests without
assuming that any trace is prime.

Let $K$ be a totally real cubic field, with real embeddings
$\sigma_1,\sigma_2,\sigma_3$. Fix a unit $w\in\mathcal O_K^\times$ with
\[
 \sigma_1(w)>0,\qquad \sigma_2(w)<0,\qquad \sigma_3(w)>0.
\]
Its norm is $-1$. Let $\mathfrak m$ be a nonzero integral ideal, and set
\[
 E^+(\mathfrak m)=
 \{u\in\mathcal O_K^\times:
       \sigma_i(u)>0\ (1\le i\le3),\ u\equiv1\pmod{\mathfrak m}\}.
\]
Its two-coordinate logarithmic lattice is
\[
 \Lambda_{\mathfrak m}=
 \{(\log\sigma_1(u),\log\sigma_2(u)):u\in E^+(\mathfrak m)\}
 \subset\mathbb R^2.
\]
Write $\Delta_{\mathfrak m}>0$ for its covolume in these coordinates.
This convention avoids an extraneous factor of $\sqrt3$ from embedding
the lattice into the plane $x+y+z=0$ in $\mathbb R^3$.

\begin{theorem}[Quantitative persistence in every admissible unit coset]
\label{pis:transfer:thm:ray}
Fix $0<a<1/2$. For $H>1$, let $B_{w,\mathfrak m,a}(H)$ count the units
$\beta\in wE^+(\mathfrak m)$ such that, on writing
\[
 B=\sigma_1(\beta),\qquad r=\sigma_2(\beta),\qquad
 s=\sigma_3(\beta),
\]
one has
\[
 1<B\le H,\qquad 0<s<-r<1,\qquad -r<B^{-a}.
\]
Then
\begin{equation}\label{pis:transfer:eq:count}
 B_{w,\mathfrak m,a}(H)
 =\frac{1-2a}{4\Delta_{\mathfrak m}}(\log H)^2
     +O_{K,w,\mathfrak m,a}(\log H+1).
\end{equation}
Every counted $B$ is a cubic Pisot number of norm $-1$. In addition,
\[
 -r<s^{a/(1-a)},\qquad \beta\equiv w\pmod{\mathfrak m}.
\]
In particular, for $a=17/40$, the two strict bounds
\[
 -r<\min(B^{-17/40},s^{17/23})
\]
hold for quadratically many units in logarithmic height in every such
coset.
\end{theorem}

\begin{proof}
The subgroup $E^+(\mathfrak m)$ has finite index in the unit group:
positivity imposes only a finite sign condition, and reduction modulo
$\mathfrak m$ maps units to the finite group
$(\mathcal O_K/\mathfrak m)^\times$. Dirichlet's unit theorem therefore
makes $\Lambda_{\mathfrak m}$ a full lattice of rank two. Positivity
removes its possible torsion, so its logarithm map is injective.

The logarithms of the coset $wE^+(\mathfrak m)$ form the translate
\[
 (\log\sigma_1(w),\log|\sigma_2(w)|)+\Lambda_{\mathfrak m}.
\]
For a point of this translate put $x=\log B$ and $y=\log(-r)$.
Because $|\operatorname N\beta|=1$, one has $\log s=-x-y$.
The asserted inequalities are exactly
\begin{equation}\label{pis:transfer:eq:triangle}
 0<x\le\log H,\qquad -x/2<y<-ax.
\end{equation}
Indeed $y>-x/2$ is $-r>s$, and $y<-ax<0$ gives $-r<1$;
then $-x-y<-x/2<0$, so $s<1$ as well.
The triangle in \eqref{pis:transfer:eq:triangle} has area
\[
 \int_0^{\log H}(1/2-a)x\,dx
       =\frac{1-2a}{4}(\log H)^2
\]
and boundary length $O_a(\log H)$. Counting translates of a fixed
fundamental parallelogram gives its area divided by
$\Delta_{\mathfrak m}$, with error $O(\log H+1)$: only cells meeting a
fixed-width neighborhood of the boundary can contribute to the error.
Strict versus non-strict boundaries changes the count by at most the
same order. This proves \eqref{pis:transfer:eq:count}.

The dominant embedding exceeds one and the other two have modulus
less than one, so $B$ is Pisot. It cannot be rational, since a rational
unit is $\pm1$; its degree is therefore three. Finally
$B(-r)s=1$, and the inequality $-r<B^{-a}$ is equivalent to
$(-r)^{1-a}<s^a$, proving the additional bound. The congruence is the
one defining the coset.
\end{proof}

\begin{corollary}[Persistence of every fixed finite trace sieve]
\label{pis:transfer:cor:sieve}
Take $\mathfrak m=M\mathcal O_K$ with integer $M\ge2$. Every unit
counted in Theorem~\ref{pis:transfer:thm:ray} satisfies
\[
 \operatorname{Tr}_{K/\mathbb Q}(\beta^j)
 \equiv\operatorname{Tr}_{K/\mathbb Q}(w^j)\pmod M
 \qquad(j\ge1).
\]
Thus, if a seed $w$ has
$\gcd(\operatorname{Tr}(w^{3^n}),M)=1$ at every level $n\ge0$,
then all the counted balanced Pisot units have the same avoidance
property at every level.
\end{corollary}
\begin{proof}
From $\beta-w\in M\mathcal O_K$ one obtains
$\beta^j-w^j\in M\mathcal O_K$ for every positive integer $j$.
The trace of an algebraic integer is an integer, so tracing gives the
stated congruence. The last assertion follows separately at every
level; no primality assertion is inferred from coprimality with $M$.
\end{proof}

The seed hypothesis in this corollary is essential. The theorem does
not assert that an arbitrary finite sieve has an avoiding seed in a
specified field. Instead it proves that, once a finite residue class
is admissible, the archimedean restrictions alone cannot eliminate it:
there are asymptotically $c(\log H)^2$ members satisfying the strict
conjugate bounds. The regulator enters only through the positive
covolume $\Delta_{\mathfrak m}$ and changes the leading constant,
not the existence of these candidates. This is stronger than an
unquantified ray construction but does not close the remaining
norm-minus-one cyclic Pisot case.


\begin{example}[An explicit nonempty congruence coset]
\label{pis:transfer:ex:seed}
Let $r$ be the largest root of $X^3-25X^2-4X+1$ and order its other
roots as $s<0<t$. The rational root test excludes $\pm1$, the
discriminant is $273^2$, and root brackets give
\[
 25<r<26,\qquad -0.294<s<-0.292,\qquad 0.135<t<0.137.
\]
Thus $K=\mathbb Q(r)$ is cyclic cubic and $r$ has signs $(+,-,+)$
and norm $-1$. Direct root enumeration of the polynomial over
$\mathbb F_2,\mathbb F_5,\mathbb F_{17}$ shows it irreducible in all
three cases; these primes are inert. The integral element
$v=(r^2+2)/3$ satisfies
\[
 v^3-213v^2+290v-99=0.
\]
This polynomial reduces to $X(X-1)(X+1)$ modulo three, so $3$ splits
completely. The residues of $r$ at its three primes are
$(-1,1,1)$ in some order, as follows from its norm $-1$ and trace
$25\equiv1\pmod3$.

Put $w=r^{66619}$, where $66619=7\cdot31\cdot307$.
At $2$, the order of $r$ divides $7$; at $5$ and $17$, the finite-field
norm identities give $r^{31}=-1$ and $r^{307}=-1$, respectively.
The relevant quotients of $66619$ are odd. Consequently, for every
odd integer $j\ge1$,
\[
 \operatorname{Tr}(w^j)\equiv
 \begin{cases}
 1&\pmod2,\\
 1&\pmod3,\\
 2&\pmod5,\\
 14&\pmod{17}.
 \end{cases}
\]
Applying Theorem~\ref{pis:transfer:thm:ray} to
$wE^+(510\mathcal O_K)$ gives asymptotically
$3(\log H)^2/(80\Delta_{510\mathcal O_K})$ strictly Saito-balanced
Pisot units whose entire cubic-power trace towers avoid
$2,3,5,17$. This asserts finite-prime avoidance, not primality.
\end{example}

% Section fragment for the integrated revision; written 21 October 2026.
% Uses thm:windows, lem:contraction, prop:fourteen, Keio009, and S14.
\section{A complete cyclic cubic norm class}
\label{pis:rescyc:sec:cyclic}
Submitted round S14 already proves that an eventually prime,
growing trace tower in a cyclic cubic field must have norm congruent
to zero or minus one modulo three \cite[Theorem~1]{S14}. Thus the
norm class excluded below is not newly discovered here. Combining
that restriction with the trace-limit classification and cyclic
Frobenius-return argument of \cite{Keio009}, we record a
positive-density formulation for the actual Pisot floors, including
both possible real floor branches. This is a synthesis of existing
ingredients, not a new return mechanism or a claim of historical
priority.

\begin{theorem}[Cyclic cubic norms congruent to one]
\label{pis:rescyc:thm:norm-one-mod-three}
Let $K/\Q$ be a cyclic cubic extension, and let $\beta\in K$ be a
cubic Pisot number. If
\[
 N=\Nm_{K/\Q}(\beta)\equiv1\pmod3,
\]
then the set
\[
 \mathcal C_\beta=
 \{n\ge0:\lfloor\beta^{3^n}\rfloor\text{ is composite}\}
\]
has positive lower density. In particular, this conclusion holds
for every cyclic cubic Pisot unit of norm $+1$.
\end{theorem}

\begin{proof}
A cyclic cubic field is totally real, since a nontrivial complex
conjugation would have order two in a group of order three. Write
the conjugates of $\beta$ as $\beta,\gamma_2,\gamma_3$, and put
\[
 U_n=\Tr_{K/\Q}(\beta^{3^n}),\qquad
 f(X)=X^3-uX^2+vX-N.
\]
The two small conjugates have different absolute values: equality
would give $\gamma_2=-\gamma_3$ and hence $\beta=u\in\Z$,
contrary to degree three. Thus their odd-power sum has a fixed
nonzero sign. For all sufficiently large $n$,
\[
 \lfloor\beta^{3^n}\rfloor=U_n-s,
 \qquad s\in\{0,1\},
\]
where $s=1$ exactly when the small conjugate of larger absolute
value is positive. Also $U_n\to+\infty$.
Let $\lambda=\lim_{n\to\infty}U_n$ in $\Z_3$.

If $s=1$, Proposition~\ref{pis:prop:fourteen} and its explicit table
show that either equality $\lambda=0$ or $\lambda=2$ would force
$3\mid N$. Since $N\equiv1\pmod3$, neither equality is possible.
Theorem~\ref{pis:thm:windows}, with the single shift $h=-1$, supplies
an eventual arithmetic progression of composite floors.

Suppose now that $s=0$. If $\lambda\notin\{1,-1\}$,
Theorem~\ref{pis:thm:windows} with $h=0$ gives the same conclusion.
We handle the two remaining limits separately.

If $\lambda=1$, the residue classification gives
\[
 f(X)\equiv(X-1)(X^2+1)\pmod3,
\]
because $N\equiv1\pmod3$. Its root $1$ is simple, so Hensel's
lemma gives an embedding $K=\Q(\beta)\hookrightarrow\Q_3$.
In a Galois extension all primes over three have the same local
degree; the degree-one completion therefore implies that three
splits completely in $K$. The three conjugates of $\beta$ now lie
in $\Q_3$ and are units, because $3\nmid N$. Each has
Teichm\"uller part $1$ or $-1$. Their sum is $\lambda=1$, so
these three parts must be $1,1,-1$ in some order. Their product
gives $N\equiv-1\pmod3$, a contradiction. Equivalently, the
displayed separable reduction has factorization type $1+2$,
which cannot occur at an unramified prime in a cyclic cubic field.

It remains to treat $\lambda=-1$. Lemma~\ref{pis:lem:contraction}
gives $U_n\equiv-1\pmod3$ for every $n$. If only finitely many
$U_n$ are prime, then every sufficiently late floor is composite,
and the density is one. Otherwise choose a sufficiently late
prime $p=U_r$ with
\[
 p\nmid3N\operatorname{disc}(f).
\]
This is possible since $U_n\to+\infty$. Frobenius in the cyclic
cubic field has cycle type $1+1+1$ or $3$. Consequently all roots
of $f$ modulo $p$ lie in $\F_{p^3}$ and are nonzero. Since
$p\equiv-1\pmod3$,
\[
 \gcd(3,p^3-1)=1.
\]
Choose $D\ge1$ with $3^D\equiv1\pmod{p^3-1}$. For every
residual root $a$ and every $j\ge0$ this gives
$a^{3^{r+jD}}=a^{3^r}$. Summing over the three roots yields
\[
 U_{r+jD}\equiv U_r=0\pmod p\qquad(j\ge0).
\]
For sufficiently large $j$ the floor equals $U_{r+jD}>p$, so
this is an arithmetic progression of composite floors. Every
progression obtained above has positive lower density, completing
the proof.
\end{proof}

\begin{remark}[What remains in the cyclic Mills branch]
For a hypothetical cyclic cubic shadow of the least Mills
constant, the dominant small conjugate is negative. An eventually
prime trace tail therefore forces $\lambda=1$: other limits are
excluded by Theorem~\ref{pis:thm:windows} or by the last part of the
proof above. If its norm is prime to three, complete splitting at
three then forces $N\equiv-1\pmod3$. In particular, a unit
shadow must have norm $-1$. The norm-$-1$ cyclic units constructed
elsewhere in this paper show why Theorem~\ref{pis:rescyc:thm:norm-one-mod-three}
does not settle the entire cyclic case or unconditional Mills
transcendence.
\end{remark}


\section{Excluding rational powers of Mills' constant from number fields}
\label{pis:res:sec:fields}

Restricted cubic results already give complete exclusion statements for
other, substantially larger classes of algebraic expressions. Throughout
this section $A$ is the least Mills constant defined in the introduction.
The cubic three-adic embedding obstruction below was already proved in
the supplied manuscript \cite{Keio031}. We include a short proof from
Theorem~\ref{pis:thm:windows} and Proposition~\ref{pis:prop:fourteen} to make the
subsequent field deductions self-contained. The ambient-field and
cyclotomic consequences are deductions from this obstruction; no
historical priority is asserted for them.

\begin{lemma}[Prime traces force a three-adic embedding]
\label{pis:res:lem:prime-embedding}
Let $\beta$ be a cubic Pisot number, with minimal polynomial
$f(X)=X^3-uX^2+vX-N$, and put
$U_n=\Tr_{\Q(\beta)/\Q}(\beta^{3^n})$.
If $U_n$ is prime for every sufficiently large $n$, then its three-adic
limit satisfies $\lambda\in\{1,-1\}$, and $\Q(\beta)$ admits an
embedding into $\Q_3$. No restriction on the norm $N$ is required.
\end{lemma}
\begin{proof}
The Pisot property gives $U_n=\beta^{3^n}+o(1)$, so $U_n\to\infty$.
If $\lambda\notin\{1,-1\}$, apply Theorem~\ref{pis:thm:windows} with the
single shift $0$. It gives infinitely many composite $U_n$, contrary
to the hypothesis. Write $\lambda=\varepsilon$, where
$\varepsilon\in\{1,-1\}$. The six corresponding coefficient classes
in Proposition~\ref{pis:prop:fourteen} satisfy
\[
 u\equiv\varepsilon,\qquad N\equiv\varepsilon v\pmod3,
 \qquad
 f(X)\equiv(X-\varepsilon)(X^2+v)\pmod3.
\]
If $v\not\equiv2\pmod3$, the residual root $\varepsilon$ is simple.
If $v\equiv2\pmod3$, the residual polynomial is
$(X-\varepsilon)^2(X+\varepsilon)$, whose root $-\varepsilon$ is simple.
Hensel's lemma therefore gives a root of $f$ in $\Z_3$. Since $f$ is
irreducible over $\Q$, evaluation at this root embeds
$\Q[X]/(f)=\Q(\beta)$ into $\Q_3$.
\end{proof}

\begin{lemma}[Pisot powers retain their field]
\label{pis:res:lem:power-field}
If $\beta$ is a Pisot number and $r$ is a positive integer, then
$\Q(\beta^r)=\Q(\beta)$.
\end{lemma}
\begin{proof}
An embedding of $\Q(\beta)$ into $\mathbb C$ that fixes
$\Q(\beta^r)$ sends $\beta$ to a conjugate $\gamma$ with
$\gamma^r=\beta^r$. Hence $|\gamma|=\beta>1$. The defining Pisot
condition leaves only $\gamma=\beta$. Separability in characteristic
zero says that the number of such embeddings is
$[\Q(\beta):\Q(\beta^r)]$, which is therefore one.
\end{proof}

\begin{theorem}[An obstruction in every ambient number field]
\label{pis:res:thm:ambient}
Let $K\subseteq\mathbb C$ be a number field. If $A^q\in K$ for some
$q\in\Q\setminus\{0\}$, then $K$ contains a totally real cubic
subfield $F$ admitting an embedding into $\Q_3$. Consequently, if
$K$ has no such subfield, then
\[
 \boxed{A^q\notin K\qquad(q\in\Q\setminus\{0\}).}
\]
Here rational powers of $A>1$ denote their positive real values.
\end{theorem}
\begin{proof}
Replacing $q$ by $-q$ if necessary, write $q=r/s$ with positive
integers $r,s$. The inclusion $A^{r/s}\in K$ makes $A$ algebraic.
Saito's reduction \cite[Theorem~1.7]{SaitoVariants} supplies
$\beta=A^{3^m}$, a totally real cubic Pisot number satisfying
\eqref{pis:eq:saitobounds}. Its two small conjugates have negative sum
at every sufficiently large odd power: the negative one has strictly
larger absolute value. Thus
\[
 \lfloor\beta^{3^n}\rfloor
 =\Tr_{\Q(\beta)/\Q}(\beta^{3^n})
\]
for every sufficiently large $n$. These integers are prime by the
definition of $A$. Lemma~\ref{pis:res:lem:prime-embedding} therefore
gives a $\Q_3$-embedding of $F=\Q(\beta)$. Finally,
\[
 \beta^r=(A^{r/s})^{s3^m}\in K,
 \qquad \Q(\beta^r)=F
\]
by Lemma~\ref{pis:res:lem:power-field}, proving $F\subseteq K$.
\end{proof}

\begin{corollary}[Cubic fields and radical expressions]
\label{pis:res:cor:cubic}
If a cubic number field $K$ has no embedding into $\Q_3$, then no
nonzero rational power of $A$ belongs to $K$. In particular this holds
when $3$ is inert or totally ramified in $K$. Equivalently, no positive
element of such a field, raised to a nonzero rational power, can
equal $A$.
\end{corollary}
\begin{proof}
Any cubic subfield of the cubic field $K$ is $K$ itself, so
Theorem~\ref{pis:res:thm:ambient} applies. Inertness and total ramification
both exclude a prime of $K$ over $3$ with ramification index and
residue degree equal to one, which an embedding into $\Q_3$ would
require.
\end{proof}

\begin{theorem}[A Galois prime-count obstruction]
\label{pis:res:thm:galois}
Let $K/\Q$ be a finite Galois extension and let $g_3(K)$ be the
number of prime ideals of $K$ above $3$. If $3\nmid g_3(K)$, then
$A^q\notin K$ for every $q\in\Q\setminus\{0\}$.
\end{theorem}
\begin{proof}
Suppose otherwise and let $F\subseteq K$ be the cubic subfield from
Theorem~\ref{pis:res:thm:ambient}. Put
$G=\operatorname{Gal}(K/\Q)$ and $H=\operatorname{Gal}(K/F)$, so
$[G:H]=3$. The $\Q_3$-embedding supplies a prime $\mathfrak p$ of
$F$ with completion $F_{\mathfrak p}=\Q_3$. Choose a prime
$\mathfrak P$ of $K$ above $\mathfrak p$, and let $D\subseteq G$ be
its decomposition group. The image of $F$ in $K_{\mathfrak P}$ lies
in $\Q_3$, so every member of $D$ fixes $F$ and $D\subseteq H$.
The usual transitive action on the primes above $3$ gives
\[
 g_3(K)=[G:D]=[G:H][H:D]=3[H:D],
\]
contrary to the hypothesis.
\end{proof}

\begin{remark}[The exact abelian scope of the local obstruction]
\label{pis:res:rem:abelian}
If $K/\Q$ is finite abelian, it contains a cyclic cubic subfield in
which $3$ splits completely if and only if $3\mid g_3(K)$.
Indeed, its decomposition group $D$ is normal. An index-three
subgroup $H$ containing $D$ exists precisely when the finite abelian
group $G/D$ has order divisible by three. Its fixed field is cyclic
cubic, is totally real because its degree is odd, and has trivial
decomposition group at $3$. This equivalence concerns the existence
of a subfield with the required local property; it does not assert
that a Mills candidate exists in that subfield.
\end{remark}

\begin{corollary}[Cyclotomic exclusions]
\label{pis:res:cor:cyclotomic}
Write $N=3^b m$, where $b\ge0$ and $(m,3)=1$. If
\[
 3\nmid\frac{\varphi(m)}{\ord_m(3)},
\]
then $A^q\notin\Q(\zeta_N)$ for every nonzero rational $q$.
Here $\zeta_N$ is a primitive $N$th root of unity, and
$\varphi(1)=\ord_1(3)=1$ by convention. In particular, for every
nonzero rational $q$,
\[
 A^q\notin\bigcup_{a,b\ge0}\Q(\zeta_{7^a3^b}),
 \qquad
 A^q\notin\bigcup_{a,b\ge0}\Q(\zeta_{13^a3^b}).
\]
\end{corollary}
\begin{proof}
In $\Q(\zeta_m)$ the prime $3$ is unramified and its Frobenius has
order $\ord_m(3)$, so it has
$\varphi(m)/\ord_m(3)$ primes above it. The $3$-power cyclotomic
factor is totally ramified at $3$. Locally, adjoining this totally
ramified factor to the unramified completion multiplies the local
degree by $\varphi(3^b)$, just as the global degree is multiplied
by $\varphi(3^b)$. Thus it does not change the number of primes:
\[
 g_3\bigl(\Q(\zeta_{3^b m})\bigr)
 =\frac{\varphi(m)}{\ord_m(3)}.
\]
Theorem~\ref{pis:res:thm:galois} proves the first assertion.

For $a\ge1$, elementary lifting of multiplicative orders gives
\[
 \ord_{7^a}(3)=6\cdot7^{a-1},\qquad
 \ord_{13^a}(3)=3\cdot13^{a-1}.
\]
Indeed, the orders modulo $7$ and $13$ are $6$ and $3$, respectively,
and $v_7(3^6-1)=v_{13}(3^3-1)=1$. The corresponding values of
$g_3$ are $1$ and $4$, neither divisible by three. The case $a=0$
has $g_3=1$. Every element of either displayed union lies in one
finite stage, proving the two exclusions.
\end{proof}

\begin{corollary}[A quotient obstruction]
\label{pis:res:cor:quotient}
Let $K/\Q$ be finite Galois. If
$\operatorname{Gal}(K/\Q)$ has no quotient isomorphic to $C_3$ or
$S_3$, then $A^q\notin K$ for every nonzero rational $q$.
This applies, in particular, when the Galois group is perfect, or
is $S_d$ with $d\ge5$.
\end{corollary}
\begin{proof}
If a cubic subfield $F$ existed, its normal closure would lie in
$K$ and have a transitive cubic Galois group, necessarily $C_3$ or
$S_3$. Restriction would give the prohibited quotient. Apply
Theorem~\ref{pis:res:thm:ambient}. Quotients of a perfect group are
perfect, whereas $C_3$ and $S_3$ are not. For $d\ge5$, the only
proper nontrivial quotient of $S_d$ is $C_2$.
\end{proof}

\begin{remark}[Composita and the remaining cubic problem]
\label{pis:res:rem:limits}
These exclusions cannot be combined by simply taking composita.
For example, $\Q(\zeta_7)$ and $\Q(\zeta_{13})$ separately satisfy
Theorem~\ref{pis:res:thm:galois}, but their compositum has
\[
 g_3\bigl(\Q(\zeta_{91})\bigr)
 =\frac{72}{\operatorname{lcm}(6,3)}=12.
\]
By Remark~\ref{pis:res:rem:abelian}, that compositum contains a cyclic
cubic subfield in which $3$ splits completely.

The section excludes entire fields and all rational powers at once,
without a unit or norm hypothesis. Nevertheless, it does not eliminate
the cubic shadows already satisfying $\lambda=\pm1$, a
$\Q_3$-embedding, and Saito's real-conjugate bounds. In the cyclic
case such a shadow must have $3$ completely split. Excluding every
remaining Mills shadow would suffice for unconditional transcendence
of $A$, even without solving the stronger question for every cubic
Pisot number. No such final exclusion is obtained here.
\end{remark}


% New section: a uniform finite-prime barrier in every admissible field.
% This file uses the main manuscript's theorem environments and macros.
\section{The exact field boundary for a fixed finite set of prime divisors}
\label{pis:res:sec:barrier}

The unit-coset result above assumes an admissible seed. Allowing a
fixed nonzero integral norm, rather than requiring norm minus one,
removes that assumption in every remaining totally real cubic field.
The resulting statement concerns divisibility by a fixed finite set
of primes; it makes no assertion that the surviving integers are prime.

For a cubic algebraic integer $\beta$, write
\[
 \lambda(\beta)=\lim_{n\to\infty}
       \Tr_{\Q(\beta)/\Q}(\beta^{3^n})\quad\hbox{in }\Z_3.
\]
Its existence and its dependence on the residual polynomial are the
three-adic convergence and classification in
Proposition~\ref{pis:prop:fourteen} and its preceding construction.

\begin{theorem}[Every admissible field survives every fixed finite prime sieve]
\label{pis:res:thm:equivalence}
Let $K$ be a totally real cubic field, with real embeddings
$\sigma_1,\sigma_2,\sigma_3$. The following are equivalent.
\begin{enumerate}
\item There is a field embedding $K\hookrightarrow\Q_3$.
\item For every finite set $P$ of rational primes, there are infinitely
many $\beta\in\mathcal O_K$ such that, on writing
\[
 B=\sigma_1(\beta),\qquad
 \sigma_2(\beta)=-\rho,\qquad \sigma_3(\beta)=s,
\]
one has
\begin{gather}
 B>1,\qquad 0<s<\rho<1,\qquad
 \rho<\min\{B^{-17/40},s^{17/23}\},
 \label{pis:res:eq:strictbounds}\\
 \lambda(\beta)=1,
 \qquad
 \gcd\left(\lfloor B^m\rfloor,\prod_{p\in P}p\right)=1
 \quad\text{for every positive odd integer }m.
 \label{pis:res:eq:avoid}
\end{gather}
\end{enumerate}
Here every $B$ in (2) is a cubic Pisot number generating $K$.
For fixed $K,P$, the elements in (2) may be chosen in one coset
$\alpha E^+(M\mathcal O_K)$, where
\[
 M=\prod_{p\in P\cup\{3\}}p,
 \qquad \Nm_{K/\Q}(\alpha)=-N,
 \qquad N\ge1,
 \qquad \gcd(N,M)=1.
\]
Thus their norm is the same integer $-N$. With
$\Delta_M$ denoting the covolume of the logarithmic lattice of
$E^+(M\mathcal O_K)$ in its first two embedding coordinates, the
number of these elements satisfying \eqref{pis:res:eq:strictbounds} and
$B\le H$ is
\begin{equation}
\label{pis:res:eq:count}
 \frac{3}{80\Delta_M}(\log H)^2
       +O_{K,\alpha,M}(\log H+1).
\end{equation}
The entire statement remains true with $\lambda(\beta)=-1$.
\end{theorem}

\begin{proof}
Assume first that $K\hookrightarrow\Q_3$. Then the degree-three
\'{e}tale algebra $K\otimes_\Q\Q_3$ is either
$\Q_3^3$ or $\Q_3\times E$, with $E/\Q_3$ a quadratic field.
The corresponding decomposition of $\mathcal O_K\otimes\Z_3$
allows the following prescriptions modulo three.
\begin{center}\small
\begin{tabular}{lll}
\toprule
Local algebra & Prescribed residues & Residual polynomial\\
\midrule
$\Q_3^3$ & $(1,1,-1)$ & $(X-1)^2(X+1)$\\
$\Q_3\times E$, $E$ ramified & $(-1,1)$ & $(X+1)(X-1)^2$\\
$\Q_3\times E$, $E$ unramified & $(1,i)$, $i^2=-1$ in $\F_9$ & $(X-1)(X^2+1)$\\
\bottomrule
\end{tabular}
\end{center}
In the ramified case the residue $1$ is the scalar element of
$\mathcal O_E/3\mathcal O_E$, so its multiplication operator has
characteristic polynomial $(X-1)^2$, including multiplicity. In the
unramified case $i\notin\F_3$, and its two residual conjugates
are $i$ and $-i$. Every listed element is a unit modulo three.

For each $p\in P\setminus\{3\}$ prescribe the scalar residue
$1\in\mathcal O_K/p\mathcal O_K$. The Chinese remainder theorem
gives a nonempty coset $\alpha_0+M\mathcal O_K$ satisfying all
these prescriptions. Under the three real embeddings,
$M\mathcal O_K$ is a full lattice in $\R^3$. Every translate of
such a lattice meets the open orthant $(+,-,+)$: a sufficiently
large multiple of a vector in that orthant has a lattice-coset
point within a fixed distance. Choose $\alpha$ in this intersection.
Its norm is $-N$ for a positive integer $N$, and all prescribed
residues are units, giving $\gcd(N,M)=1$.

Every element $\beta\in\alpha E^+(M\mathcal O_K)$ has the same
signs, norm, and residues as $\alpha$. Consequently, for every
positive integer $m$ and every $p\in P\setminus\{3\}$,
\[
 \Tr(\beta^m)\equiv3\pmod p.
\]
For every positive odd $m$, the three-adic prescriptions instead
give $\Tr(\beta^m)\equiv1\pmod3$: the residual roots are either
$1,1,-1$ or $1,i,-i$, with multiplicity. Their Teichm\"uller
lifts have sum one, so the classification in
Proposition~\ref{pis:prop:fourteen} gives $\lambda(\beta)=1$ whenever
$\beta$ generates $K$.

It remains to count the elements with the required real conjugates.
The subgroup $E^+(M\mathcal O_K)$ has finite index in the unit
group. By Dirichlet's unit theorem \cite[Chapter~5]{MilneANT},
its first two logarithmic coordinates therefore form a full
rank-two lattice of covolume $\Delta_M$. Positivity removes the
torsion kernel. The corresponding coordinates of
$\alpha E^+(M\mathcal O_K)$ form a translate of this lattice.
Put $x=\log B$, $y=\log\rho$, and $L=\log N\ge0$.
The norm identity gives
\[
 \log s=L-x-y.
\]
For $a=17/40$, the inequalities
$1<B\le H$, $0<s<\rho<1$, and $\rho<B^{-a}$ are exactly
\begin{equation}
\label{pis:res:eq:triangle}
 0<x\le\log H,\qquad \frac{L-x}{2}<y<-ax.
\end{equation}
This region has area
\[
 \frac{1-2a}{4}(\log H)^2+O_N(\log H+1)
\]
and boundary length $O_N(\log H+1)$. Counting fundamental
parallelograms, as in Theorem~\ref{pis:transfer:thm:ray}, gives
\eqref{pis:res:eq:count}. The leading constant is positive.

For every counted point, the dominant embedding exceeds one and
the other two have modulus less than one. Thus $B$ is Pisot and
cannot be rational; because $[K:\Q]=3$, it generates $K$.
The additional bound in \eqref{pis:res:eq:strictbounds} follows from
the integer norm:
\[
 \rho^{40}<B^{-17}\le N^{17}B^{-17}
     =(\rho s)^{17},
 \qquad\text{so}\qquad \rho<s^{17/23}.
\]
Finally, if $m$ is positive and odd, then
$-\rho^m+s^m\in(-1,0)$. Therefore
\[
 \lfloor B^m\rfloor=\Tr_{K/\Q}(\beta^m),
\]
and the preceding residual calculations prove \eqref{pis:res:eq:avoid}.

Conversely, a single element from (2) generates $K$ and has
$\lambda=1$. Proposition~\ref{pis:prop:fourteen} shows that its
minimal polynomial modulo three is one of
\[
 X^2(X-1),\qquad (X-1)(X^2+1),\qquad (X-1)^2(X+1).
\]
Each has a simple linear root over $\F_3$. Hensel's lemma lifts
that root to $\Q_3$, and the resulting root of the irreducible
minimal polynomial defines an embedding $K\hookrightarrow\Q_3$.

For the variant with $\lambda=-1$, negate only the prescribed
residues at three. The real sign choice is still made independently
in the same open orthant. The odd traces are now $-1$ modulo three,
and the Teichm\"uller sum is $-1$; all other arguments are unchanged.
The reverse implication follows by negating the three residual
polynomials in the preceding paragraph.
\end{proof}

\begin{corollary}[Normalized survivors in every real interval]
\label{pis:res:cor:normalized}
Fix an admissible field $K$, a finite prime set $P$, and a nonempty
open interval $I\subset(1,\infty)$. For every sufficiently large
integer $k$, there are elements $\beta$ from
Theorem~\ref{pis:res:thm:equivalence} such that
\[
 B^{1/3^k}\in I,
 \qquad
 \gcd\left(\left\lfloor
    (B^{1/3^k})^{3^{k+n}}\right\rfloor,
    \prod_{p\in P}p\right)=1\quad(n\ge0).
\]
If $I$ is contained in an open interval on which a prescribed finite
list of prime floors $\lfloor t^{3^j}\rfloor$ is constant, these
normalized survivors have that same finite prime prefix.
\end{corollary}
\begin{proof}
Choose $1<a<b$ with $(a,b)\subset I$. Subtracting
\eqref{pis:res:eq:count} at $H=b^{3^k}$ and $H=a^{3^k}$ gives
\[
 \frac{3}{80\Delta_M}3^{2k}
       \bigl((\log b)^2-(\log a)^2\bigr)+O_{K,\alpha,M,a,b}(3^k),
\]
which is positive for all sufficiently large $k$. Any resulting
$a^{3^k}<B\le b^{3^k}$ can be obtained with strict upper endpoint
as well, by choosing a smaller inner interval. The displayed tail
identity is simply
$(B^{1/3^k})^{3^{k+n}}=B^{3^n}$. The assertion about the finite
prefix follows from membership in $I$.
\end{proof}

\begin{remark}[Scope of the obstruction]
\label{pis:res:rem:scope}
The quantifiers in Theorem~\ref{pis:res:thm:equivalence} are
\[
 \text{for every finite }P\text{, there exists a fixed }N
 \text{ and infinitely many }\beta.
\]
The integer $N$ and the coset may depend on $P$; no one element
is asserted to avoid every prime. The normalized construction
likewise asserts neither an all-prime tail nor a prime condition
at the intervening indices between a fixed prefix and its prescribed
tail. What is ruled out is a field-uniform proof using only
Saito's strict conjugate bounds, the limiting trace condition, and
divisibility by one fixed finite set of primes. The theorem does not
rule out adaptive choices of divisors depending on $\beta$, other
arithmetic information, or an argument using the least-prime property
specific to the minimum defining Mills' constant.
\end{remark}


\section{The remaining Mills alternative}
The results yield several rigorous forms of progress, but each has
a precise quantifier. The simultaneous-window theorem closes a
residue-defined subclass. The coefficient sieve controls the size
of a possible exceptional set in boxes. The Pell certificates
exclude entire arithmetic classes of parameters in one family.
The conditional construction realizes every prescribed \emph{finite}
pattern, and the unit-lattice theorem preserves \emph{fixed finite}
local data. The new residual-field theorem supplies its own seed
in every three-adically admissible cubic field; this still preserves
a finite-divisor quantifier, not primality. None of these statements exchanges a finite quantifier
for an infinite one.

Three specific obstructions remain.
\begin{enumerate}
\item At a nonintegral three-adic limit, the prime-factor order
budget works for every fixed shift. At $\lambda=\pm1$, its key
distance from $\pm1$ is zero for the unshifted trace. These are
exactly the residual classes in which a new argument is required.
\item A coefficient set of stretched-exponentially small proportion
can still contain a particular algebraic candidate. An ambient
box estimate gives no automatic exclusion on a fixed field or on
the exponentially sparse Pell orbit. The separate Pell calculation
is needed precisely for this reason.
\item An infinite collection of algebraic numbers, each with a
long finite prime prefix, need not contain one number with an
infinite prime tower. Nor does it determine the minimum defining
Mills' constant. The minimum requires a global prime-chain argument
in addition to local polynomial information.
\end{enumerate}
Consequently the remaining full targets are unchanged: Saito's
cubic Pisot Problem~1.1, the exclusion of the Mills cubic alternative,
and irrationality of the least square-clock constant. What is
established here is a stronger collection of restricted theorems, quantitative bounds, and exact exclusions, with reusable proofs.

\paragraph{A concrete remaining target using minimality.}
For a surviving cubic candidate let $T_n=\Tr(\beta^{3^n})$.
A hypothetical algebraic least Mills constant requires both eventual
primality of $T_n$ and, for every sufficiently large $n$,
\[
 T_{n+1}=\operatorname{nextprime}(T_n^3).
\]
The latter property is an inherited observation from the supplied
manuscript [010]. To see why it holds, the prime-interval bound used
in \cite{SaitoVariants} extends every sufficiently large prime $q$
to a prime child strictly between $q^3$ and $(q+1)^3-1$.
An earlier prime child of $T_n$ could therefore be continued to an
infinite nested prime chain below the defining minimum, a contradiction.
Consequently it would suffice to prove the following restricted,
still unproved assertion: every surviving candidate either has
infinitely many composite $T_n$, or admits arbitrarily large $n$
with a prime strictly between $T_n^3$ and $T_{n+1}$.
This target uses the least-prime selection rule that the fixed-prime
survival theorem does not encode. It would settle Mills transcendence
without requiring composite floors for all cubic Pisot numbers.

\paragraph{Verification boundary.}
The present edition replays the complete Python certificate below and
adds independent checks of the new structural extensions. The following
additional audit description is inherited from the source manuscript:
its central symbolic arguments were checked in parallel. The residual-field proofs were checked independently,
including the ramified local algebra at three, all-rational-power
quantifiers, and the distinction between a finite-prime sieve and an
all-prime tower. Exact arithmetic replayed the two finite-field
factorizations of the balanced seed, the old 26 second-orbit Pell
witnesses and five first-orbit witnesses, the three new witnesses,
and the completeness of the new residue cover. Additional exhaustive
small-prime checks test the shifted local counts; they are consistency
checks, not substitutes for the all-degree proof. No large numerical
prime certificate from an earlier submission is needed for the new
theorems. A compact standard-library certificate is included below.

\section{Prime clocks in arbitrary degree and polynomial observables}
\label{newclock:sec}

The simultaneous-return argument extends beyond cubic traces.  The
degree enters only through a finite-field order budget.  This section
is independent of the phase-embedding hypotheses used in the
Euler--Bessel part of the paper.

Let $\ell$ be a prime, let $d\ge1$, and let
$\alpha_1,\ldots,\alpha_d$ be the roots, counted with multiplicity,
of a monic polynomial in $\mathbb Z[X]$.  Define
\[
 \boldsymbol a_n=(a_{1,n},\ldots,a_{d,n}),\qquad
 a_{j,n}=e_j(\alpha_1^{\ell^n},\ldots,\alpha_d^{\ell^n}),
\]
where $e_j$ is the $j$th elementary symmetric polynomial.  All
$a_{j,n}$ are integers.  Put
\[
 r_\ell=\begin{cases}2,&\ell=2,\\ \ell-1,&\ell>2,\end{cases}
 \qquad b_{\ell,d}=\lfloor\log_\ell d\rfloor.
\]

\begin{lemma}[Coefficient contraction at every prime clock]
\label{newclock:contraction}
There is an integral polynomial map $F_\ell:\mathbb Z^d\to
\mathbb Z^d$ such that $\boldsymbol a_{n+1}=F_\ell(\boldsymbol
a_n)$ and, as a polynomial identity modulo $\ell$,
\[
 (F_\ell)_j(X_1,\ldots,X_d)\equiv X_j^\ell\pmod\ell.
\]
The sequence $\boldsymbol a_n$ converges to a vector
$\boldsymbol\Lambda\in\mathbb Z_\ell^d$, and
\[
 \boldsymbol\Lambda-\boldsymbol a_n
     \in\ell^{n+1}\mathbb Z_\ell^d.
\]
The limit depends only on $\boldsymbol a_0\pmod\ell$.
\end{lemma}

\begin{proof}
The fundamental theorem of symmetric polynomials supplies the
integral map.  In characteristic $\ell$, the identity
$e_j(T_1^\ell,\ldots,T_d^\ell)=e_j(T_1,\ldots,T_d)^\ell$
proves the displayed polynomial congruence.  Consequently every
first partial derivative of $(F_\ell)_j$ is divisible by $\ell$.
If $\boldsymbol x\equiv\boldsymbol y\pmod{\ell^s}$ with $s\ge1$,
the linear terms of the integral Taylor expansion are divisible
by $\ell^{s+1}$ and all higher terms by $\ell^{2s}$.  Thus
\[
 F_\ell(\boldsymbol x)\equiv F_\ell(\boldsymbol y)
          \pmod{\ell^{s+1}}.
\]
Fermat's congruence gives $F_\ell(\boldsymbol a_0)\equiv
\boldsymbol a_0\pmod\ell$.  Iteration yields
$\boldsymbol a_{n+1}-\boldsymbol a_n\in
\ell^{n+1}\mathbb Z^d$ and the asserted limit and error bound.
Applying the same contraction to two initial vectors congruent
modulo $\ell$ proves the last assertion.
\end{proof}

\begin{lemma}[A uniform residual-order bound]
\label{newclock:budget}
Let $q\ne\ell$ be a rational prime.  Then
\[
 v_\ell\!\left(\operatorname{lcm}_{1\le f\le d}(q^f-1)\right)
 \le v_\ell(q^{r_\ell}-1)+b_{\ell,d}.
\]
\end{lemma}
\begin{proof}
For odd $\ell$, put $h=\operatorname{ord}_{\mathbb F_\ell^\times}(q)$.
If $h\nmid f$, the relevant valuation is zero.  Otherwise the
lifting-the-exponent identity gives
\[
 v_\ell(q^f-1)=v_\ell(q^h-1)+v_\ell(f/h)
              =v_\ell(q^{\ell-1}-1)+v_\ell(f/h),
\]
which has the required bound.  For $\ell=2$, the prime $q$ is
odd.  If $f$ is odd then $v_2(q^f-1)=v_2(q-1)\le v_2(q^2-1)$.
If $f$ is even, the same elementary identity gives
$v_2(q^f-1)=v_2(q^2-1)+v_2(f)-1$.  Taking the maximum over
$1\le f\le d$ proves the assertion, with a harmless slack in
the dyadic bound.
\end{proof}

\begin{theorem}[Simultaneous composite polynomial observables]
\label{newclock:polynomial}
Let $P_1,\ldots,P_s\in\mathbb Z[X_1,\ldots,X_d]$, and suppose
\[
 Q_j(n):=P_j(\boldsymbol a_n)\longrightarrow+\infty
       \qquad(1\le j\le s).
\]
Write $\lambda_j=P_j(\boldsymbol\Lambda)$.  If
\[
 \lambda_j^{r_\ell}\ne1\qquad(1\le j\le s),
\]
there are integers $n_0\ge0$ and $D\ge1$ for which every
$Q_j(n_0+mD)$ is composite for all sufficiently large integers
$m$.  In particular the simultaneous composite indices have
positive lower density.
\end{theorem}

\begin{proof}
If $\ell\mid\lambda_j$, then $\ell\mid Q_j(n)$ at every index,
because the coefficient vector is constant modulo $\ell$.
This supplies a fixed recurring prime for that observable.

For the remaining indices $j$, define the finite integer
$K_j=v_\ell(\lambda_j^{r_\ell}-1)$.  Choose one sufficiently
large $n_0$ so that all $Q_j(n_0)>1$, so that
\[
 v_\ell\bigl(Q_j(n_0)^{r_\ell}-1\bigr)=K_j,
 \qquad n_0\ge K_j+b_{\ell,d}
\]
for every unit limit, and so that the limits have stabilized to
these valuations.  Such a choice follows from
Lemma~\ref{newclock:contraction} and the assumed real growth.

Fix one unit limit.  Some prime divisor $q_j$ of $Q_j(n_0)$
satisfies
\[
 v_\ell(q_j^{r_\ell}-1)\le K_j.
\]
Indeed, if every prime divisor had valuation at least $K_j+1$,
raising the prime factorization of $Q_j(n_0)$ to the power
$r_\ell$ would give $Q_j(n_0)^{r_\ell}\equiv1
\pmod{\ell^{K_j+1}}$, a contradiction.  Here $q_j\ne\ell$
because $Q_j(n_0)$ is an $\ell$-adic unit.

Reduce the original degree-$d$ polynomial modulo $q_j$.
Every nonzero residual root belongs to $\mathbb F_{q_j^f}$
for some $f\le d$.  By Lemma~\ref{newclock:budget}, the
$\ell$-primary part of its multiplicative order divides
$\ell^{n_0}$.  Its $\ell^{n_0}$th power therefore has order
prime to $\ell$.  There is a common $D_j\ge1$ such that raising
all these powered roots to $\ell^{D_j}$ fixes each of them.
Zero roots remain zero, and repeated roots retain their
multiplicities.  Hence the entire coefficient vector satisfies
\[
 \boldsymbol a_{n_0+mD_j}\equiv\boldsymbol a_{n_0}
                  \pmod{q_j}\qquad(m\ge0),
\]
and $q_j\mid Q_j(n_0+mD_j)$.

Take $D$ to be the least common multiple of the finitely many
$D_j$, using $D=1$ if all observables have $\ell$-divisible
limits.  All selected primes recur on this common progression.
Since every $Q_j(n)$ tends to $+\infty$, each selected divisor
is eventually proper.  The progression has density $1/D$ in
the nonnegative indices.
\end{proof}

\begin{corollary}[Pisot floor windows in every degree]
\label{newclock:floors}
Let $\beta$ be any Pisot number, of any degree, and let
$T_n=\operatorname{Tr}_{\mathbb Q(\beta)/\mathbb Q}
(\beta^{\ell^n})$.  Write $\lambda=\lim_n T_n\in\mathbb Z_\ell$.
For an integer $H\ge0$, assume
\[
 (\lambda+h)^{r_\ell}\ne1
       \qquad(-H-1\le h\le H).
\]
There is an arithmetic progression of indices on which,
eventually, every integer
\[
 \lfloor\beta^{\ell^n}\rfloor+h,\qquad |h|\le H,
\]
is composite.  For $\ell=2$ or $\ell=3$, a single sufficient
condition for this conclusion at every fixed $H$ is
$\lambda\notin\mathbb Z$.
\end{corollary}
\begin{proof}
Apply Theorem~\ref{newclock:polynomial} to the linear observables
$P_h(\boldsymbol X)=X_1+h$ for $-H-1\le h\le H$.
The Pisot property gives $T_n-\beta^{\ell^n}\to0$ in
$\mathbb R$, and $T_n\to+\infty$.  Thus the floor belongs
to $\{T_n-1,T_n\}$ for every sufficiently large $n$,
regardless of the sign of the error.  Both possible floor
windows are contained in the selected trace window.  If
$\ell=2$ or $3$, then $r_\ell=2$, so the equation $x^{r_\ell}=1$
has only the two solutions $1$ and $-1$; equality at an
integer shift would force $\lambda\in\mathbb Z$.
\end{proof}

\begin{example}[A quartic window theorem]
\label{newclock:tetranacci}
Let $\beta$ be the real root greater than one of
\[
 f(X)=X^4-X^3-X^2-X-1.
\]
For every fixed $H$, all integers in the radius-$H$ window
around $\lfloor\beta^{3^n}\rfloor$ are simultaneously
composite on an eventual arithmetic progression of indices.

Here are complete verifications of the hypotheses.  Since
$f(1)<0<f(2)$, a root lies in $(1,2)$.  The identity
$(X-1)f(X)=X^5-2X^4+1$ and Rouch\'e's theorem show that
for every $r>1$ sufficiently close to one, the latter polynomial
has exactly four roots in $|X|<r$: on that circle,
$r^5+1<2r^4$.  One of these roots is $1$, so $f$ has
three roots in $|X|<r$.  It has no root on $|X|=1$,
since equality in $X^5+1=2X^4$ would force both
$X^5=1$ and $X^4=1$, hence $X=1$, while $f(1)\ne0$.
Its other three roots therefore lie strictly in the unit
disc.  A nontrivial monic integral factor without the root
$\beta$ would have a nonzero integral constant term of
absolute value less than one.  Thus $f$ is irreducible and
$\beta$ is a quartic Pisot unit.

Modulo three,
\[
 f(X)=(X-1)(X^3-X+1),
\]
and the cubic factor has no root in $\mathbb F_3$.
The three-adic trace limit is the sum of the Teichm\"uller
lifts of the four residual roots.  To see this directly,
work in an unramified splitting field: each unit root is
its Teichm\"uller lift times a principal unit, whose
$3^n$th power tends to one; Frobenius permutes the
multiset of Teichm\"uller lifts.  The linear factor
contributes $1$.  The cubic factor contributes
\[
 \varepsilon(\zeta^a+\zeta^{3a}+\zeta^{9a}),
 \qquad\varepsilon\in\{1,-1\},\quad
 a\not\equiv0\pmod{13},
\]
where $\zeta$ is a primitive thirteenth root of unity in
the unramified cubic extension of $\mathbb Q_3$.
This sum is not rational.  Otherwise a polynomial of
degree at most twelve, with at most four nonzero
coefficients, would vanish at $\zeta$, contradicting
its minimal polynomial $1+X+\cdots+X^{12}$ over
$\mathbb Q$.  Thus $\lambda\notin\mathbb Z$, and
Corollary~\ref{newclock:floors} applies.
\end{example}

\begin{proposition}[An explicit quartic progression]
\label{newclock:tetranacci-explicit}
For the quartic Pisot number in Example~\ref{newclock:tetranacci},
the three integers
\[
 \lfloor\beta^{3^n}\rfloor-1,\qquad
 \lfloor\beta^{3^n}\rfloor,\qquad
 \lfloor\beta^{3^n}\rfloor+1
\]
are all composite for every sufficiently large
$n\equiv4\pmod{68}$.
\end{proposition}
\begin{proof}
Write $S_m=\sum_{i=1}^4\alpha_i^m$ and $T_n=S_{3^n}$.
Newton's identities give
\[
 S_0=4,\quad S_1=1,\quad S_2=3,\quad S_3=7,\qquad
 S_{m+4}=S_{m+3}+S_{m+2}+S_{m+1}+S_m,
\]
and hence the exact value
$T_4=S_{81}=121808162096862030748321$.
The following table records recurring divisors of four consecutive
trace shifts; the period refers to the index $n$ after its initial
transient.
\begin{center}
\begin{tabular}{rrrr}
\toprule
Shift $h$ & Prime $q$ & $T_4\bmod q$ & Period\\
\midrule
$-2$ & $17$ & $2$ & $68$\\
$-1$ & $3$ & $1$ & $1$\\
$0$ & $11$ & $0$ & $4$\\
$1$ & $2$ & $1$ & $4$\\
\bottomrule
\end{tabular}
\end{center}
Here is a finite verification of the periods. Modulo two,
$f(X)=\Phi_5(X)$, so the nonzero root orders divide five
and $3^4\equiv1\pmod5$. Modulo eleven,
\[
 f(X)=(X^2+X-3)(X^2-2X+4).
\]
The roots lie in $\mathbb F_{11^2}^\times$; after one cubing
their orders divide forty, and $3^4\equiv1\pmod{40}$.
Modulo seventeen,
\[
 f(X)=(X-6)(X^3+5X^2-5X+3),
\]
where the cubic has no root in $\mathbb F_{17}$ and is
irreducible.  The root orders thus divide $17^3-1=4912$,
and $3^{68}\equiv1\pmod{4912}$.  Finally the Frobenius
identity gives $T_n\equiv S_1=1\pmod3$ for every $n$.
All four listed primes therefore divide their indicated
shifts when $n\equiv4\pmod{68}$.  The shifts grow to
infinity, making the divisors proper, and the Pisot floor
belongs to $\{T_n-1,T_n\}$ eventually.  Its radius-one
window is consequently contained in these four composite
trace shifts.
\end{proof}

\begin{remark}[What the extension does and does not exclude]
The theorem enlarges the degree, the prime clock, and the
class of observables without assuming prime initial values.
Its excluded roots of unity are genuine limitations of
this argument: for odd $\ell$, an $\ell$-adic unit may
satisfy $\lambda^{\ell-1}=1$ without being the rational
number $1$ or $-1$.  Therefore the noninteger-limit
simplification in Corollary~\ref{newclock:floors} is
asserted only for clocks two and three.  This extension
does not exclude the residual cubic alternatives with
trace limit $1$ or $-1$, and hence does not settle the
full cubic Pisot question or Mills' transcendence.
No external priority claim for the general extension is
made here.
\end{remark}
\section{The predecessor branch survives every fixed modulus}
\label{new:sec:predecessor}
The three-adic embedding criterion for a negative dominant small
conjugate does not extend unchanged to the other floor branch.  In
fact the predecessor branch admits the following uniform construction
in every totally real cubic field.  The modulus may contain arbitrary
prime powers, and the conclusion holds at every positive integer
exponent, rather than only at cubic-power exponents.

\begin{theorem}[Constant invertible floor residues in every field]
\label{new:thm:predecessor}
Let $K$ be a totally real cubic field with ordered embeddings
$\sigma_1,\sigma_2,\sigma_3$, let $M\ge2$, and fix $0<a<1/2$.
There exist a positive integer $C$, a positive integer $M'$ divisible
by $M$, and a residue $c_M\in(\mathbb Z/M\mathbb Z)^\times$
such that infinitely many
\[
 \beta\in C E^+(M'\mathcal O_K)
\]
have embeddings $B=\sigma_1(\beta)$, $r=\sigma_2(\beta)$,
$s=\sigma_3(\beta)$ satisfying
\[
 B>1,\qquad 0<s<r<B^{-a},\qquad r+s<1.
\]
Each such $B$ is a cubic Pisot number generating $K$, has norm $C^3$
and three-adic trace limit zero, and satisfies
\begin{equation}\label{new:eq:allpowers}
 \lfloor B^m\rfloor=\operatorname{Tr}_{K/\mathbb Q}(\beta^m)-1
       \equiv c_M\pmod M\qquad(m\ge1).
\end{equation}
The residue can be chosen to be $-1$ modulo the $2$- and $3$-primary
parts of $M$, and $2$ modulo every other prime-power factor of $M$.
If $\Delta_{M'}$ is the two-coordinate logarithmic covolume of
$E^+(M'\mathcal O_K)$, the number of elements in this coset with
$B\le H$ and the displayed inequalities is
\[
 \frac{1-2a}{4\Delta_{M'}}(\log H)^2+O_{K,C,M',a}(\log H+1).
\]
\end{theorem}
\begin{proof}
Put $M'=\operatorname{lcm}(M,6)$. By the Chinese remainder theorem,
choose $C>0$ such that $C\equiv0$ modulo the full $2$- and $3$-primary
parts of $M'$, and $C\equiv1$ modulo every other prime-power factor
of $M'$. Every element of $C E^+(M'\mathcal O_K)$ is totally positive,
has norm $C^3$, and is congruent to the scalar $C$ modulo
$M'\mathcal O_K$. Consequently, for every $m\ge1$,
\[
 \operatorname{Tr}(\beta^m)-1\equiv
 \begin{cases}
 -1\pmod{p^e},&p\in\{2,3\},\ p^e\mid M',\\
 2\pmod{p^e},&p\ge5,\ p^e\mid M'.
 \end{cases}
\]
All these residues are units. Since $3\mid C$, the residual
characteristic polynomial at three is $X^3$, and the three-adic trace
limit is zero.

Write $x=\log B$, $y=\log r$, and $L=3\log C$. The third embedding
has logarithm $L-x-y$. Apart from the extra condition $r+s<1$, the
required region is
\[
 0<x\le\log H,\qquad (L-x)/2<y<-ax.
\]
Its area is $(1-2a)(\log H)^2/4+O_{C,a}(\log H+1)$ and its boundary
length is $O_{C,a}(\log H+1)$. Its logarithmic coset is a translate
of a rank-two lattice of covolume $\Delta_{M'}$. Counting fundamental
parallelograms gives the stated asymptotic. Moreover $r+s<2B^{-a}<1$
once $B>2^{1/a}$, so imposing this extra condition removes only a
bounded portion of the region and does not alter the asymptotic.

The resulting dominant embedding is greater than one, and its two
other conjugates lie in $(0,1)$. It cannot be rational, and hence has
degree three. Finally $0<r^m+s^m\le r+s<1$ for every $m\ge1$.
The integer trace is $B^m+r^m+s^m$, which proves both the floor
identity and its congruence in \eqref{new:eq:allpowers}.
\end{proof}

\begin{corollary}[No field-uniform finite divisor cover in this branch]
For every totally real cubic field $K$ and every finite set $P$ of
rational primes, infinitely many cubic Pisot numbers in $K$ have
$\gcd(\lfloor B^m\rfloor,\prod_{p\in P}p)=1$ for every $m\ge1$.
This remains true in fields having no embedding into $\mathbb Q_3$.
\end{corollary}
\begin{proof}
Take $M$ divisible by every member of $P$ and apply
Theorem~\ref{new:thm:predecessor}.
\end{proof}
The norm is fixed within each construction but may depend on $M$;
it is divisible by three. Thus the corollary does not contradict the
prime-to-three norm hypothesis in Theorem~\ref{pis:local:thm:noembedding}.
It also does not assert primality of any value: the surviving floors
may have prime divisors outside the prescribed finite set. In
particular it neither disproves Saito's problem nor eliminates its
remaining cases.

\subsection{An eventual-compositeness refinement of the parity criterion}
\begin{proposition}\label{new:prop:eventual-composite}
Let $\beta$ be a cubic Pisot number with minimal polynomial
$X^3-uX^2+vX-N$, and suppose $3\mid u$. If either
\[
 N\equiv1\pmod2,\quad u\equiv v\pmod2,
 \qquad\text{or}\qquad
 N\equiv0,\quad u\equiv1,\quad v\equiv0\pmod2,
\]
then $\lfloor\beta^{3^n}\rfloor$ is composite for every sufficiently
large $n$. Nonreal conjugates are allowed. If instead $N$ is odd
and $u\not\equiv v\pmod2$, the same argument guarantees composite
floors on one parity class of indices, and hence lower density at
least $1/2$.
\end{proposition}
\begin{proof}
The trace $U_n$ remains zero modulo three. With odd norm the
coefficient pair modulo two follows
$00\mapsto11\mapsto11$ and $01\leftrightarrow10$; thus equal
initial parities make $U_n$ odd eventually, and opposite initial
parities make it odd on one parity class. With even norm the map
is $(a,b)\mapsto(a+ab,b)$, so the state $(1,0)$ is fixed. In each
case $U_n$ is odd at the asserted indices. The Pisot estimate puts
the sufficiently late floor in $\{U_n,U_n-1\}$. At those indices
$U_n$ is divisible by three and $U_n-1$ is divisible by two; both
are eventually larger than these divisors.
\end{proof}

\section{Conclusions and remaining exact targets}
The composite-window theorem now applies in every degree and for every
prime clock, subject to its explicit local-limit exclusion. The finite
winding theorem replaces infinitely many exact specializations by a
degree-dependent finite test in the mixed two-block setting, under the
same independence hypotheses as its intersection theorem. The derivative
determinant is obtained by a direct differential-algebra calculation.
None of these statements removes the surviving $\lambda=\pm1$ cubic
Mills branch, supplies the missing joint-independence proof, or identifies
the canonical winding parameter arithmetically. Those remain separate
requirements for stronger conclusions.

\appendix

\section{The inherited Pell reduction and the exact finite cover}
For self-contained use of the new certificate we record the earlier
reduction and cover. The substitutions $x=8u+5$, $y=4w$ transform
$4u^2+5u+23=2w^2$ into
\[
 x^2-2y^2=-343,\qquad x\equiv5\pmod8,\quad y\equiv0\pmod4.
\]
Multiplying $x+y\sqrt2$ by powers of $3-2\sqrt2$ reduces it uniquely
to $[\sqrt{343},(3+2\sqrt2)\sqrt{343})$. This bounds
$0<x<53$, $0<y<40$. Enumeration gives exactly
$(7,14),(13,16),(25,22),(35,28)$. Their four-term residue cycles
modulo eight retain only the first seed at odd powers of
$3+2\sqrt2$, and the second at even powers. These become the two
orbits from $(9,14)$ and $(1,4)$ under $T$. The latter initial point
is omitted because $u=1$ is a reducible degeneration.

For the first orbit $a_j$, the residues have periods two, ten, and
five modulo $3,19,29$ respectively:
\[
\begin{array}{c|l}
3&0,1\\
19&9,8,17,1,18,4,5,15,12,14\\
29&9,12,13,15,24.
\end{array}
\]
Even indices give trace zero modulo three. The five odd classes
$j=1,3,5,7,9\pmod{10}$ are covered by
$(p,c,n)=(29,12,41),(29,15,5),(19,4,17),(29,13,1),(19,14,7)$.
Every listed zero is beyond its root-order return budget.
For the second orbit odd $j$ give trace zero modulo three. Writing
$j=2k$, the following old rows cover all $k$ except
$0,2187\pmod{6561}$; a pair $(a,n)$ means a recurring zero at
trace index $n$ whenever $k\equiv a\pmod m$.
\begin{center}\small
\begin{tabular}{r r l}
\toprule
$p$&$m$&$(a,n)$\\\midrule
5&3&$(1,3)$\\
7&3&$(2,8)$\\
179&9&$(3,52)$\\
199&9&$(6,126)$\\
107&27&$(9,4),(18,25)$\\
487&81&$(27,366),(54,1333)$\\
4373&729&$(162,2379),(486,7323),(567,3228)$\\
20411&729&$(81,7662),(405,34139)$\\
8747&2187&$(243,7161),(324,2814),(648,105),(729,830)$\\
&&$(1053,2330),(1701,3456),(1782,4165),(2106,12465)$\\
52489&2187&$(1377,54889)$\\
39367&6561&$(3645,45679),(4374,15588)$\\
629857&2187&$(972,647491),(1458,678551)$\\
\bottomrule
\end{tabular}
\end{center}
The three new rows of Theorem~\ref{pis:thm:pell-new} reduce the survivor
set to $k\equiv0\pmod{19683}$.

To verify the valuation formula used there, set
$(17+12\sqrt2)^{2k}=C_k+D_k\sqrt2$. Then
\[
 b_{2k}-1=\frac{13(C_k-1)+32D_k}{8},\quad
 D_{3k}=D_k(3+8D_k^2).
\]
For $3\nmid k$, reduction modulo nine gives $v_3(D_k)=1$.
Repeated cubing gives $v_3(D_k)=1+v_3(k)$ for every $k>0$.
Since $C_k^2-2D_k^2=1$ and $C_k\equiv1\pmod3$,
$v_3(C_k-1)=2v_3(D_k)$. The term involving $D_k$ therefore
determines the numerator's valuation, proving
$v_3(b_j-1)=1+v_3(j)$ for even $j>0$.

\section{A reproducible exact certificate}
The following Python 3 code uses only the standard library. It
checks primality by trial division, modular periods in the quadratic
ring, cubic trace zeros by binary polynomial powering, and the full
residue cover. It never constructs a trace of millions-digit exponent.
The algebraic reasons why these finite checks imply infinite
progressions are given in the proofs above.

\begin{lstlisting}[language=Python]
from math import isqrt

rows=[(5,3,[(1,3)]),(7,3,[(2,8)]),(179,9,[(3,52)]),(199,9,[(6,126)]),(107,27,[(9,4),(18,25)]),(487,81,[(27,366),(54,1333)]),(4373,729,[(162,2379),(486,7323),(567,3228)]),(20411,729,[(81,7662),(405,34139)]),(8747,2187,[(243,7161),(324,2814),(648,105),(729,830),(1053,2330),(1701,3456),(1782,4165),(2106,12465)]),(52489,2187,[(1377,54889)]),(39367,6561,[(3645,45679),(4374,15588)]),(629857,2187,[(972,647491),(1458,678551)])]
def v3(n):
 a=0
 while n%3==0 and n: a+=1;n//=3
 return a

def qm(x,y,p):return ((x[0]*y[0]+2*x[1]*y[1])%p,(x[0]*y[1]+x[1]*y[0])%p)
def qp(x,n,p):
 r=(1,0)
 while n:
  if n&1:r=qm(r,x,p)
  x=qm(x,x,p);n//=2
 return r

def b(j,p):
 x,y=qm((13,16),qp((17,12),j,p),p)
 return ((x-5)*pow(8,-1,p))%p

def pm(a,b,u,p):
 c=[0]*5
 for i,x in enumerate(a):
  for j,y in enumerate(b):c[i+j]=(c[i+j]+x*y)%p
 for k in (4,3):
  v=c[k];c[k]=0;c[k-1]=(c[k-1]+u*v)%p;c[k-2]=(c[k-2]+v)%p;c[k-3]=(c[k-3]-v)%p
 return c[:3]
def pp(a,n,u,p):
 r=[1,0,0]
 while n:
  if n&1:r=pm(r,a,u,p)
  a=pm(a,a,u,p);n//=2
 return r

def trace(n,u,p):
 M=(p-1)*(p*p-1)*(p**3-1); E=pow(3,n,M); A,B,C=pp([0,1,0],E,u,p)
 return (3*A+u*B+(u*u+2)*C)%p
for p,m,items in rows:
 assert p>=2 and all(p%d for d in range(2,isqrt(p)+1)),p
 assert qp((17,12),2*m,p)==(1,0),(p,m)
 for a,n in items:
  assert n>=v3((p-1)*(p*p-1)*(p**3-1)),(p,n)
  assert trace(n,b(2*a,p),p)==0,(p,m,a,n,b(2*a,p))
surv=[k for k in range(6561) if all(all(k%m!=a for a,n in items) for p,m,items in rows)]
assert surv==[0,2187],surv
for p,c,n in [(29,12,41),(29,15,5),(19,4,17),(29,13,1),(19,14,7)]:
 assert trace(n,c,p)==0
 assert n>=v3((p-1)*(p*p-1)*(p**3-1))
print('All 26 second-orbit witnesses and 5 first-orbit witnesses verified by exact arithmetic.')
print('Surviving k residues modulo6561:',surv)

# Three additional witnesses: (prime, parameter period, class, trace index).
new=[(8279983,6561,2187,4863973),
     (27398737,19683,6561,4929815),
     (15037811,19683,13122,549606)]
for p,m,a,n in new:
 assert all(p%d for d in range(2,isqrt(p)+1))
 assert qp((17,12),2*m,p)==(1,0)
 assert n>=v3((p-1)*(p*p-1)*(p**3-1))
 assert trace(n,b(2*a,p),p)==0
remaining=[k for k in range(19683)
 if all(all(k%m!=a for a,n in items) for p,m,items in rows)
 and all(k%m!=a for p,m,a,n in new)]
assert remaining==[0]
# Reduced Pell seeds and first-orbit parameter periods.
seeds=[(x,y) for x in range(1,53) for y in range(1,40)
       if x*x-2*y*y==-343]
assert seeds==[(7,14),(13,16),(25,22),(35,28)]
def orbit(seed,p,length):
 u,w=seed;out=[]
 for j in range(length):
  out.append(u%p);u,w=(17*u+12*w+10)%p,(24*u+17*w+15)%p
 assert (u,w)==(seed[0]%p,seed[1]%p)
 return out
assert orbit((9,14),3,2)==[0,1]
assert orbit((9,14),19,10)==[9,8,17,1,18,4,5,15,12,14]
assert orbit((9,14),29,5)==[9,12,13,15,24]
assert orbit((1,4),3,2)==[1,0]
print('All certificates verified; surviving k modulo19683:',remaining)

\end{lstlisting}

\section{Source map and attribution}
The earlier archive cited by the source manuscript is research material, not a collection
of independently published papers. Round identifiers below refer to
\texttt{Race244\_Topic02(1).zip}; S means Showa. The earlier numbered
manuscripts refer to the previously supplied 36-paper collection.
\begin{center}\small
\begin{tabular}{p{.17\textwidth}p{.72\textwidth}}
\toprule
Source & Use and status\\\midrule
S12, S27 & Local factor blocks and uniform irreducibility, reproduced and audited.\\
S28 & Finite prime/composite patterns under H; exact multiplicities are added here.\\
S35, S36 & Complete Pell reduction and inherited finite cover; three new rows added.\\
S37 & Fixed-norm trace-zero counts and unshifted quantitative sieve; extended to all shifts and odd degrees.\\
S14 & Inherited all-norm cyclic prime-trace restrictions; used in the positive-density floor synthesis.\\
S20 & Prior norm-one trace theorem in discriminant $1524$; used only in the explicitly identified field corollary.\\
Heisei R05 & A cyclic cubic seed for finite-sieve survival; the coset-counting theorem is proved here.\\
Old [007], [008] & Previously resolved dyadic Fibonacci and nonreversible recurrence problems.\\
Old [009], [031] & Earlier local obstructions, the three-adic embedding lemma, nonreal case, qualitative density and dimension results.\\
\bottomrule
\end{tabular}
\end{center}
The source manuscript reports that other topics were screened for useful connections.
The pure-cubic sixth-power observation recovers results already
subsumed by the old nonreal-cubic argument. No proved identity in
the Gieseking, integral-period, or elliptic-resolvent notes supplies
a missing prime-divisor theorem for a surviving Mills candidate.
No such speculative connection is used as a proof input.


\begin{thebibliography}{99}
\bibitem{EB} Keio, \emph{Euler--Gompertz and Bessel Values: Joint
Independence, Integral Intersections, and Critical Limits}, revised
integrated research manuscript, 22 October 2026. Named labels refer to
its current \LaTeX\ source. Not included in the three attachments for this edition; its joint-independence claims are explicit inputs here.

\bibitem{FR} S. Fischler and T. Rivoal, \emph{The ring of values of
$E$-functions is not a field}, preprint, 2026, arXiv:2610.16603.
\url{https://arxiv.org/abs/2610.16603}.

\bibitem{FRG} S.~Fischler and T.~Rivoal, \emph{Relations between values of
arithmetic Gevrey series, and applications to values of the Gamma function},
arXiv:2301.13518.

\bibitem{MZ} T.~Matala-aho and W.~Zudilin, \emph{Euler's factorial series and
global relations}, J. Number Theory \textbf{186} (2018), 202--210.

\bibitem{MS} M.~R.~Murty and N.~Saradha, \emph{Euler--Lehmer constants and a
conjecture of Erd\H{o}s}, J. Number Theory (2010),
doi:10.1016/j.jnt.2010.07.004.

\bibitem{WL} Keio, \emph{The winding line: continuous phase embeddings
and two harmonic Bessel coordinates at dependent scales}, supplied research
manuscript, October 2026; incorporated into Part I of the present edition.

\bibitem{BR} Keio, \emph{A monic resultant for Bessel derivative
coefficients}, research note, 22 October 2026, 5 pp. Not included in the three attachments; its all-derivative nonvanishing assertion is an explicit input where invoked.

\bibitem{BakerIII} A.~Baker, \emph{Linear forms in the logarithms of
algebraic numbers (III)}, Mathematika \textbf{14} (1967), 220--228.
\href{https://doi.org/10.1112/S0025579300003843}{doi:10.1112/S0025579300003843}.

\bibitem{DLMF} NIST Digital Library of Mathematical Functions,
Chapters 6 and 10, especially Sections 6.2, 10.29, 10.31, 10.32,
10.34, and 10.40. \url{https://dlmf.nist.gov/}.

\bibitem{MilneANT} J. S. Milne, \emph{Algebraic Number Theory},
version 3.08 (2020), Chapters 5--8.
\url{https://www.jmilne.org/math/CourseNotes/ANT.pdf}.

\bibitem{S14} Showa, \emph{Conductor and local restrictions on cyclic
cubic prime traces}, submitted constant-race round 14,
20 October 2026. Internal supplied research manuscript.

\bibitem{SaitoIntervals} K. Saito,
\emph{Intervals without primes near an iterated linear recurrence sequence},
International Journal of Number Theory \textbf{22} (2026), 1269--1279.
Problem numbering used here: arXiv:2504.14968v1.
\url{https://arxiv.org/html/2504.14968v1}.

\bibitem{SaitoIrrational} K. Saito, \emph{Mills' constant is irrational},
Mathematika \textbf{71} (2025), e70027.
Question numbering used here: arXiv:2404.19461v2.
\url{https://arxiv.org/html/2404.19461v2}.

\bibitem{SaitoVariants} K. Saito,
\emph{Transcendency of variants of Mills' constant},
The Ramanujan Journal \textbf{70} (2026), article 69.
Theorem numbering used here: arXiv:2508.16068v3, Theorems 1.7--1.8.
\url{https://arxiv.org/html/2508.16068v3}.

\bibitem{SaitoDimension} K. Saito,
\emph{Prime-representing functions and Hausdorff dimension},
Acta Mathematica Hungarica \textbf{165} (2021), 203--217.
Question numbering: arXiv:2102.04038v1.
\url{https://arxiv.org/abs/2102.04038}.

\bibitem{Kowalski} E. Kowalski,
\emph{The principle of the large sieve}, arXiv:math/0610021v3,
Section 5. \url{https://arxiv.org/abs/math/0610021}.

\bibitem{Milne} J. S. Milne, \emph{Algebraic Geometry},
version 6.10 (2024), B\'ezout's theorem.
\url{https://www.jmilne.org/math/CourseNotes/AG.pdf}.

\bibitem{Sutherland} A. V. Sutherland,
\emph{Global class field theory, the Chebotarev density theorem},
MIT 18.785, Lecture 28 (2021).
\url{https://math.mit.edu/classes/18.785/2021fa/LectureNotes28.pdf}.

\bibitem{Balady} S. Balady, \emph{Families of cyclic cubic fields},
arXiv:1511.04783v1 (2015).
\url{https://arxiv.org/html/1511.04783v1}.

\bibitem{Schinzel} A. Schinzel and W. Sierpi\'nski,
\emph{Sur certaines hypoth\`eses concernant les nombres premiers},
Acta Arithmetica \textbf{4} (1958), 185--208.
Hypothesis H is used as a conjectural assumption, not as a theorem.

\bibitem{Race} Showa, Heisei, and Taisho,
\emph{Constant proof race: Mills--Pisot submissions},
user-supplied research notes, 21 October 2026. In particular Showa
rounds 12, 20, 27, 28, 35, 36, 37 and Heisei round 05.
Source map above identifies the exact uses. Not external publications.

\bibitem{Keio007} Keio,
\emph{Dyadic Lucas Prime Cascades and Composite Fibonacci Shifts},
user-supplied manuscript [007], Theorem 1.1 and Corollary 1.2.

\bibitem{Keio008} Keio,
\emph{Additive Clocks, Non-Reversible Recurrences, and Transcendental
Prime-Representing Constants}, user-supplied manuscript [008],
Theorem 4.3.

\bibitem{Keio009} Keio,
\emph{Composite Terms in Sparse Cubic Pisot Floors
(Alternating Shifts, 3-Adic Limits, and Projective Return)},
user-supplied manuscript [009].

\bibitem{Keio031} Keio,
\emph{Prime-Representing Constants and Related Structures
(Encodings, Arithmetic, Geometry, and Modular Obstructions)},
user-supplied manuscript [031]. In particular Sections 14--15 and
the earlier cubic-floor and density arguments.

\bibitem{LM95} L. Lorch and M. E. Muldoon, \emph{Transcendentality of zeros
of higher derivatives of functions involving Bessel functions},
International Journal of Mathematics and Mathematical Sciences
\textbf{18} (1995), no. 3, 551--560.
\href{https://doi.org/10.1155/S0161171295000706}{doi:10.1155/S0161171295000706}.

\bibitem{SourceExt} Keio, \emph{Extensions of the winding line: three harmonic
coordinates and affine Bessel projections}, supplied research supplement,
22 October 2026; incorporated into Part I.

\bibitem{SourcePisot} Keio, \emph{Cubic Pisot Floors: Composite Windows,
Sparse Exceptions, and Finite-Sieve Barriers}, supplied revised research
manuscript, 21 October 2026; incorporated into Part II.

\end{thebibliography}
\end{document}
